% ----------------------------------------------------------------------
%  1.4. Структура кристаллов
% ----------------------------------------------------------------------
\begin{orientir}
Как описывают конкретные кристаллические структуры. Главная идея~—
\vazhno{плотнейшие упаковки шаров и их пустоты}: большинство структур
металлов и многих соединений~— это плотная упаковка одних атомов, в
пустотах которой сидят другие. Разберём структурные типы элементов
(1.4.1--1.4.2), соединений AB, AB$_2$ и сложных (1.4.3--1.4.5),
сверхструктуры (1.4.6), фуллерены и квазикристаллы (1.4.7) и структуру
поверхности (1.4.8).
\end{orientir}

Чтобы описать структуру, указывают пространственную группу, координаты
атомов в ячейке, координационные числа и многогранники.

\begin{opredelenie}[Координация]
\term{Координационное число} (к.\,ч.)~— число ближайших атомов (ионов)
данного сорта, находящихся на одинаковом расстоянии от центрального.
\term{Координационный многогранник}~— многогранник, вершины которого
заняты этими соседями; число его вершин равно к.\,ч.
\end{opredelenie}

Кристаллы, структуры которых одинаковы с точностью до подобия, относят к
одному \term{структурному типу}; его называют по типичному представителю
(«тип меди», «тип NaCl»). В символике Strukturbericht буква указывает
стехиометрию: A~— элементы, B~— соединения AB, C~— AB$_2$, D~—
A$_m$B$_n$, E~— более сложные, L~— сверхструктуры; за буквой следует номер
(A1~— тип меди, B1~— тип NaCl и~т.\,д.).

\subsection{Плотнейшие упаковки}\label{ssec:packings}

Будем строить структуру из одинаковых несжимаемых шаров. В плотном
плоском слое (рис.~\ref{fig:cp-layer}) каждый шар касается шести соседей и
окружён шестью лунками, а каждая лунка~— тремя шарами. Лунок вдвое больше,
чем шаров; они двух сортов~— B (треугольники вершиной вниз) и C (вершиной
вверх). Следующий слой можно положить либо в лунки B, либо в лунки C.

\begin{figure}[ht]
\centering
\begin{subfigure}[b]{0.5\textwidth}\centering
\begin{tikzpicture}[scale=0.62]
  \foreach \i in {-3,...,3}{\foreach \j in {-2,...,2}{
    \pgfmathsetmacro{\px}{\i+0.5*\j}\pgfmathsetmacro{\py}{0.866*\j}
    \fill[atomA, draw=atomLine, line width=0.3pt] (\px,\py) circle (0.5);}}
  \foreach \i in {-2,...,2}{\foreach \j in {-1,0,1}{
    \pgfmathsetmacro{\px}{\i+0.5*\j+0.5}\pgfmathsetmacro{\py}{0.866*\j+0.2887}
    \node[font=\scriptsize\bfseries, fmRed] at (\px,\py) {B};
    \pgfmathsetmacro{\qx}{\i+0.5*\j+1.0}\pgfmathsetmacro{\qy}{0.866*\j+0.5774}
    \node[font=\scriptsize\bfseries, fmBlue] at (\qx,\qy) {C};}}
  \node[font=\scriptsize\bfseries] at (0,0) {A};
\end{tikzpicture}
\caption{}
\end{subfigure}\hfill
\begin{subfigure}[b]{0.46\textwidth}\centering
\begin{tikzpicture}[x=1cm, y=0.55cm]
  \foreach \k/\l in {0/A,1/B,2/A,3/B,4/A}{
    \node[font=\small] at (0,\k) {\l};}
  \foreach \k/\l in {0/A,1/B,2/C,3/A,4/B}{
    \node[font=\small] at (2.6,\k) {\l};}
  \draw[fmGray] (-0.5,-0.6) -- (-0.5,4.6) (2.1,-0.6) -- (2.1,4.6);
  \node[font=\small, align=center] at (0,-1.6) {ГПУ\\$\ldots ABAB\ldots$};
  \node[font=\small, align=center] at (2.6,-1.6) {ГЦК\\$\ldots ABCABC\ldots$};
\end{tikzpicture}
\caption{}
\end{subfigure}
\caption{Плотнейшие упаковки: \textit{а}~— плотный слой шаров A и лунки
двух сортов B и C; \textit{б}~— две последовательности укладки слоёв:
двухслойная (гексагональная) и трёхслойная (кубическая)}
\label{fig:cp-layer}
\end{figure}

Во втором слое пустоты уже двух типов:
\begin{itemize}
\item \term{тетраэдрическая} ($T$)~— над лункой первого слоя лежит шар
второго; пустота окружена четырьмя шарами по вершинам тетраэдра;
\item \term{октаэдрическая} ($O$)~— лунка второго слоя над лункой первого;
пустота окружена шестью шарами по вершинам октаэдра.
\end{itemize}
Октаэдрических пустот столько же, сколько шаров, тетраэдрических~— вдвое
больше. В октаэдрическую пустоту помещается шар радиусом $0{,}414R$, в
тетраэдрическую~— $0{,}225R$ ($R$~— радиус шаров упаковки).

Третий слой можно положить в лунки $T$ (над шарами первого слоя)~—
получится \emph{двухслойная} упаковка $\ldots ABAB\ldots$,~— или в лунки $O$~—
\emph{трёхслойная} $\ldots ABCABC\ldots$ Возможны и другие чередования (кроме
повторения одной буквы), но плотнейших всего две.

\begin{glavnoe}[Две плотнейшие упаковки]
\begin{itemize}
\item \emph{Гексагональная} (ГПУ, $\ldots ABAB\ldots$): к.\,ч.~=~12 (6 в слое, по 3
сверху и снизу); идеальное отношение осей
$c/a=\sqrt{8/3}=1{,}633$; октаэдрические пустоты образуют сквозные каналы
вдоль оси $c$ (пути быстрой диффузии примесей).
\item \emph{Кубическая} ($\ldots ABCABC\ldots$): шары~— в узлах ГЦК-решётки,
плотные слои~— $\{111\}$, плотнейшие ряды~— $\nfam{110}$, к.\,ч.~=~12;
сквозных каналов нет.
\end{itemize}
Коэффициент компактности (доля объёма, занятая шарами) у обеих
$\pi/\sqrt{18}=74{,}05\,\%$; у ОЦК~— $68\,\%$, у простой кубической~— $52\,\%$.
\end{glavnoe}

Идея плотнейших упаковок очень плодотворна: в большинстве структур более
крупные атомы образуют плотную (или почти плотную) упаковку, а структуры
различаются тем, какие и сколько пустот заняты другими атомами.

\subsection{Структурные типы химических элементов}\label{ssec:elem-types}

\paragraph{Тип меди (A1, ГЦК)}
Al, Ni, Pb, $\gamma$-Fe, Au, Ag, Ir, Pt и многие другие. Пространственная
группа $Fm\br3m$; базис: $0,0,0$; $\frac12,\frac12,0$; $\frac12,0,\frac12$;
$0,\frac12,\frac12$ (кратко: $0,0,0$; $\frac12,\frac12,0\,[3]$, где
«[3]» означает три циклические перестановки). Атомов в ячейке $Z=4$, к.\,ч.~=~12,
координационный многогранник~— \emph{кубооктаэдр}. Это кубическая
плотнейшая упаковка слоёв $\{111\}$.

Пустоты ГЦК-структуры (рис.~\ref{fig:fcc-voids}):
\begin{itemize}
\item \emph{октаэдрические}~— в центре куба $\frac12,\frac12,\frac12$ и на серединах
рёбер $\frac12,0,0\,[3]$; на ячейку $1+12/4=4$; радиус
$r=(\sqrt2-1)R=0{,}414R$;
\item \emph{тетраэдрические}~— в центрах восьми «малых кубиков»
$\frac14,\frac14,\frac14$ и т.\,д.; на ячейку 8; $r=(\sqrt{3/2}-1)R=0{,}225R$.
\end{itemize}
Атомы : октапустоты : тетрапустоты $=1:1:2$, размеры
$1:0{,}414:0{,}225$. Многие соединения строятся как плотнейшая упаковка
крупных анионов, в пустотах которой сидят мелкие катионы.

\begin{figure}[ht]
\centering
\begin{subfigure}[b]{0.3\textwidth}\centering
\begin{tikzpicture}[cubeproj, scale=2.2]
  \CubeEdges[atomLine!60]
  \FCCback[2.6mm] \FCCmid[2.6mm] \FCCfront[2.6mm]
  \foreach \p in {(0.5,0.5,0.5),(0.5,0,0),(0.5,1,0),(0.5,0,1),(0.5,1,1),(0,0.5,0),(1,0.5,0),(0,0.5,1),(1,0.5,1),(0,0,0.5),(1,0,0.5),(0,1,0.5),(1,1,0.5)}
    \node[circle, fill=fmRed, inner sep=0pt, minimum size=4.2pt] at \p {};
\end{tikzpicture}
\caption{октаэдрические}
\end{subfigure}\hspace{6mm}
\begin{subfigure}[b]{0.3\textwidth}\centering
\begin{tikzpicture}[cubeproj, scale=2.2]
  \CubeEdges[atomLine!60]
  \FCCback[2.6mm] \FCCmid[2.6mm] \FCCfront[2.6mm]
  \foreach \x in {0.25,0.75}{\foreach \y in {0.25,0.75}{\foreach \z in {0.25,0.75}{
    \node[circle, fill=fmGreen!70!black, inner sep=0pt, minimum size=4.2pt] at (\x,\y,\z) {};}}}
\end{tikzpicture}
\caption{тетраэдрические}
\end{subfigure}\hspace{6mm}
\begin{subfigure}[b]{0.3\textwidth}\centering
\begin{tikzpicture}[cubeproj, scale=2.2]
  \CubeEdges[atomLine!60]
  \fill[fmBlue!20, draw=fmBlue, line width=0.5pt] (1,0,0) -- (0,1,0) -- (0,0,1) -- cycle;
  \fill[fmBlue!10, draw=fmBlue!60, line width=0.4pt, dashed] (1,1,0) -- (0,1,1) -- (1,0,1) -- cycle;
  \FCCback[2.6mm] \FCCmid[2.6mm] \FCCfront[2.6mm]
\end{tikzpicture}
\caption{плотные слои $(111)$}
\end{subfigure}
\caption{ГЦК-структура: \textit{а}~— октаэдрические пустоты (центр куба и
середины рёбер); \textit{б}~— тетраэдрические пустоты; \textit{в}~—
плотноупакованные плоскости $(111)$}
\label{fig:fcc-voids}
\end{figure}

\paragraph{Тип $\alpha$-вольфрама (A2, ОЦК)}
Cr, V, Mo, Nb, Ta, $\alpha$-Fe, щелочные металлы, высокотемпературные
$\beta$-Ti, $\beta$-Zr, $\beta$-Hf. Группа $Im\br3m$; базис $0,0,0$; $\frac12,\frac12,\frac12$;
$Z=2$; к.\,ч.~=~8 (многогранник~— куб). Вторая координационная сфера
(6~атомов на расстоянии $a$) лишь на 15\,\% дальше первой ($a\sqrt3/2$), поэтому
часто говорят о к.\,ч.~=~8+6=14. Пустоты ОЦК-структуры
(рис.~\ref{fig:bcc-voids}):
\begin{itemize}
\item \emph{тетраэдрические} (самые крупные!)~— в точках $\frac12,\frac14,0$ и
эквивалентных, в центрах слегка искажённых тетраэдров; их 12 на ячейку
(по 6 на атом); $r=(\sqrt{5/3}-1)R=0{,}291R$;
\item \emph{октаэдрические}~— в центрах граней и серединах рёбер; 6 на
ячейку ($6\cdot\frac12+12\cdot\frac14$); пустота~— сплюснутый октаэдр,
$r=(2/\sqrt3-1)R=0{,}155R$.
\end{itemize}
Атомы : октапустоты : тетрапустоты $=1:3:6$. Хотя тетрапустоты крупнее,
углерод в $\alpha$-железе сидит в \emph{октаэдрических}: сплюснутый октаэдр
раздвигает лишь два ближайших атома (вдоль оси сплющивания), и это
энергетически выгоднее; отсюда тетрагональность мартенсита
(п.~\ref{ssec:martensite-cryst}).

\begin{figure}[ht]
\centering
\begin{tikzpicture}[cubeproj, scale=2.6]
  \CubeEdges[atomLine!60]
  \foreach \p in {(0,0,0),(1,0,0),(0,1,0),(1,1,0),(0,0,1),(1,0,1),(0,1,1),(1,1,1),(0.5,0.5,0.5)}
    \node[atom=3.6mm] at \p {};
  \foreach \p in {(0.5,0.5,0),(0.5,0.5,1),(0.5,0,0.5),(0.5,1,0.5),(0,0.5,0.5),(1,0.5,0.5),
                  (0.5,0,0),(0.5,1,0),(0,0.5,0),(1,0.5,0),(0.5,0,1),(0.5,1,1),(0,0.5,1),(1,0.5,1),(0,0,0.5),(1,0,0.5),(0,1,0.5),(1,1,0.5)}
    \node[circle, fill=fmRed, inner sep=0pt, minimum size=4pt] at \p {};
  \foreach \p in {(0.5,0.25,0),(0.5,0.75,0),(0.25,0.5,0),(0.75,0.5,0),(0.5,0,0.25),(0.5,0,0.75),(0.25,0,0.5),(0.75,0,0.5),
                  (1,0.5,0.25),(1,0.5,0.75),(1,0.25,0.5),(1,0.75,0.5)}
    \node[regular polygon, regular polygon sides=3, fill=fmGreen!70!black, inner sep=0pt, minimum size=5.5pt] at \p {};
\end{tikzpicture}
\caption{ОЦК-структура: октаэдрические пустоты (красные кружки~— центры
граней и середины рёбер) и тетраэдрические пустоты (зелёные треугольники;
показаны только на трёх видимых гранях)}
\label{fig:bcc-voids}
\end{figure}

\paragraph{Тип магния (A3, ГПУ)}
Be, Mg, Zn, Cd, Re, Os, $\alpha$-Ti, $\alpha$-Zr, $\alpha$-Hf, Co и др. Группа
$P6_3/mmc$, решётка Браве гексагональная примитивная; в ячейке 2 атома:
$0,0,0$ и $\frac23,\frac13,\frac12$ (рис.~\ref{fig:hcp}); к.\,ч.~=~12
(многогранник~— гексагональный, или «двойниковый», кубооктаэдр).
Октаэдрические пустоты: $\frac13,\frac23,\frac14$ и $\frac13,\frac23,\frac34$ (2 на ячейку,
$r=0{,}414R$); тетраэдрические: $0,0,\frac38$; $0,0,\frac58$; $\frac23,\frac13,\frac18$;
$\frac23,\frac13,\frac78$ (4 на ячейку, $r=0{,}225R$)~— соотношения те же, что в ГЦК.
У большинства ГПУ-металлов $c/a=1{,}57$--$1{,}63$, близко к идеальному;
исключения~— цинк (1,86) и кадмий (1,89).

\begin{figure}[ht]
\centering
\begin{tikzpicture}[cubeproj, scale=1.25]
  \def\c{1.633}
  \foreach \k in {0,...,5}{
    \pgfmathsetmacro{\xa}{cos(60*\k)}\pgfmathsetmacro{\ya}{sin(60*\k)}
    \pgfmathsetmacro{\xb}{cos(60*\k+60)}\pgfmathsetmacro{\yb}{sin(60*\k+60)}
    \pgfmathtruncatemacro{\back}{(\k==1 || \k==2) ? 1 : 0}
    \ifnum\back=1
      \draw[atomLine!50, dashed] (\xa,\ya,0) -- (\xb,\yb,0);
    \else
      \draw[atomLine] (\xa,\ya,0) -- (\xb,\yb,0);
    \fi
    \draw[atomLine] (\xa,\ya,\c) -- (\xb,\yb,\c);
    \ifnum\k=2 \draw[atomLine!50, dashed] (\xa,\ya,0) -- (\xa,\ya,\c);
    \else \draw[atomLine] (\xa,\ya,0) -- (\xa,\ya,\c);\fi}
  \foreach \k in {0,...,5}{\node[atom=3.4mm] at ({cos(60*\k)},{sin(60*\k)},0) {};}
  \node[atom=3.4mm] at (0,0,0) {};
  \foreach \k in {0,1,2}{\node[atomB=3.4mm] at ({0.577*cos(120*\k+30)},{0.577*sin(120*\k+30)},{\c/2}) {};}
  \foreach \k in {0,...,5}{\node[atom=3.4mm] at ({cos(60*\k)},{sin(60*\k)},\c) {};}
  \node[atom=3.4mm] at (0,0,\c) {};
  \node[font=\small, anchor=west] at (1.15,0,\c) {A};
  \node[font=\small, anchor=west] at (1.15,0,{\c/2}) {B};
  \node[font=\small, anchor=west] at (1.15,0,0) {A};
\end{tikzpicture}
\caption{ГПУ-структура (тип магния): гексагональная призма из трёх
элементарных ячеек; слои A (голубые атомы) и B (оранжевые) чередуются
вдоль оси $c$}
\label{fig:hcp}
\end{figure}

\paragraph{Тип алмаза (A4)}
Алмаз, Si, Ge, серое олово ($\alpha$-Sn). Группа $Fd\br3m$; решётка Браве ГЦК,
базис из двух атомов $0,0,0$ и $\frac14,\frac14,\frac14$, т.\,е. атомы занимают узлы
ГЦК-решётки и половину тетраэдрических пустот (через одну)
(рис.~\ref{fig:diamond-graphite}\,\textit{а}). Иначе: две ГЦК-решётки,
сдвинутые на $\frac14$ пространственной диагонали. $Z=8$, к.\,ч.~=~4,
многогранник~— тетраэдр\ispr{В оригинале для алмаза указано
к.\,ч.~=~12; при тетраэдрическом координационном многограннике
к.\,ч.~=~4.}. Центр симметрии лежит посередине связи между соседними
атомами. Гексагональная модификация~— \emph{лонсдейлит}: двухслойная
упаковка, в которой атомами заполнена половина тетрапустот.

\paragraph{Тип графита (A9)}
Группа $P6_3/mmc$, гексагональная решётка, 4 атома в ячейке, к.\,ч.~=~3
(многогранник~— правильный треугольник). Структура слоистая
(рис.~\ref{fig:diamond-graphite}\,\textit{б}): внутри слоя прочные ковалентные
связи (расстояние 0,142~нм), между слоями~— слабые ван-дер-ваальсовы
(0,335~нм). Отсюда огромная анизотропия: электропроводность в плоскости
слоёв в $\sim10^4$--$10^5$ раз больше, чем вдоль оси $c$\ispr{В оригинале
сказано наоборот: «электропроводность вдоль оси $c$ в $10^5$ раз больше, чем в
поперечных направлениях». Проводимость графита велика именно в плоскости
слоёв.}, а лёгкое скольжение по $(0001)$ делает графит смазкой. Есть и
ромбоэдрический политип с трёхслойной укладкой ($R\br3m$); ещё одна
форма углерода~— линейный полимер \emph{карбин}.

\begin{figure}[ht]
\centering
\begin{subfigure}[b]{0.4\textwidth}\centering
\begin{tikzpicture}[cubeproj, scale=1.9, anion/.style={atom=3.2mm}, cation/.style={atom=3.2mm}]
  \CubeEdges
  % сгенерировано tools/gen_ionic_cells.py — не править вручную
\draw[bondthick] (0.25,0.25,0.25) -- (0,0,0);
\draw[bondthick] (0.25,0.25,0.25) -- (0,0.5,0.5);
\draw[bondthick] (0.25,0.25,0.25) -- (0.5,0,0.5);
\draw[bondthick] (0.25,0.25,0.25) -- (0.5,0.5,0);
\draw[bondthick] (0.75,0.75,0.25) -- (0.5,0.5,0);
\draw[bondthick] (0.75,0.75,0.25) -- (0.5,1,0.5);
\draw[bondthick] (0.75,0.75,0.25) -- (1,0.5,0.5);
\draw[bondthick] (0.75,0.75,0.25) -- (1,1,0);
\draw[bondthick] (0.75,0.25,0.75) -- (0.5,0,0.5);
\draw[bondthick] (0.75,0.25,0.75) -- (0.5,0.5,1);
\draw[bondthick] (0.75,0.25,0.75) -- (1,0,1);
\draw[bondthick] (0.75,0.25,0.75) -- (1,0.5,0.5);
\draw[bondthick] (0.25,0.75,0.75) -- (0,0.5,0.5);
\draw[bondthick] (0.25,0.75,0.75) -- (0,1,1);
\draw[bondthick] (0.25,0.75,0.75) -- (0.5,0.5,1);
\draw[bondthick] (0.25,0.75,0.75) -- (0.5,1,0.5);
\node[anion] at (0,1,0) {};
\node[anion] at (1,1,0) {};
\node[anion] at (0.5,1,0.5) {};
\node[anion] at (0,1,1) {};
\node[anion] at (1,1,1) {};
\node[cation] at (0.75,0.75,0.25) {};
\node[cation] at (0.25,0.75,0.75) {};
\node[anion] at (0.5,0.5,0) {};
\node[anion] at (0,0.5,0.5) {};
\node[anion] at (1,0.5,0.5) {};
\node[anion] at (0.5,0.5,1) {};
\node[cation] at (0.25,0.25,0.25) {};
\node[cation] at (0.75,0.25,0.75) {};
\node[anion] at (0,0,0) {};
\node[anion] at (1,0,0) {};
\node[anion] at (0.5,0,0.5) {};
\node[anion] at (0,0,1) {};
\node[anion] at (1,0,1) {};

\end{tikzpicture}
\caption{}
\end{subfigure}\hfill
\begin{subfigure}[b]{0.55\textwidth}\centering
\begin{tikzpicture}[scale=0.75]
  \clip (-3.4,-2.6) rectangle (3.4,2.6);
  \foreach \i in {-2,...,2}{\foreach \j in {-2,...,2}{
    \pgfmathsetmacro{\cx}{1.732*\i+0.866*\j}\pgfmathsetmacro{\cy}{1.5*\j}
    \foreach \k in {0,...,5}{
      \draw[fmBlue!40, line width=0.6pt, dashed] ($(\cx+0.866,\cy+0.5)+({cos(60*\k+30)},{sin(60*\k+30)})$) -- ($(\cx+0.866,\cy+0.5)+({cos(60*\k+90)},{sin(60*\k+90)})$);
      \draw[atomLine, line width=0.9pt] ($(\cx,\cy)+({cos(60*\k+30)},{sin(60*\k+30)})$) -- ($(\cx,\cy)+({cos(60*\k+90)},{sin(60*\k+90)})$);}}}
\end{tikzpicture}
\caption{}
\end{subfigure}
\caption{\textit{а}~— структура алмаза (атомы в узлах ГЦК-решётки и в
половине тетраэдрических пустот, линии~— ковалентные связи);
\textit{б}~— графит, вид вдоль оси $c$: сетки из шестиугольников
соседних слоёв A (сплошные) и B (штриховые) сдвинуты друг относительно
друга}
\label{fig:diamond-graphite}
\end{figure}

\paragraph{Тип $\beta$-вольфрама (A15)}
Группа $Pm\br3n$, кубическая примитивная решётка, 8 атомов в ячейке двух
сортов (рис.~\ref{fig:a15}): W$_\text{I}$ в $0,0,0$ и $\frac12,\frac12,\frac12$ (окружены
12 атомами W$_\text{II}$), W$_\text{II}$ парами на гранях ($0,\frac14,\frac12$; $0,\frac34,\frac12$
и~т.\,д.; к.\,ч.~=~14: $2+4+8$), образующими цепочки вдоль трёх осей куба.
Название историческое: за «$\beta$-вольфрам» приняли оксид W$_3$O. Сейчас
тип интересен тем, что его имеют многие низкотемпературные
сверхпроводники: Nb$_3$Sn ($T_\text{с}=18$~К), Nb$_3$Ge (23~К).

\begin{figure}[ht]
\centering
\begin{tikzpicture}[cubeproj, scale=2.4]
  \CubeEdges[atomLine!60]
  \foreach \p in {(0,0,0),(1,0,0),(0,1,0),(1,1,0),(0,0,1),(1,0,1),(0,1,1),(1,1,1),(0.5,0.5,0.5)}
    \node[atom=3.4mm] at \p {};
  \foreach \a/\b in {{(0,0.25,0.5)}/{(0,0.75,0.5)}, {(1,0.25,0.5)}/{(1,0.75,0.5)},
                     {(0.5,0,0.25)}/{(0.5,0,0.75)}, {(0.5,1,0.25)}/{(0.5,1,0.75)},
                     {(0.25,0.5,0)}/{(0.75,0.5,0)}, {(0.25,0.5,1)}/{(0.75,0.5,1)}}{
    \draw[fmAmber!80!black, line width=1.4pt] \a -- \b;
    \node[atomB=3.0mm] at \a {}; \node[atomB=3.0mm] at \b {};}
\end{tikzpicture}
\caption{Структура A15 ($\beta$-W, Nb$_3$Sn): атомы W$_\text{I}$ (Sn) в вершинах и
центре куба, атомы W$_\text{II}$ (Nb)~— парами на гранях; пары образуют три
семейства непересекающихся цепочек вдоль осей куба}
\label{fig:a15}
\end{figure}

\subsection{Структурные типы соединений AB}\label{ssec:ab-types}

Ионные структуры устойчивы, если каждый ион касается только ионов
противоположного знака; поэтому координация определяется отношением
радиусов катиона и аниона $r_\text{к}/r_\text{а}$ (подробно~— гл.~3, п.~3.1.5):
при $0{,}41$--$0{,}73$ выгоден октаэдр (к.\,ч.~=~6), при $0{,}73$--$1{,}0$~— куб
(к.\,ч.~=~8), при $0{,}225$--$0{,}41$~— тетраэдр (к.\,ч.~=~4).

\begin{figure}[ht]
\centering
\begin{subfigure}[b]{0.3\textwidth}\centering
\begin{tikzpicture}[cubeproj, scale=1.8, anion/.style={atom=4.4mm}, cation/.style={atomB=2.8mm}]
  \CubeEdges[draw opacity=0]
  % сгенерировано tools/gen_ionic_cells.py — не править вручную
\draw[bondline] (0,0,0) -- (0,0,0.5);
\draw[bondline] (0,0,0) -- (0,0.5,0);
\draw[bondline] (0,0,0) -- (0.5,0,0);
\draw[bondline] (0,0,0.5) -- (0,0,1);
\draw[bondline] (0,0,0.5) -- (0,0.5,0.5);
\draw[bondline] (0,0,0.5) -- (0.5,0,0.5);
\draw[bondline] (0,0,1) -- (0,0.5,1);
\draw[bondline] (0,0,1) -- (0.5,0,1);
\draw[bondline] (0,0.5,0) -- (0,0.5,0.5);
\draw[bondline] (0,0.5,0) -- (0,1,0);
\draw[bondline] (0,0.5,0) -- (0.5,0.5,0);
\draw[bondline] (0,0.5,0.5) -- (0,0.5,1);
\draw[bondline] (0,0.5,0.5) -- (0,1,0.5);
\draw[bondline] (0,0.5,0.5) -- (0.5,0.5,0.5);
\draw[bondline] (0,0.5,1) -- (0,1,1);
\draw[bondline] (0,0.5,1) -- (0.5,0.5,1);
\draw[bondline] (0,1,0) -- (0,1,0.5);
\draw[bondline] (0,1,0) -- (0.5,1,0);
\draw[bondline] (0,1,0.5) -- (0,1,1);
\draw[bondline] (0,1,0.5) -- (0.5,1,0.5);
\draw[bondline] (0,1,1) -- (0.5,1,1);
\draw[bondline] (0.5,0,0) -- (0.5,0,0.5);
\draw[bondline] (0.5,0,0) -- (0.5,0.5,0);
\draw[bondline] (0.5,0,0) -- (1,0,0);
\draw[bondline] (0.5,0,0.5) -- (0.5,0,1);
\draw[bondline] (0.5,0,0.5) -- (0.5,0.5,0.5);
\draw[bondline] (0.5,0,0.5) -- (1,0,0.5);
\draw[bondline] (0.5,0,1) -- (0.5,0.5,1);
\draw[bondline] (0.5,0,1) -- (1,0,1);
\draw[bondline] (0.5,0.5,0) -- (0.5,0.5,0.5);
\draw[bondline] (0.5,0.5,0) -- (0.5,1,0);
\draw[bondline] (0.5,0.5,0) -- (1,0.5,0);
\draw[bondline] (0.5,0.5,0.5) -- (0.5,0.5,1);
\draw[bondline] (0.5,0.5,0.5) -- (0.5,1,0.5);
\draw[bondline] (0.5,0.5,0.5) -- (1,0.5,0.5);
\draw[bondline] (0.5,0.5,1) -- (0.5,1,1);
\draw[bondline] (0.5,0.5,1) -- (1,0.5,1);
\draw[bondline] (0.5,1,0) -- (0.5,1,0.5);
\draw[bondline] (0.5,1,0) -- (1,1,0);
\draw[bondline] (0.5,1,0.5) -- (0.5,1,1);
\draw[bondline] (0.5,1,0.5) -- (1,1,0.5);
\draw[bondline] (0.5,1,1) -- (1,1,1);
\draw[bondline] (1,0,0) -- (1,0,0.5);
\draw[bondline] (1,0,0) -- (1,0.5,0);
\draw[bondline] (1,0,0.5) -- (1,0,1);
\draw[bondline] (1,0,0.5) -- (1,0.5,0.5);
\draw[bondline] (1,0,1) -- (1,0.5,1);
\draw[bondline] (1,0.5,0) -- (1,0.5,0.5);
\draw[bondline] (1,0.5,0) -- (1,1,0);
\draw[bondline] (1,0.5,0.5) -- (1,0.5,1);
\draw[bondline] (1,0.5,0.5) -- (1,1,0.5);
\draw[bondline] (1,0.5,1) -- (1,1,1);
\draw[bondline] (1,1,0) -- (1,1,0.5);
\draw[bondline] (1,1,0.5) -- (1,1,1);
\node[anion] at (0,1,0) {};
\node[cation] at (0.5,1,0) {};
\node[anion] at (1,1,0) {};
\node[cation] at (0,1,0.5) {};
\node[anion] at (0.5,1,0.5) {};
\node[cation] at (1,1,0.5) {};
\node[anion] at (0,1,1) {};
\node[cation] at (0.5,1,1) {};
\node[anion] at (1,1,1) {};
\node[cation] at (0,0.5,0) {};
\node[anion] at (0.5,0.5,0) {};
\node[cation] at (1,0.5,0) {};
\node[anion] at (0,0.5,0.5) {};
\node[cation] at (0.5,0.5,0.5) {};
\node[anion] at (1,0.5,0.5) {};
\node[cation] at (0,0.5,1) {};
\node[anion] at (0.5,0.5,1) {};
\node[cation] at (1,0.5,1) {};
\node[anion] at (0,0,0) {};
\node[cation] at (0.5,0,0) {};
\node[anion] at (1,0,0) {};
\node[cation] at (0,0,0.5) {};
\node[anion] at (0.5,0,0.5) {};
\node[cation] at (1,0,0.5) {};
\node[anion] at (0,0,1) {};
\node[cation] at (0.5,0,1) {};
\node[anion] at (1,0,1) {};

\end{tikzpicture}
\caption{NaCl (B1)}
\end{subfigure}\hfill
\begin{subfigure}[b]{0.3\textwidth}\centering
\begin{tikzpicture}[cubeproj, scale=1.8, anion/.style={atom=4.4mm}, cation/.style={atomB=4.0mm}]
  \CubeEdges
  % сгенерировано tools/gen_ionic_cells.py — не править вручную
\draw[bondline] (0.5,0.5,0.5) -- (0,0,0);
\draw[bondline] (0.5,0.5,0.5) -- (0,0,1);
\draw[bondline] (0.5,0.5,0.5) -- (0,1,0);
\draw[bondline] (0.5,0.5,0.5) -- (0,1,1);
\draw[bondline] (0.5,0.5,0.5) -- (1,0,0);
\draw[bondline] (0.5,0.5,0.5) -- (1,0,1);
\draw[bondline] (0.5,0.5,0.5) -- (1,1,0);
\draw[bondline] (0.5,0.5,0.5) -- (1,1,1);
\node[anion] at (0,1,0) {};
\node[anion] at (1,1,0) {};
\node[anion] at (0,1,1) {};
\node[anion] at (1,1,1) {};
\node[cation] at (0.5,0.5,0.5) {};
\node[anion] at (0,0,0) {};
\node[anion] at (1,0,0) {};
\node[anion] at (0,0,1) {};
\node[anion] at (1,0,1) {};

\end{tikzpicture}
\caption{CsCl (B2)}
\end{subfigure}\hfill
\begin{subfigure}[b]{0.3\textwidth}\centering
\begin{tikzpicture}[cubeproj, scale=1.8, anion/.style={atom=4.2mm}, cation/.style={atomB=2.8mm}]
  \CubeEdges
  % сгенерировано tools/gen_ionic_cells.py — не править вручную
\draw[bondthick] (0.25,0.25,0.25) -- (0,0,0);
\draw[bondthick] (0.25,0.25,0.25) -- (0,0.5,0.5);
\draw[bondthick] (0.25,0.25,0.25) -- (0.5,0,0.5);
\draw[bondthick] (0.25,0.25,0.25) -- (0.5,0.5,0);
\draw[bondthick] (0.75,0.75,0.25) -- (0.5,0.5,0);
\draw[bondthick] (0.75,0.75,0.25) -- (0.5,1,0.5);
\draw[bondthick] (0.75,0.75,0.25) -- (1,0.5,0.5);
\draw[bondthick] (0.75,0.75,0.25) -- (1,1,0);
\draw[bondthick] (0.75,0.25,0.75) -- (0.5,0,0.5);
\draw[bondthick] (0.75,0.25,0.75) -- (0.5,0.5,1);
\draw[bondthick] (0.75,0.25,0.75) -- (1,0,1);
\draw[bondthick] (0.75,0.25,0.75) -- (1,0.5,0.5);
\draw[bondthick] (0.25,0.75,0.75) -- (0,0.5,0.5);
\draw[bondthick] (0.25,0.75,0.75) -- (0,1,1);
\draw[bondthick] (0.25,0.75,0.75) -- (0.5,0.5,1);
\draw[bondthick] (0.25,0.75,0.75) -- (0.5,1,0.5);
\node[anion] at (0,1,0) {};
\node[anion] at (1,1,0) {};
\node[anion] at (0.5,1,0.5) {};
\node[anion] at (0,1,1) {};
\node[anion] at (1,1,1) {};
\node[cation] at (0.75,0.75,0.25) {};
\node[cation] at (0.25,0.75,0.75) {};
\node[anion] at (0.5,0.5,0) {};
\node[anion] at (0,0.5,0.5) {};
\node[anion] at (1,0.5,0.5) {};
\node[anion] at (0.5,0.5,1) {};
\node[cation] at (0.25,0.25,0.25) {};
\node[cation] at (0.75,0.25,0.75) {};
\node[anion] at (0,0,0) {};
\node[anion] at (1,0,0) {};
\node[anion] at (0.5,0,0.5) {};
\node[anion] at (0,0,1) {};
\node[anion] at (1,0,1) {};

\end{tikzpicture}
\caption{ZnS, сфалерит (B3)}
\end{subfigure}
\caption{Структурные типы соединений AB (голубые~— анионы, оранжевые~—
катионы)}
\label{fig:ab-types}
\end{figure}

\begin{itemize}
\item \term{NaCl (B1)}, $Fm\br3m$: анионы Cl$^-$ в узлах ГЦК-решётки ($0,0,0$;
$\frac12,\frac12,0\,[3]$), катионы Na$^+$ во \emph{всех октаэдрических} пустотах
($\frac12,\frac12,\frac12$; $\frac12,0,0\,[3]$); тетрапустоты свободны. К.\,ч.~=~6 для
обоих ионов (октаэдр), во второй сфере~— 12 одноимённых. Так же построены
оксиды TiO, MnO, FeO, NiO, карбиды и нитриды переходных металлов IV--V
групп (TiC, VN\ldots). Для NaCl $r_\text{Na}/r_\text{Cl}=0{,}98/1{,}81=0{,}54$~—
октаэдрический интервал.
\item \term{CsCl (B2)}, $Pm\br3m$: Cl$^-$ в вершинах куба, Cs$^+$ в центре;
к.\,ч.~=~8 (куб). Решётка Браве~— \emph{простая} кубическая (центр и вершины
заняты разными ионами). $r_\text{Cs}/r_\text{Cl}=1{,}65/1{,}81=0{,}91$~— кубический
интервал, к.\,ч.~=~8\ispr{В оригинале для CsCl сказано, что отношение
радиусов 0,91 «соответствует ОЦК структуре с к.\,ч.~=~6»; правильно~—
к.\,ч.~=~8, а решётка Браве простая кубическая.}. Этот тип имеют многие
интерметаллиды с близкими радиусами атомов (упорядоченная $\beta'$-латунь
CuZn, NiAl, FeAl).
\item \term{Сфалерит ZnS (B3)}, $F\br43m$: S в узлах ГЦК (кубическая плотнейшая
упаковка), Zn в \emph{половине тетраэдрических} пустот ($\frac14,\frac14,\frac14$;
$\frac34,\frac34,\frac14\,[3]$); к.\,ч.~=~4 у обоих (тетраэдры). Центра симметрии нет~—
структура полярна (вдоль $\nfam{111}$), что важно для пьезоэффекта. Тот же
тип~— у GaAs, InSb, CdTe, кубического BN.
\item \term{Вюрцит ZnS (B4)}, $P6_3mc$: то же, но на основе
\emph{гексагональной} плотнейшей упаковки атомов S; Zn в половине
тетрапустот, к.\,ч.~=~4. Так построены ZnO, GaN, AlN, BeO.
\item \term{NiAs (B8)}, $P6_3/mmc$: As~— гексагональная плотнейшая упаковка,
Ni во \emph{всех октаэдрических} пустотах ($0,0,0$; $0,0,\frac12$), к.\,ч.~=~6, $Z=2$.
Характерен для соединений переходных металлов с S, Se, Te, As, Sb (FeS, CoS,
NiSb).
\end{itemize}

\subsection{Структурные типы соединений AB$_2$}\label{ssec:ab2-types}

\term{Флюорит CaF$_2$ (C1)}, $Fm\br3m$ (рис.~\ref{fig:fluorite-perov}\,\textit{а}):
ионы Ca$^{2+}$ в узлах ГЦК-решётки, ионы F$^-$ во \emph{всех тетраэдрических}
пустотах ($\frac14,\frac14,\frac14$ и т.\,д.). Каждый F$^-$ окружён тетраэдром из четырёх
Ca$^{2+}$, каждый Ca$^{2+}$~— кубом из восьми F$^-$: к.\,ч.(F)~=~4, к.\,ч.(Ca)~=~8.
Необычно, что плотнейшую упаковку образуют здесь \emph{катионы}. Так
построены UO$_2$, ThO$_2$, ZrO$_2$ (кубический), PuO$_2$~— основа ядерного
топлива. В \emph{антифлюорите} (Li$_2$O, Na$_2$O, Mg$_2$Si) катионы и анионы
меняются местами: плотную упаковку образуют анионы, катионы заполняют все
тетрапустоты.

\subsection{Структурные типы сложных соединений}\label{ssec:complex-types}

\term{Перовскит CaTiO$_3$ (E2$_1$)}, идеализированно $Pm\br3m$
(рис.~\ref{fig:fluorite-perov}\,\textit{б}): Ca в вершинах куба, Ti в центре,
O в центрах граней. К.\,ч.: Ti по O~— 6 (октаэдр TiO$_6$), Ca по O~— 12, O по
Ti~— 2, O по Ca~— 4. Реальный CaTiO$_3$ слегка искажён (октаэдры повёрнуты)
и имеет ромбическую симметрию. Перовскит~— одна из важнейших структур
материаловедения: сегнетоэлектрики (BaTiO$_3$), а также слоистые структуры,
производные от перовскита,~— высокотемпературные сверхпроводники
YBa$_2$Cu$_3$O$_{7-\delta}$ ($T_\text{с}=93$~К), Bi$_2$Sr$_2$CaCu$_2$O$_{8+x}$ (до 94~К),
Tl$_2$Ba$_2$Ca$_2$Cu$_3$O$_{10+x}$ (125~К), HgBa$_2$Ca$_2$Cu$_3$O$_{8+\delta}$ (133~К).

\begin{figure}[ht]
\centering
\begin{subfigure}[b]{0.4\textwidth}\centering
\begin{tikzpicture}[cubeproj, scale=1.9, anion/.style={atomB=3.4mm}, cation/.style={atom=4.4mm}]
  \CubeEdges
  % сгенерировано tools/gen_ionic_cells.py — не править вручную
\draw[bondline] (0.25,0.25,0.25) -- (0,0,0);
\draw[bondline] (0.25,0.25,0.25) -- (0,0.5,0.5);
\draw[bondline] (0.25,0.25,0.25) -- (0.5,0,0.5);
\draw[bondline] (0.25,0.25,0.25) -- (0.5,0.5,0);
\draw[bondline] (0.25,0.25,0.75) -- (0,0,1);
\draw[bondline] (0.25,0.25,0.75) -- (0,0.5,0.5);
\draw[bondline] (0.25,0.25,0.75) -- (0.5,0,0.5);
\draw[bondline] (0.25,0.25,0.75) -- (0.5,0.5,1);
\draw[bondline] (0.25,0.75,0.25) -- (0,0.5,0.5);
\draw[bondline] (0.25,0.75,0.25) -- (0,1,0);
\draw[bondline] (0.25,0.75,0.25) -- (0.5,0.5,0);
\draw[bondline] (0.25,0.75,0.25) -- (0.5,1,0.5);
\draw[bondline] (0.25,0.75,0.75) -- (0,0.5,0.5);
\draw[bondline] (0.25,0.75,0.75) -- (0,1,1);
\draw[bondline] (0.25,0.75,0.75) -- (0.5,0.5,1);
\draw[bondline] (0.25,0.75,0.75) -- (0.5,1,0.5);
\draw[bondline] (0.75,0.25,0.25) -- (0.5,0,0.5);
\draw[bondline] (0.75,0.25,0.25) -- (0.5,0.5,0);
\draw[bondline] (0.75,0.25,0.25) -- (1,0,0);
\draw[bondline] (0.75,0.25,0.25) -- (1,0.5,0.5);
\draw[bondline] (0.75,0.25,0.75) -- (0.5,0,0.5);
\draw[bondline] (0.75,0.25,0.75) -- (0.5,0.5,1);
\draw[bondline] (0.75,0.25,0.75) -- (1,0,1);
\draw[bondline] (0.75,0.25,0.75) -- (1,0.5,0.5);
\draw[bondline] (0.75,0.75,0.25) -- (0.5,0.5,0);
\draw[bondline] (0.75,0.75,0.25) -- (0.5,1,0.5);
\draw[bondline] (0.75,0.75,0.25) -- (1,0.5,0.5);
\draw[bondline] (0.75,0.75,0.25) -- (1,1,0);
\draw[bondline] (0.75,0.75,0.75) -- (0.5,0.5,1);
\draw[bondline] (0.75,0.75,0.75) -- (0.5,1,0.5);
\draw[bondline] (0.75,0.75,0.75) -- (1,0.5,0.5);
\draw[bondline] (0.75,0.75,0.75) -- (1,1,1);
\node[cation] at (0,1,0) {};
\node[cation] at (1,1,0) {};
\node[cation] at (0.5,1,0.5) {};
\node[cation] at (0,1,1) {};
\node[cation] at (1,1,1) {};
\node[anion] at (0.25,0.75,0.25) {};
\node[anion] at (0.75,0.75,0.25) {};
\node[anion] at (0.25,0.75,0.75) {};
\node[anion] at (0.75,0.75,0.75) {};
\node[cation] at (0.5,0.5,0) {};
\node[cation] at (0,0.5,0.5) {};
\node[cation] at (1,0.5,0.5) {};
\node[cation] at (0.5,0.5,1) {};
\node[anion] at (0.25,0.25,0.25) {};
\node[anion] at (0.75,0.25,0.25) {};
\node[anion] at (0.25,0.25,0.75) {};
\node[anion] at (0.75,0.25,0.75) {};
\node[cation] at (0,0,0) {};
\node[cation] at (1,0,0) {};
\node[cation] at (0.5,0,0.5) {};
\node[cation] at (0,0,1) {};
\node[cation] at (1,0,1) {};

\end{tikzpicture}
\caption{CaF$_2$}
\end{subfigure}\hspace{6mm}
\begin{subfigure}[b]{0.4\textwidth}\centering
\begin{tikzpicture}[cubeproj, scale=2.1]
  \CubeEdges[atomLine!60]
  \draw[fmBlue!50, line width=0.5pt] (0.5,0.5,0) -- (0.5,0.5,1) (0,0.5,0.5) -- (1,0.5,0.5) (0.5,0,0.5) -- (0.5,1,0.5);
  \foreach \p in {(0.5,1,0.5),(0.5,0.5,0),(0,0.5,0.5)} \node[atomI=2.8mm] at \p {};
  \node[atomB=3.2mm] at (0.5,0.5,0.5) {};
  \foreach \p in {(1,0.5,0.5),(0.5,0.5,1),(0.5,0,0.5)} \node[atomI=2.8mm] at \p {};
  \foreach \p in {(0,0,0),(1,0,0),(0,1,0),(1,1,0),(0,0,1),(1,0,1),(0,1,1),(1,1,1)} \node[atom=4.2mm] at \p {};
\end{tikzpicture}
\caption{CaTiO$_3$}
\end{subfigure}
\caption{\textit{а}~— флюорит CaF$_2$ (голубые~— Ca$^{2+}$, оранжевые~— F$^-$);
\textit{б}~— идеальный перовскит CaTiO$_3$: Ca (голубые) в вершинах, Ti
(оранжевый) в центре октаэдра из атомов O (красные) в центрах граней}
\label{fig:fluorite-perov}
\end{figure}

\subsection{Сверхструктуры в твёрдых растворах замещения}\label{ssec:superstructures}

Напомним (гл.~3, п.~3.3.2): твёрдые растворы бывают замещения,
внедрения (атомы H, C, N, O, B в пустотах металлической решётки) и
вычитания (на основе соединений, часть узлов пуста). Растворение атомов
меньшего размера уменьшает период решётки, большего~— увеличивает; в
непрерывных растворах замещения период часто меняется линейно с
составом~— \term{закон Вегарда}:
\begin{equation}\label{eq:vegard}
a=a_1c_1+a_2c_2=a_1+(a_2-a_1)c_2 ,
\end{equation}
где $a_i$~— периоды компонентов, $c_i$~— атомные доли.

При упорядочении на рентгенограммах появляются дополнительные,
\emph{сверхструктурные}, линии~— признак сверхструктуры (сверхрешётки).
Важнейшие сверхструктуры (рис.~\ref{fig:superstructures}):
\begin{itemize}
\item \term{Cu$_3$Au (L1$_2$)}: ниже $\approx390\,^\circ$C из неупорядоченной ГЦК
($Fm\br3m$) образуется структура с Au в вершинах куба и Cu в центрах граней;
группа $Pm\br3m$ (решётка стала примитивной).
\item \term{CuAu\,I (L1$_0$)}: ниже $\approx385\,^\circ$C слои $(001)$ заняты
поочерёдно атомами Cu и Au; решётка тетрагональна ($c/a\approx0{,}93$),
$P4/mmm$.
\item \term{CuPt (L1$_1$)}: упорядочение по плоскостям $(111)$ с
ромбоэдрическим искажением.
\item \term{CuAu\,II}~— сверхструктура с \emph{длинным периодом}: десять
ячеек CuAu\,I вдоль оси $b$, причём посередине периода слои меди и золота
меняются местами~— периодически расположенные антифазные границы
(одномерная длиннопериодическая сверхструктура).
\item На основе ОЦК: \term{$\beta'$-латунь CuZn (B2)} ниже $\approx465\,^\circ$C
(переход $Im\br3m\to Pm\br3m$); на ячейке с удвоенным периодом (8 ОЦК-ячеек,
центры которых делятся на позиции $X$ и $Y$, чередующиеся в шахматном
порядке): \term{FeAl} (Fe в вершинах, Al в $X$ и $Y$~— та же B2),
\term{Fe$_3$Al (D0$_3$)} (Fe в вершинах и $X$, Al в $Y$), \term{Cu$_2$MnAl (L2$_1$,
сплав Гейслера)} (Cu в вершинах, Mn в $X$, Al в $Y$; ферромагнетик из
неферромагнитных компонентов, см. гл.~3, п.~3.6.8).
\end{itemize}

\begin{figure}[ht]
\centering
\begin{subfigure}[b]{0.24\textwidth}\centering
\begin{tikzpicture}[cubeproj, scale=1.55]
  \CubeEdges
  \foreach \p in {(0,1,0),(1,1,0),(0,1,1),(1,1,1)} \node[atomB=3.8mm] at \p {};
  \node[atom=3.2mm] at (0.5,1,0.5) {};
  \foreach \p in {(0,0.5,0.5),(1,0.5,0.5),(0.5,0.5,0),(0.5,0.5,1)} \node[atom=3.2mm] at \p {};
  \foreach \p in {(0,0,0),(1,0,0),(0,0,1),(1,0,1)} \node[atomB=3.8mm] at \p {};
  \node[atom=3.2mm] at (0.5,0,0.5) {};
\end{tikzpicture}
\caption{Cu$_3$Au}
\end{subfigure}\hfill
\begin{subfigure}[b]{0.24\textwidth}\centering
\begin{tikzpicture}[cubeproj, x={(1cm,0cm)}, y={(0.42cm,0.30cm)}, z={(0cm,0.93cm)}, scale=1.55]
  \CubeEdges
  \foreach \p in {(0,1,0),(1,1,0),(0,1,1),(1,1,1)} \node[atomB=3.8mm] at \p {};
  \node[atom=3.2mm] at (0.5,1,0.5) {};
  \foreach \p in {(0,0.5,0.5),(1,0.5,0.5)} \node[atom=3.2mm] at \p {};
  \foreach \p in {(0.5,0.5,0),(0.5,0.5,1)} \node[atomB=3.8mm] at \p {};
  \foreach \p in {(0,0,0),(1,0,0),(0,0,1),(1,0,1)} \node[atomB=3.8mm] at \p {};
  \node[atom=3.2mm] at (0.5,0,0.5) {};
\end{tikzpicture}
\caption{CuAu\,I}
\end{subfigure}\hfill
\begin{subfigure}[b]{0.24\textwidth}\centering
\begin{tikzpicture}[cubeproj, scale=1.55]
  \CubeEdges
  \foreach \p in {(0,1,0),(1,1,0),(0,1,1),(1,1,1)} \node[atom=3.6mm] at \p {};
  \node[atomI=3.6mm] at (0.5,0.5,0.5) {};
  \foreach \p in {(0,0,0),(1,0,0),(0,0,1),(1,0,1)} \node[atom=3.6mm] at \p {};
\end{tikzpicture}
\caption{CuZn (B2)}
\end{subfigure}\hfill
\begin{subfigure}[b]{0.24\textwidth}\centering
\begin{tikzpicture}[cubeproj, scale=0.78]
  \draw[atomLine!50, dashed] (0,0,0)--(0,2,0)--(2,2,0) (0,2,0)--(0,2,2);
  \draw[atomLine] (0,0,0)--(2,0,0)--(2,0,2)--(0,0,2)--cycle (2,0,0)--(2,2,0)--(2,2,2)--(2,0,2) (0,0,2)--(0,2,2)--(2,2,2);
  \foreach \x in {0,1,2}{\foreach \y in {2,1,0}{\foreach \z in {0,1,2}{\node[atom=2.2mm] at (\x,\y,\z) {};}}}
  \foreach \x in {0.5,1.5}{\foreach \y in {1.5,0.5}{\foreach \z in {0.5,1.5}{
    \pgfmathtruncatemacro{\pp}{mod(round(\x+\y+\z-0.5),2)}
    \ifnum\pp=1 \node[atom=2.2mm] at (\x,\y,\z) {};\else \node[atomI=2.6mm] at (\x,\y,\z) {};\fi}}}
\end{tikzpicture}
\caption{Fe$_3$Al (D0$_3$)}
\end{subfigure}
\caption{Сверхструктуры: \textit{а}~— Cu$_3$Au (Au оранжевые~— в вершинах,
Cu~— в центрах граней); \textit{б}~— CuAu\,I (слои Au и Cu вдоль оси $c$);
\textit{в}~— CuZn (атомы разного сорта в вершинах и в центре куба);
\textit{г}~— Fe$_3$Al в удвоенной ОЦК-ячейке (Fe~— голубые, Al~— красные в
позициях $Y$)}
\label{fig:superstructures}
\end{figure}

\subsection{Фуллерены, фуллериты, квазикристаллы}\label{ssec:fullerenes}

\paragraph{Фуллерены}
При лазерном испарении графита в струе гелия образуются устойчивые
молекулы из большого чётного числа атомов углерода, чаще всего C$_{60}$ и
C$_{70}$. Форма молекулы C$_{60}$ напоминает купола архитектора Р.~Фуллера~—
отсюда название \term{фуллерены}. Атомы C$_{60}$ сидят в вершинах
\emph{усечённого икосаэдра} из 20 шестиугольников и 12 правильных
пятиугольников (рис.~\ref{fig:c60}); симметрия молекулы икосаэдрическая,
оси 5-го порядка проходят через центры пятиугольников. Все атомы лежат на
сфере радиусом $R=0{,}357$~нм; валентные электроны распределены почти
равномерно по сферическому слою толщиной около $8a_0$ ($a_0=0{,}0529$~нм~—
боровский радиус), а внутри есть полость радиусом около $2a_0$, почти
свободная от электронов. В неё можно поместить атом~— фуллерены
рассматривают как «контейнеры» для отдельных атомов и молекул.

\begin{figure}[ht]
\centering
\begin{subfigure}[b]{0.42\textwidth}\centering
\tikzset{c60front/.style={atomLine, line width=0.7pt},
  c60back/.style={atomLine!30, line width=0.4pt},
  c60pent/.style={fmAmber!45}}
\begin{tikzpicture}[scale=1.25]
  \fill[fmBlue!8] (0,0) circle (1.25);
  % создано tools/gen_cryst.py: молекула C60 (передние рёбра — сплошные)
\fill[c60pent] (0.009,-1.110) -- (0.451,-1.110) -- (0.757,-0.803) -- (0.506,-0.613) -- (0.043,-0.803) -- cycle;
\fill[c60pent] (-0.793,0.847) -- (-0.570,0.692) -- (-0.107,0.882) -- (-0.045,1.154) -- (-0.468,1.132) -- cycle;
\fill[c60pent] (0.188,0.148) -- (0.578,-0.137) -- (0.903,0.148) -- (0.713,0.610) -- (0.272,0.610) -- cycle;
\fill[c60pent] (-0.770,-0.539) -- (-0.347,-0.517) -- (-0.274,-0.042) -- (-0.653,0.230) -- (-0.960,-0.077) -- cycle;
\draw[c60back] (-0.188,-0.148) -- (0.274,0.042);
\draw[c60back] (-0.578,0.137) -- (-0.188,-0.148);
\draw[c60back] (0.274,0.042) -- (0.347,0.517);
\draw[c60back] (-0.272,-0.610) -- (-0.188,-0.148);
\draw[c60back] (0.274,0.042) -- (0.653,-0.230);
\draw[c60back] (-0.578,0.137) -- (-0.506,0.613);
\draw[c60back] (-0.043,0.803) -- (0.347,0.517);
\draw[c60back] (-0.903,-0.148) -- (-0.578,0.137);
\draw[c60back] (-0.272,-0.610) -- (0.107,-0.882);
\draw[c60back] (-0.506,0.613) -- (-0.043,0.803);
\draw[c60back] (0.570,-0.692) -- (0.653,-0.230);
\draw[c60back] (0.347,0.517) -- (0.770,0.539);
\draw[c60back] (-0.713,-0.610) -- (-0.272,-0.610);
\draw[c60back] (0.653,-0.230) -- (0.960,0.077);
\draw[c60back] (0.107,-0.882) -- (0.570,-0.692);
\draw[c60back] (-0.903,-0.148) -- (-0.713,-0.610);
\draw[c60back] (0.770,0.539) -- (0.960,0.077);
\draw[c60back] (-0.757,0.803) -- (-0.506,0.613);
\draw[c60back] (-0.009,1.110) -- (-0.043,0.803);
\draw[c60back] (0.045,-1.154) -- (0.107,-0.882);
\draw[c60back] (0.570,-0.692) -- (0.793,-0.847);
\draw[c60back] (-1.154,0.042) -- (-0.903,-0.148);
\draw[c60back] (0.804,0.847) -- (0.770,0.539);
\draw[c60back] (-0.776,-0.882) -- (-0.713,-0.610);
\draw[c60back] (0.960,0.077) -- (1.183,-0.077);
\draw[c60back] (-0.757,0.803) -- (-0.451,1.110);
\draw[c60back] (-1.082,0.517) -- (-0.757,0.803);
\draw[c60back] (-0.451,1.110) -- (-0.009,1.110);
\draw[c60back] (-0.009,1.110) -- (0.414,1.132);
\draw[c60back] (-1.154,0.042) -- (-1.082,0.517);
\draw[c60back] (0.414,1.132) -- (0.804,0.847);
\draw[c60back] (-0.397,-1.154) -- (0.045,-1.154);
\draw[c60back] (0.045,-1.154) -- (0.468,-1.132);
\draw[c60back] (0.793,-0.847) -- (1.100,-0.539);
\draw[c60back] (0.468,-1.132) -- (0.793,-0.847);
\draw[c60back] (-0.776,-0.882) -- (-0.397,-1.154);
\draw[c60back] (1.100,-0.539) -- (1.183,-0.077);
\draw[c60back] (-1.217,-0.230) -- (-1.154,0.042);
\draw[c60back] (0.804,0.847) -- (1.028,0.692);
\draw[c60back] (-1.028,-0.692) -- (-0.776,-0.882);
\draw[c60back] (1.217,0.230) -- (1.183,-0.077);
\draw[c60back] (-0.468,1.132) -- (-0.451,1.110);
\draw[c60back] (-1.100,0.539) -- (-1.082,0.517);
\draw[c60back] (0.397,1.154) -- (0.414,1.132);
\draw[c60back] (-1.217,-0.230) -- (-1.028,-0.692);
\draw[c60front] (1.028,0.692) -- (1.217,0.230);
\draw[c60front] (-0.414,-1.132) -- (-0.397,-1.154);
\draw[c60front] (1.082,-0.517) -- (1.100,-0.539);
\draw[c60front] (0.451,-1.110) -- (0.468,-1.132);
\draw[c60front] (-1.183,0.077) -- (-1.217,-0.230);
\draw[c60front] (0.776,0.882) -- (1.028,0.692);
\draw[c60front] (-1.028,-0.692) -- (-0.804,-0.847);
\draw[c60front] (1.154,-0.042) -- (1.217,0.230);
\draw[c60front] (-1.183,0.077) -- (-1.100,0.539);
\draw[c60front] (0.397,1.154) -- (0.776,0.882);
\draw[c60front] (-0.793,0.847) -- (-0.468,1.132);
\draw[c60front] (-1.100,0.539) -- (-0.793,0.847);
\draw[c60front] (-0.468,1.132) -- (-0.045,1.154);
\draw[c60front] (-0.045,1.154) -- (0.397,1.154);
\draw[c60front] (-0.804,-0.847) -- (-0.414,-1.132);
\draw[c60front] (1.082,-0.517) -- (1.154,-0.042);
\draw[c60front] (-0.414,-1.132) -- (0.009,-1.110);
\draw[c60front] (0.009,-1.110) -- (0.451,-1.110);
\draw[c60front] (0.757,-0.803) -- (1.082,-0.517);
\draw[c60front] (0.451,-1.110) -- (0.757,-0.803);
\draw[c60front] (-1.183,0.077) -- (-0.960,-0.077);
\draw[c60front] (0.713,0.610) -- (0.776,0.882);
\draw[c60front] (-0.770,-0.539) -- (-0.804,-0.847);
\draw[c60front] (0.903,0.148) -- (1.154,-0.042);
\draw[c60front] (-0.793,0.847) -- (-0.570,0.692);
\draw[c60front] (-0.107,0.882) -- (-0.045,1.154);
\draw[c60front] (0.043,-0.803) -- (0.009,-1.110);
\draw[c60front] (0.506,-0.613) -- (0.757,-0.803);
\draw[c60front] (-0.960,-0.077) -- (-0.770,-0.539);
\draw[c60front] (0.713,0.610) -- (0.903,0.148);
\draw[c60front] (-0.570,0.692) -- (-0.107,0.882);
\draw[c60front] (-0.960,-0.077) -- (-0.653,0.230);
\draw[c60front] (0.272,0.610) -- (0.713,0.610);
\draw[c60front] (-0.770,-0.539) -- (-0.347,-0.517);
\draw[c60front] (-0.653,0.230) -- (-0.570,0.692);
\draw[c60front] (0.043,-0.803) -- (0.506,-0.613);
\draw[c60front] (-0.107,0.882) -- (0.272,0.610);
\draw[c60front] (0.578,-0.137) -- (0.903,0.148);
\draw[c60front] (-0.347,-0.517) -- (0.043,-0.803);
\draw[c60front] (0.506,-0.613) -- (0.578,-0.137);
\draw[c60front] (-0.653,0.230) -- (-0.274,-0.042);
\draw[c60front] (0.188,0.148) -- (0.272,0.610);
\draw[c60front] (-0.347,-0.517) -- (-0.274,-0.042);
\draw[c60front] (0.188,0.148) -- (0.578,-0.137);
\draw[c60front] (-0.274,-0.042) -- (0.188,0.148);
\fill[atomLine] (1.028,0.692) circle (0.03);
\fill[atomLine] (1.217,0.230) circle (0.03);
\fill[atomLine] (-0.468,1.132) circle (0.03);
\fill[atomLine] (-1.100,0.539) circle (0.03);
\fill[atomLine] (0.397,1.154) circle (0.03);
\fill[atomLine] (-0.414,-1.132) circle (0.03);
\fill[atomLine] (1.082,-0.517) circle (0.03);
\fill[atomLine] (0.451,-1.110) circle (0.03);
\fill[atomLine] (-1.183,0.077) circle (0.03);
\fill[atomLine] (0.776,0.882) circle (0.03);
\fill[atomLine] (-0.804,-0.847) circle (0.03);
\fill[atomLine] (-0.793,0.847) circle (0.03);
\fill[atomLine] (1.154,-0.042) circle (0.03);
\fill[atomLine] (-0.045,1.154) circle (0.03);
\fill[atomLine] (0.009,-1.110) circle (0.03);
\fill[atomLine] (0.757,-0.803) circle (0.03);
\fill[atomLine] (-0.960,-0.077) circle (0.03);
\fill[atomLine] (-0.770,-0.539) circle (0.03);
\fill[atomLine] (0.713,0.610) circle (0.03);
\fill[atomLine] (0.903,0.148) circle (0.03);
\fill[atomLine] (-0.570,0.692) circle (0.03);
\fill[atomLine] (-0.107,0.882) circle (0.03);
\fill[atomLine] (0.043,-0.803) circle (0.03);
\fill[atomLine] (0.506,-0.613) circle (0.03);
\fill[atomLine] (-0.653,0.230) circle (0.03);
\fill[atomLine] (0.272,0.610) circle (0.03);
\fill[atomLine] (-0.347,-0.517) circle (0.03);
\fill[atomLine] (0.578,-0.137) circle (0.03);
\fill[atomLine] (-0.274,-0.042) circle (0.03);
\fill[atomLine] (0.188,0.148) circle (0.03);

\end{tikzpicture}
\caption{}
\end{subfigure}\hfill
\begin{subfigure}[b]{0.52\textwidth}\centering
\tikzset{pthin/.style={fill=fmBlue!45}, pthick/.style={fill=fmAmber!35},
  pedge/.style={atomLine, line width=0.25pt}}
\begin{tikzpicture}[scale=0.92]
  \clip (0,0) circle (3.3);
  % создано tools/gen_cryst.py: мозаика Пенроуза
\fill[pthin] (2.762,0.401) -- (2.470,0.307) -- (2.470,0.496) -- cycle;
\fill[pthick] (2.470,0.307) -- (2.762,0.401) -- (2.470,0.000) -- cycle;
\fill[pthick] (2.179,0.401) -- (1.998,0.649) -- (2.470,0.496) -- cycle;
\fill[pthick] (2.470,0.307) -- (2.179,0.401) -- (2.470,0.000) -- cycle;
\fill[pthin] (2.179,0.401) -- (2.470,0.307) -- (2.470,0.496) -- cycle;
\fill[pthick] (2.942,0.649) -- (2.762,0.401) -- (2.762,0.897) -- cycle;
\fill[pthick] (2.942,0.956) -- (2.942,0.649) -- (3.234,1.051) -- cycle;
\fill[pthin] (2.942,0.649) -- (2.942,0.956) -- (2.762,0.897) -- cycle;
\fill[pthick] (2.290,0.744) -- (1.998,0.649) -- (2.470,0.496) -- cycle;
\fill[pthick] (2.582,0.649) -- (2.290,0.744) -- (2.762,0.897) -- cycle;
\fill[pthin] (2.290,0.744) -- (2.582,0.649) -- (2.470,0.496) -- cycle;
\fill[pthin] (2.762,0.401) -- (2.582,0.649) -- (2.470,0.496) -- cycle;
\fill[pthick] (2.582,0.649) -- (2.762,0.401) -- (2.762,0.897) -- cycle;
\fill[pthick] (2.942,-0.153) -- (3.234,-0.248) -- (2.762,-0.401) -- cycle;
\fill[pthick] (2.650,-0.248) -- (2.942,-0.153) -- (2.470,0.000) -- cycle;
\fill[pthin] (2.942,-0.153) -- (2.650,-0.248) -- (2.762,-0.401) -- cycle;
\fill[pthick] (3.053,-0.803) -- (3.234,-1.051) -- (3.234,-0.555) -- cycle;
\fill[pthick] (3.053,-0.496) -- (3.053,-0.803) -- (2.762,-0.401) -- cycle;
\fill[pthin] (3.053,-0.803) -- (3.053,-0.496) -- (3.234,-0.555) -- cycle;
\fill[pthin] (3.234,-0.248) -- (3.053,-0.496) -- (3.234,-0.555) -- cycle;
\fill[pthick] (3.053,-0.496) -- (3.234,-0.248) -- (2.762,-0.401) -- cycle;
\fill[pthick] (2.942,0.649) -- (2.762,0.401) -- (3.234,0.555) -- cycle;
\fill[pthick] (3.234,0.744) -- (2.942,0.649) -- (3.234,1.051) -- cycle;
\fill[pthin] (2.942,0.649) -- (3.234,0.744) -- (3.234,0.555) -- cycle;
\fill[pthick] (3.234,0.059) -- (3.234,-0.248) -- (2.942,0.153) -- cycle;
\fill[pthick] (3.053,0.307) -- (3.234,0.059) -- (3.234,0.555) -- cycle;
\fill[pthin] (3.234,0.059) -- (3.053,0.307) -- (2.942,0.153) -- cycle;
\fill[pthin] (2.762,0.401) -- (3.053,0.307) -- (2.942,0.153) -- cycle;
\fill[pthick] (3.053,0.307) -- (2.762,0.401) -- (3.234,0.555) -- cycle;
\fill[pthin] (2.762,0.401) -- (2.762,0.095) -- (2.942,0.153) -- cycle;
\fill[pthick] (2.762,0.095) -- (2.762,0.401) -- (2.470,0.000) -- cycle;
\fill[pthick] (2.942,-0.153) -- (3.234,-0.248) -- (2.942,0.153) -- cycle;
\fill[pthick] (2.762,0.095) -- (2.942,-0.153) -- (2.470,0.000) -- cycle;
\fill[pthin] (2.942,-0.153) -- (2.762,0.095) -- (2.942,0.153) -- cycle;
\fill[pthick] (1.235,0.095) -- (1.235,0.401) -- (1.527,0.000) -- cycle;
\fill[pthick] (1.415,-0.153) -- (1.235,0.095) -- (1.235,-0.401) -- cycle;
\fill[pthin] (1.235,0.095) -- (1.415,-0.153) -- (1.527,0.000) -- cycle;
\fill[pthick] (1.818,0.401) -- (1.998,0.649) -- (1.527,0.496) -- cycle;
\fill[pthick] (1.527,0.307) -- (1.818,0.401) -- (1.527,0.000) -- cycle;
\fill[pthin] (1.818,0.401) -- (1.527,0.307) -- (1.527,0.496) -- cycle;
\fill[pthin] (1.235,0.401) -- (1.527,0.307) -- (1.527,0.496) -- cycle;
\fill[pthick] (1.527,0.307) -- (1.235,0.401) -- (1.527,0.000) -- cycle;
\fill[pthick] (0.472,-0.153) -- (0.763,-0.248) -- (0.472,0.153) -- cycle;
\fill[pthick] (0.292,0.095) -- (0.472,-0.153) -- (0.000,0.000) -- cycle;
\fill[pthin] (0.472,-0.153) -- (0.292,0.095) -- (0.472,0.153) -- cycle;
\fill[pthick] (0.944,0.307) -- (1.235,0.401) -- (0.944,0.000) -- cycle;
\fill[pthick] (0.763,0.059) -- (0.944,0.307) -- (0.472,0.153) -- cycle;
\fill[pthin] (0.944,0.307) -- (0.763,0.059) -- (0.944,0.000) -- cycle;
\fill[pthin] (0.763,-0.248) -- (0.763,0.059) -- (0.944,0.000) -- cycle;
\fill[pthick] (0.763,0.059) -- (0.763,-0.248) -- (0.472,0.153) -- cycle;
\fill[pthin] (0.763,-0.248) -- (1.055,-0.153) -- (0.944,0.000) -- cycle;
\fill[pthick] (1.055,-0.153) -- (0.763,-0.248) -- (1.235,-0.401) -- cycle;
\fill[pthick] (1.235,0.095) -- (1.235,0.401) -- (0.944,0.000) -- cycle;
\fill[pthick] (1.055,-0.153) -- (1.235,0.095) -- (1.235,-0.401) -- cycle;
\fill[pthin] (1.235,0.095) -- (1.055,-0.153) -- (0.944,0.000) -- cycle;
\fill[pthick] (2.470,-0.496) -- (2.470,-0.803) -- (2.762,-0.401) -- cycle;
\fill[pthick] (2.650,-0.248) -- (2.470,-0.496) -- (2.470,0.000) -- cycle;
\fill[pthin] (2.470,-0.496) -- (2.650,-0.248) -- (2.762,-0.401) -- cycle;
\fill[pthick] (3.053,-0.803) -- (3.234,-1.051) -- (2.762,-0.897) -- cycle;
\fill[pthick] (2.762,-0.708) -- (3.053,-0.803) -- (2.762,-0.401) -- cycle;
\fill[pthin] (3.053,-0.803) -- (2.762,-0.708) -- (2.762,-0.897) -- cycle;
\fill[pthin] (2.470,-0.803) -- (2.762,-0.708) -- (2.762,-0.897) -- cycle;
\fill[pthick] (2.762,-0.708) -- (2.470,-0.803) -- (2.762,-0.401) -- cycle;
\fill[pthick] (1.707,-0.248) -- (1.998,-0.153) -- (1.707,-0.555) -- cycle;
\fill[pthick] (1.527,-0.496) -- (1.707,-0.248) -- (1.235,-0.401) -- cycle;
\fill[pthin] (1.707,-0.248) -- (1.527,-0.496) -- (1.707,-0.555) -- cycle;
\fill[pthick] (2.179,-0.708) -- (2.470,-0.803) -- (2.179,-0.401) -- cycle;
\fill[pthick] (1.998,-0.460) -- (2.179,-0.708) -- (1.707,-0.555) -- cycle;
\fill[pthin] (2.179,-0.708) -- (1.998,-0.460) -- (2.179,-0.401) -- cycle;
\fill[pthin] (1.998,-0.153) -- (1.998,-0.460) -- (2.179,-0.401) -- cycle;
\fill[pthick] (1.998,-0.460) -- (1.998,-0.153) -- (1.707,-0.555) -- cycle;
\fill[pthin] (1.998,-0.153) -- (2.290,-0.248) -- (2.179,-0.401) -- cycle;
\fill[pthick] (2.290,-0.248) -- (1.998,-0.153) -- (2.470,0.000) -- cycle;
\fill[pthick] (2.470,-0.496) -- (2.470,-0.803) -- (2.179,-0.401) -- cycle;
\fill[pthick] (2.290,-0.248) -- (2.470,-0.496) -- (2.470,0.000) -- cycle;
\fill[pthin] (2.470,-0.496) -- (2.290,-0.248) -- (2.179,-0.401) -- cycle;
\fill[pthin] (1.998,-0.153) -- (2.179,0.095) -- (1.998,0.153) -- cycle;
\fill[pthick] (2.179,0.095) -- (1.998,-0.153) -- (2.470,0.000) -- cycle;
\fill[pthick] (2.179,0.401) -- (1.998,0.649) -- (1.998,0.153) -- cycle;
\fill[pthick] (2.179,0.095) -- (2.179,0.401) -- (2.470,0.000) -- cycle;
\fill[pthin] (2.179,0.401) -- (2.179,0.095) -- (1.998,0.153) -- cycle;
\fill[pthick] (1.707,-0.248) -- (1.998,-0.153) -- (1.527,0.000) -- cycle;
\fill[pthick] (1.415,-0.153) -- (1.707,-0.248) -- (1.235,-0.401) -- cycle;
\fill[pthin] (1.707,-0.248) -- (1.415,-0.153) -- (1.527,0.000) -- cycle;
\fill[pthick] (1.818,0.401) -- (1.998,0.649) -- (1.998,0.153) -- cycle;
\fill[pthick] (1.818,0.095) -- (1.818,0.401) -- (1.527,0.000) -- cycle;
\fill[pthin] (1.818,0.401) -- (1.818,0.095) -- (1.998,0.153) -- cycle;
\fill[pthin] (1.998,-0.153) -- (1.818,0.095) -- (1.998,0.153) -- cycle;
\fill[pthick] (1.818,0.095) -- (1.998,-0.153) -- (1.527,0.000) -- cycle;
\fill[pthin] (2.470,1.299) -- (2.179,1.204) -- (2.290,1.051) -- cycle;
\fill[pthick] (2.179,1.204) -- (2.470,1.299) -- (1.998,1.452) -- cycle;
\fill[pthick] (1.998,0.956) -- (1.998,0.649) -- (2.290,1.051) -- cycle;
\fill[pthick] (2.179,1.204) -- (1.998,0.956) -- (1.998,1.452) -- cycle;
\fill[pthin] (1.998,0.956) -- (2.179,1.204) -- (2.290,1.051) -- cycle;
\fill[pthick] (2.762,1.204) -- (2.470,1.299) -- (2.762,0.897) -- cycle;
\fill[pthick] (2.942,0.956) -- (2.762,1.204) -- (3.234,1.051) -- cycle;
\fill[pthin] (2.762,1.204) -- (2.942,0.956) -- (2.762,0.897) -- cycle;
\fill[pthick] (2.290,0.744) -- (1.998,0.649) -- (2.290,1.051) -- cycle;
\fill[pthick] (2.470,0.992) -- (2.290,0.744) -- (2.762,0.897) -- cycle;
\fill[pthin] (2.290,0.744) -- (2.470,0.992) -- (2.290,1.051) -- cycle;
\fill[pthin] (2.470,1.299) -- (2.470,0.992) -- (2.290,1.051) -- cycle;
\fill[pthick] (2.470,0.992) -- (2.470,1.299) -- (2.762,0.897) -- cycle;
\fill[pthick] (2.290,1.853) -- (2.470,2.101) -- (1.998,1.948) -- cycle;
\fill[pthick] (1.998,1.759) -- (2.290,1.853) -- (1.998,1.452) -- cycle;
\fill[pthin] (2.290,1.853) -- (1.998,1.759) -- (1.998,1.948) -- cycle;
\fill[pthick] (1.998,2.444) -- (1.998,2.751) -- (2.290,2.349) -- cycle;
\fill[pthick] (2.179,2.196) -- (1.998,2.444) -- (1.998,1.948) -- cycle;
\fill[pthin] (1.998,2.444) -- (2.179,2.196) -- (2.290,2.349) -- cycle;
\fill[pthin] (2.470,2.101) -- (2.179,2.196) -- (2.290,2.349) -- cycle;
\fill[pthick] (2.179,2.196) -- (2.470,2.101) -- (1.998,1.948) -- cycle;
\fill[pthick] (2.762,1.204) -- (2.470,1.299) -- (2.942,1.452) -- cycle;
\fill[pthick] (3.053,1.299) -- (2.762,1.204) -- (3.234,1.051) -- cycle;
\fill[pthin] (2.762,1.204) -- (3.053,1.299) -- (2.942,1.452) -- cycle;
\fill[pthick] (2.650,1.853) -- (2.470,2.101) -- (2.470,1.605) -- cycle;
\fill[pthick] (2.650,1.547) -- (2.650,1.853) -- (2.942,1.452) -- cycle;
\fill[pthin] (2.650,1.853) -- (2.650,1.547) -- (2.470,1.605) -- cycle;
\fill[pthin] (2.470,1.299) -- (2.650,1.547) -- (2.470,1.605) -- cycle;
\fill[pthick] (2.650,1.547) -- (2.470,1.299) -- (2.942,1.452) -- cycle;
\fill[pthin] (2.470,1.299) -- (2.290,1.547) -- (2.470,1.605) -- cycle;
\fill[pthick] (2.290,1.547) -- (2.470,1.299) -- (1.998,1.452) -- cycle;
\fill[pthick] (2.290,1.853) -- (2.470,2.101) -- (2.470,1.605) -- cycle;
\fill[pthick] (2.290,1.547) -- (2.290,1.853) -- (1.998,1.452) -- cycle;
\fill[pthin] (2.290,1.853) -- (2.290,1.547) -- (2.470,1.605) -- cycle;
\fill[pthick] (1.055,0.649) -- (1.235,0.401) -- (1.235,0.897) -- cycle;
\fill[pthick] (1.055,0.956) -- (1.055,0.649) -- (0.763,1.051) -- cycle;
\fill[pthin] (1.055,0.649) -- (1.055,0.956) -- (1.235,0.897) -- cycle;
\fill[pthick] (1.707,0.744) -- (1.998,0.649) -- (1.527,0.496) -- cycle;
\fill[pthick] (1.415,0.649) -- (1.707,0.744) -- (1.235,0.897) -- cycle;
\fill[pthin] (1.707,0.744) -- (1.415,0.649) -- (1.527,0.496) -- cycle;
\fill[pthin] (1.235,0.401) -- (1.415,0.649) -- (1.527,0.496) -- cycle;
\fill[pthick] (1.415,0.649) -- (1.235,0.401) -- (1.235,0.897) -- cycle;
\fill[pthick] (0.292,0.401) -- (0.472,0.649) -- (0.472,0.153) -- cycle;
\fill[pthick] (0.292,0.095) -- (0.292,0.401) -- (0.000,0.000) -- cycle;
\fill[pthin] (0.292,0.401) -- (0.292,0.095) -- (0.472,0.153) -- cycle;
\fill[pthick] (0.944,0.307) -- (1.235,0.401) -- (0.763,0.555) -- cycle;
\fill[pthick] (0.652,0.401) -- (0.944,0.307) -- (0.472,0.153) -- cycle;
\fill[pthin] (0.944,0.307) -- (0.652,0.401) -- (0.763,0.555) -- cycle;
\fill[pthin] (0.472,0.649) -- (0.652,0.401) -- (0.763,0.555) -- cycle;
\fill[pthick] (0.652,0.401) -- (0.472,0.649) -- (0.472,0.153) -- cycle;
\fill[pthin] (0.472,0.649) -- (0.763,0.744) -- (0.763,0.555) -- cycle;
\fill[pthick] (0.763,0.744) -- (0.472,0.649) -- (0.763,1.051) -- cycle;
\fill[pthick] (1.055,0.649) -- (1.235,0.401) -- (0.763,0.555) -- cycle;
\fill[pthick] (0.763,0.744) -- (1.055,0.649) -- (0.763,1.051) -- cycle;
\fill[pthin] (1.055,0.649) -- (0.763,0.744) -- (0.763,0.555) -- cycle;
\fill[pthick] (1.707,1.853) -- (1.527,2.101) -- (1.998,1.948) -- cycle;
\fill[pthick] (1.998,1.759) -- (1.707,1.853) -- (1.998,1.452) -- cycle;
\fill[pthin] (1.707,1.853) -- (1.998,1.759) -- (1.998,1.948) -- cycle;
\fill[pthick] (1.998,2.444) -- (1.998,2.751) -- (1.707,2.349) -- cycle;
\fill[pthick] (1.818,2.196) -- (1.998,2.444) -- (1.998,1.948) -- cycle;
\fill[pthin] (1.998,2.444) -- (1.818,2.196) -- (1.707,2.349) -- cycle;
\fill[pthin] (1.527,2.101) -- (1.818,2.196) -- (1.707,2.349) -- cycle;
\fill[pthick] (1.818,2.196) -- (1.527,2.101) -- (1.998,1.948) -- cycle;
\fill[pthick] (1.235,1.204) -- (1.527,1.299) -- (1.055,1.452) -- cycle;
\fill[pthick] (0.944,1.299) -- (1.235,1.204) -- (0.763,1.051) -- cycle;
\fill[pthin] (1.235,1.204) -- (0.944,1.299) -- (1.055,1.452) -- cycle;
\fill[pthick] (1.346,1.853) -- (1.527,2.101) -- (1.527,1.605) -- cycle;
\fill[pthick] (1.346,1.547) -- (1.346,1.853) -- (1.055,1.452) -- cycle;
\fill[pthin] (1.346,1.853) -- (1.346,1.547) -- (1.527,1.605) -- cycle;
\fill[pthin] (1.527,1.299) -- (1.346,1.547) -- (1.527,1.605) -- cycle;
\fill[pthick] (1.346,1.547) -- (1.527,1.299) -- (1.055,1.452) -- cycle;
\fill[pthin] (1.527,1.299) -- (1.707,1.547) -- (1.527,1.605) -- cycle;
\fill[pthick] (1.707,1.547) -- (1.527,1.299) -- (1.998,1.452) -- cycle;
\fill[pthick] (1.707,1.853) -- (1.527,2.101) -- (1.527,1.605) -- cycle;
\fill[pthick] (1.707,1.547) -- (1.707,1.853) -- (1.998,1.452) -- cycle;
\fill[pthin] (1.707,1.853) -- (1.707,1.547) -- (1.527,1.605) -- cycle;
\fill[pthin] (1.527,1.299) -- (1.818,1.204) -- (1.707,1.051) -- cycle;
\fill[pthick] (1.818,1.204) -- (1.527,1.299) -- (1.998,1.452) -- cycle;
\fill[pthick] (1.998,0.956) -- (1.998,0.649) -- (1.707,1.051) -- cycle;
\fill[pthick] (1.818,1.204) -- (1.998,0.956) -- (1.998,1.452) -- cycle;
\fill[pthin] (1.998,0.956) -- (1.818,1.204) -- (1.707,1.051) -- cycle;
\fill[pthick] (1.235,1.204) -- (1.527,1.299) -- (1.235,0.897) -- cycle;
\fill[pthick] (1.055,0.956) -- (1.235,1.204) -- (0.763,1.051) -- cycle;
\fill[pthin] (1.235,1.204) -- (1.055,0.956) -- (1.235,0.897) -- cycle;
\fill[pthick] (1.707,0.744) -- (1.998,0.649) -- (1.707,1.051) -- cycle;
\fill[pthick] (1.527,0.992) -- (1.707,0.744) -- (1.235,0.897) -- cycle;
\fill[pthin] (1.707,0.744) -- (1.527,0.992) -- (1.707,1.051) -- cycle;
\fill[pthin] (1.527,1.299) -- (1.527,0.992) -- (1.707,1.051) -- cycle;
\fill[pthick] (1.527,0.992) -- (1.527,1.299) -- (1.235,0.897) -- cycle;
\fill[pthin] (0.472,2.751) -- (0.472,2.444) -- (0.292,2.503) -- cycle;
\fill[pthick] (0.472,2.444) -- (0.472,2.751) -- (0.763,2.349) -- cycle;
\fill[pthick] (0.292,2.196) -- (0.000,2.101) -- (0.292,2.503) -- cycle;
\fill[pthick] (0.472,2.444) -- (0.292,2.196) -- (0.763,2.349) -- cycle;
\fill[pthin] (0.292,2.196) -- (0.472,2.444) -- (0.292,2.503) -- cycle;
\fill[pthick] (0.292,2.999) -- (0.472,2.751) -- (0.000,2.904) -- cycle;
\fill[pthick] (0.000,3.093) -- (0.292,2.999) -- (0.000,3.400) -- cycle;
\fill[pthin] (0.292,2.999) -- (0.000,3.093) -- (0.000,2.904) -- cycle;
\fill[pthick] (0.000,2.408) -- (0.000,2.101) -- (0.292,2.503) -- cycle;
\fill[pthick] (0.180,2.656) -- (0.000,2.408) -- (0.000,2.904) -- cycle;
\fill[pthin] (0.000,2.408) -- (0.180,2.656) -- (0.292,2.503) -- cycle;
\fill[pthin] (0.472,2.751) -- (0.180,2.656) -- (0.292,2.503) -- cycle;
\fill[pthick] (0.180,2.656) -- (0.472,2.751) -- (0.000,2.904) -- cycle;
\fill[pthick] (1.055,2.751) -- (1.235,2.999) -- (1.235,2.503) -- cycle;
\fill[pthick] (1.055,2.444) -- (1.055,2.751) -- (0.763,2.349) -- cycle;
\fill[pthin] (1.055,2.751) -- (1.055,2.444) -- (1.235,2.503) -- cycle;
\fill[pthick] (1.707,2.656) -- (1.998,2.751) -- (1.527,2.904) -- cycle;
\fill[pthick] (1.415,2.751) -- (1.707,2.656) -- (1.235,2.503) -- cycle;
\fill[pthin] (1.707,2.656) -- (1.415,2.751) -- (1.527,2.904) -- cycle;
\fill[pthin] (1.235,2.999) -- (1.415,2.751) -- (1.527,2.904) -- cycle;
\fill[pthick] (1.415,2.751) -- (1.235,2.999) -- (1.235,2.503) -- cycle;
\fill[pthick] (0.292,2.999) -- (0.472,2.751) -- (0.472,3.247) -- cycle;
\fill[pthick] (0.292,3.305) -- (0.292,2.999) -- (0.000,3.400) -- cycle;
\fill[pthin] (0.292,2.999) -- (0.292,3.305) -- (0.472,3.247) -- cycle;
\fill[pthick] (0.944,3.093) -- (1.235,2.999) -- (0.763,2.845) -- cycle;
\fill[pthick] (0.652,2.999) -- (0.944,3.093) -- (0.472,3.247) -- cycle;
\fill[pthin] (0.944,3.093) -- (0.652,2.999) -- (0.763,2.845) -- cycle;
\fill[pthin] (0.472,2.751) -- (0.652,2.999) -- (0.763,2.845) -- cycle;
\fill[pthick] (0.652,2.999) -- (0.472,2.751) -- (0.472,3.247) -- cycle;
\fill[pthin] (0.472,2.751) -- (0.763,2.656) -- (0.763,2.845) -- cycle;
\fill[pthick] (0.763,2.656) -- (0.472,2.751) -- (0.763,2.349) -- cycle;
\fill[pthick] (1.055,2.751) -- (1.235,2.999) -- (0.763,2.845) -- cycle;
\fill[pthick] (0.763,2.656) -- (1.055,2.751) -- (0.763,2.349) -- cycle;
\fill[pthin] (1.055,2.751) -- (0.763,2.656) -- (0.763,2.845) -- cycle;
\fill[pthick] (0.292,1.204) -- (0.000,1.299) -- (0.472,1.452) -- cycle;
\fill[pthick] (0.583,1.299) -- (0.292,1.204) -- (0.763,1.051) -- cycle;
\fill[pthin] (0.292,1.204) -- (0.583,1.299) -- (0.472,1.452) -- cycle;
\fill[pthick] (0.180,1.853) -- (0.000,2.101) -- (0.000,1.605) -- cycle;
\fill[pthick] (0.180,1.547) -- (0.180,1.853) -- (0.472,1.452) -- cycle;
\fill[pthin] (0.180,1.853) -- (0.180,1.547) -- (0.000,1.605) -- cycle;
\fill[pthin] (0.000,1.299) -- (0.180,1.547) -- (0.000,1.605) -- cycle;
\fill[pthick] (0.180,1.547) -- (0.000,1.299) -- (0.472,1.452) -- cycle;
\fill[pthick] (0.292,0.401) -- (0.472,0.649) -- (0.000,0.496) -- cycle;
\fill[pthick] (0.000,0.307) -- (0.292,0.401) -- (0.000,0.000) -- cycle;
\fill[pthin] (0.292,0.401) -- (0.000,0.307) -- (0.000,0.496) -- cycle;
\fill[pthick] (0.000,0.992) -- (0.000,1.299) -- (0.292,0.897) -- cycle;
\fill[pthick] (0.180,0.744) -- (0.000,0.992) -- (0.000,0.496) -- cycle;
\fill[pthin] (0.000,0.992) -- (0.180,0.744) -- (0.292,0.897) -- cycle;
\fill[pthin] (0.472,0.649) -- (0.180,0.744) -- (0.292,0.897) -- cycle;
\fill[pthick] (0.180,0.744) -- (0.472,0.649) -- (0.000,0.496) -- cycle;
\fill[pthin] (0.472,0.649) -- (0.472,0.956) -- (0.292,0.897) -- cycle;
\fill[pthick] (0.472,0.956) -- (0.472,0.649) -- (0.763,1.051) -- cycle;
\fill[pthick] (0.292,1.204) -- (0.000,1.299) -- (0.292,0.897) -- cycle;
\fill[pthick] (0.472,0.956) -- (0.292,1.204) -- (0.763,1.051) -- cycle;
\fill[pthin] (0.292,1.204) -- (0.472,0.956) -- (0.292,0.897) -- cycle;
\fill[pthick] (1.235,2.196) -- (1.527,2.101) -- (1.235,2.503) -- cycle;
\fill[pthick] (1.055,2.444) -- (1.235,2.196) -- (0.763,2.349) -- cycle;
\fill[pthin] (1.235,2.196) -- (1.055,2.444) -- (1.235,2.503) -- cycle;
\fill[pthick] (1.707,2.656) -- (1.998,2.751) -- (1.707,2.349) -- cycle;
\fill[pthick] (1.527,2.408) -- (1.707,2.656) -- (1.235,2.503) -- cycle;
\fill[pthin] (1.707,2.656) -- (1.527,2.408) -- (1.707,2.349) -- cycle;
\fill[pthin] (1.527,2.101) -- (1.527,2.408) -- (1.707,2.349) -- cycle;
\fill[pthick] (1.527,2.408) -- (1.527,2.101) -- (1.235,2.503) -- cycle;
\fill[pthick] (0.763,1.547) -- (0.763,1.853) -- (1.055,1.452) -- cycle;
\fill[pthick] (0.944,1.299) -- (0.763,1.547) -- (0.763,1.051) -- cycle;
\fill[pthin] (0.763,1.547) -- (0.944,1.299) -- (1.055,1.452) -- cycle;
\fill[pthick] (1.346,1.853) -- (1.527,2.101) -- (1.055,1.948) -- cycle;
\fill[pthick] (1.055,1.759) -- (1.346,1.853) -- (1.055,1.452) -- cycle;
\fill[pthin] (1.346,1.853) -- (1.055,1.759) -- (1.055,1.948) -- cycle;
\fill[pthin] (0.763,1.853) -- (1.055,1.759) -- (1.055,1.948) -- cycle;
\fill[pthick] (1.055,1.759) -- (0.763,1.853) -- (1.055,1.452) -- cycle;
\fill[pthin] (0.763,1.853) -- (0.944,2.101) -- (1.055,1.948) -- cycle;
\fill[pthick] (0.944,2.101) -- (0.763,1.853) -- (0.763,2.349) -- cycle;
\fill[pthick] (1.235,2.196) -- (1.527,2.101) -- (1.055,1.948) -- cycle;
\fill[pthick] (0.944,2.101) -- (1.235,2.196) -- (0.763,2.349) -- cycle;
\fill[pthin] (1.235,2.196) -- (0.944,2.101) -- (1.055,1.948) -- cycle;
\fill[pthin] (0.763,1.853) -- (0.583,2.101) -- (0.472,1.948) -- cycle;
\fill[pthick] (0.583,2.101) -- (0.763,1.853) -- (0.763,2.349) -- cycle;
\fill[pthick] (0.292,2.196) -- (0.000,2.101) -- (0.472,1.948) -- cycle;
\fill[pthick] (0.583,2.101) -- (0.292,2.196) -- (0.763,2.349) -- cycle;
\fill[pthin] (0.292,2.196) -- (0.583,2.101) -- (0.472,1.948) -- cycle;
\fill[pthick] (0.763,1.547) -- (0.763,1.853) -- (0.472,1.452) -- cycle;
\fill[pthick] (0.583,1.299) -- (0.763,1.547) -- (0.763,1.051) -- cycle;
\fill[pthin] (0.763,1.547) -- (0.583,1.299) -- (0.472,1.452) -- cycle;
\fill[pthick] (0.180,1.853) -- (0.000,2.101) -- (0.472,1.948) -- cycle;
\fill[pthick] (0.472,1.759) -- (0.180,1.853) -- (0.472,1.452) -- cycle;
\fill[pthin] (0.180,1.853) -- (0.472,1.759) -- (0.472,1.948) -- cycle;
\fill[pthin] (0.763,1.853) -- (0.472,1.759) -- (0.472,1.948) -- cycle;
\fill[pthick] (0.472,1.759) -- (0.763,1.853) -- (0.472,1.452) -- cycle;
\fill[pthin] (-0.472,2.751) -- (-0.472,2.444) -- (-0.292,2.503) -- cycle;
\fill[pthick] (-0.472,2.444) -- (-0.472,2.751) -- (-0.763,2.349) -- cycle;
\fill[pthick] (-0.292,2.196) -- (0.000,2.101) -- (-0.292,2.503) -- cycle;
\fill[pthick] (-0.472,2.444) -- (-0.292,2.196) -- (-0.763,2.349) -- cycle;
\fill[pthin] (-0.292,2.196) -- (-0.472,2.444) -- (-0.292,2.503) -- cycle;
\fill[pthick] (-0.292,2.999) -- (-0.472,2.751) -- (0.000,2.904) -- cycle;
\fill[pthick] (0.000,3.093) -- (-0.292,2.999) -- (0.000,3.400) -- cycle;
\fill[pthin] (-0.292,2.999) -- (0.000,3.093) -- (0.000,2.904) -- cycle;
\fill[pthick] (0.000,2.408) -- (0.000,2.101) -- (-0.292,2.503) -- cycle;
\fill[pthick] (-0.180,2.656) -- (0.000,2.408) -- (0.000,2.904) -- cycle;
\fill[pthin] (0.000,2.408) -- (-0.180,2.656) -- (-0.292,2.503) -- cycle;
\fill[pthin] (-0.472,2.751) -- (-0.180,2.656) -- (-0.292,2.503) -- cycle;
\fill[pthick] (-0.180,2.656) -- (-0.472,2.751) -- (0.000,2.904) -- cycle;
\fill[pthick] (-1.055,2.751) -- (-1.235,2.999) -- (-1.235,2.503) -- cycle;
\fill[pthick] (-1.055,2.444) -- (-1.055,2.751) -- (-0.763,2.349) -- cycle;
\fill[pthin] (-1.055,2.751) -- (-1.055,2.444) -- (-1.235,2.503) -- cycle;
\fill[pthick] (-1.707,2.656) -- (-1.998,2.751) -- (-1.527,2.904) -- cycle;
\fill[pthick] (-1.415,2.751) -- (-1.707,2.656) -- (-1.235,2.503) -- cycle;
\fill[pthin] (-1.707,2.656) -- (-1.415,2.751) -- (-1.527,2.904) -- cycle;
\fill[pthin] (-1.235,2.999) -- (-1.415,2.751) -- (-1.527,2.904) -- cycle;
\fill[pthick] (-1.415,2.751) -- (-1.235,2.999) -- (-1.235,2.503) -- cycle;
\fill[pthick] (-0.292,2.999) -- (-0.472,2.751) -- (-0.472,3.247) -- cycle;
\fill[pthick] (-0.292,3.305) -- (-0.292,2.999) -- (0.000,3.400) -- cycle;
\fill[pthin] (-0.292,2.999) -- (-0.292,3.305) -- (-0.472,3.247) -- cycle;
\fill[pthick] (-0.944,3.093) -- (-1.235,2.999) -- (-0.763,2.845) -- cycle;
\fill[pthick] (-0.652,2.999) -- (-0.944,3.093) -- (-0.472,3.247) -- cycle;
\fill[pthin] (-0.944,3.093) -- (-0.652,2.999) -- (-0.763,2.845) -- cycle;
\fill[pthin] (-0.472,2.751) -- (-0.652,2.999) -- (-0.763,2.845) -- cycle;
\fill[pthick] (-0.652,2.999) -- (-0.472,2.751) -- (-0.472,3.247) -- cycle;
\fill[pthin] (-0.472,2.751) -- (-0.763,2.656) -- (-0.763,2.845) -- cycle;
\fill[pthick] (-0.763,2.656) -- (-0.472,2.751) -- (-0.763,2.349) -- cycle;
\fill[pthick] (-1.055,2.751) -- (-1.235,2.999) -- (-0.763,2.845) -- cycle;
\fill[pthick] (-0.763,2.656) -- (-1.055,2.751) -- (-0.763,2.349) -- cycle;
\fill[pthin] (-1.055,2.751) -- (-0.763,2.656) -- (-0.763,2.845) -- cycle;
\fill[pthick] (-0.292,1.204) -- (0.000,1.299) -- (-0.472,1.452) -- cycle;
\fill[pthick] (-0.583,1.299) -- (-0.292,1.204) -- (-0.763,1.051) -- cycle;
\fill[pthin] (-0.292,1.204) -- (-0.583,1.299) -- (-0.472,1.452) -- cycle;
\fill[pthick] (-0.180,1.853) -- (0.000,2.101) -- (0.000,1.605) -- cycle;
\fill[pthick] (-0.180,1.547) -- (-0.180,1.853) -- (-0.472,1.452) -- cycle;
\fill[pthin] (-0.180,1.853) -- (-0.180,1.547) -- (0.000,1.605) -- cycle;
\fill[pthin] (0.000,1.299) -- (-0.180,1.547) -- (0.000,1.605) -- cycle;
\fill[pthick] (-0.180,1.547) -- (0.000,1.299) -- (-0.472,1.452) -- cycle;
\fill[pthick] (-0.292,0.401) -- (-0.472,0.649) -- (0.000,0.496) -- cycle;
\fill[pthick] (0.000,0.307) -- (-0.292,0.401) -- (0.000,0.000) -- cycle;
\fill[pthin] (-0.292,0.401) -- (0.000,0.307) -- (0.000,0.496) -- cycle;
\fill[pthick] (0.000,0.992) -- (0.000,1.299) -- (-0.292,0.897) -- cycle;
\fill[pthick] (-0.180,0.744) -- (0.000,0.992) -- (0.000,0.496) -- cycle;
\fill[pthin] (0.000,0.992) -- (-0.180,0.744) -- (-0.292,0.897) -- cycle;
\fill[pthin] (-0.472,0.649) -- (-0.180,0.744) -- (-0.292,0.897) -- cycle;
\fill[pthick] (-0.180,0.744) -- (-0.472,0.649) -- (0.000,0.496) -- cycle;
\fill[pthin] (-0.472,0.649) -- (-0.472,0.956) -- (-0.292,0.897) -- cycle;
\fill[pthick] (-0.472,0.956) -- (-0.472,0.649) -- (-0.763,1.051) -- cycle;
\fill[pthick] (-0.292,1.204) -- (0.000,1.299) -- (-0.292,0.897) -- cycle;
\fill[pthick] (-0.472,0.956) -- (-0.292,1.204) -- (-0.763,1.051) -- cycle;
\fill[pthin] (-0.292,1.204) -- (-0.472,0.956) -- (-0.292,0.897) -- cycle;
\fill[pthick] (-1.235,2.196) -- (-1.527,2.101) -- (-1.235,2.503) -- cycle;
\fill[pthick] (-1.055,2.444) -- (-1.235,2.196) -- (-0.763,2.349) -- cycle;
\fill[pthin] (-1.235,2.196) -- (-1.055,2.444) -- (-1.235,2.503) -- cycle;
\fill[pthick] (-1.707,2.656) -- (-1.998,2.751) -- (-1.707,2.349) -- cycle;
\fill[pthick] (-1.527,2.408) -- (-1.707,2.656) -- (-1.235,2.503) -- cycle;
\fill[pthin] (-1.707,2.656) -- (-1.527,2.408) -- (-1.707,2.349) -- cycle;
\fill[pthin] (-1.527,2.101) -- (-1.527,2.408) -- (-1.707,2.349) -- cycle;
\fill[pthick] (-1.527,2.408) -- (-1.527,2.101) -- (-1.235,2.503) -- cycle;
\fill[pthick] (-0.763,1.547) -- (-0.763,1.853) -- (-1.055,1.452) -- cycle;
\fill[pthick] (-0.944,1.299) -- (-0.763,1.547) -- (-0.763,1.051) -- cycle;
\fill[pthin] (-0.763,1.547) -- (-0.944,1.299) -- (-1.055,1.452) -- cycle;
\fill[pthick] (-1.346,1.853) -- (-1.527,2.101) -- (-1.055,1.948) -- cycle;
\fill[pthick] (-1.055,1.759) -- (-1.346,1.853) -- (-1.055,1.452) -- cycle;
\fill[pthin] (-1.346,1.853) -- (-1.055,1.759) -- (-1.055,1.948) -- cycle;
\fill[pthin] (-0.763,1.853) -- (-1.055,1.759) -- (-1.055,1.948) -- cycle;
\fill[pthick] (-1.055,1.759) -- (-0.763,1.853) -- (-1.055,1.452) -- cycle;
\fill[pthin] (-0.763,1.853) -- (-0.944,2.101) -- (-1.055,1.948) -- cycle;
\fill[pthick] (-0.944,2.101) -- (-0.763,1.853) -- (-0.763,2.349) -- cycle;
\fill[pthick] (-1.235,2.196) -- (-1.527,2.101) -- (-1.055,1.948) -- cycle;
\fill[pthick] (-0.944,2.101) -- (-1.235,2.196) -- (-0.763,2.349) -- cycle;
\fill[pthin] (-1.235,2.196) -- (-0.944,2.101) -- (-1.055,1.948) -- cycle;
\fill[pthin] (-0.763,1.853) -- (-0.583,2.101) -- (-0.472,1.948) -- cycle;
\fill[pthick] (-0.583,2.101) -- (-0.763,1.853) -- (-0.763,2.349) -- cycle;
\fill[pthick] (-0.292,2.196) -- (0.000,2.101) -- (-0.472,1.948) -- cycle;
\fill[pthick] (-0.583,2.101) -- (-0.292,2.196) -- (-0.763,2.349) -- cycle;
\fill[pthin] (-0.292,2.196) -- (-0.583,2.101) -- (-0.472,1.948) -- cycle;
\fill[pthick] (-0.763,1.547) -- (-0.763,1.853) -- (-0.472,1.452) -- cycle;
\fill[pthick] (-0.583,1.299) -- (-0.763,1.547) -- (-0.763,1.051) -- cycle;
\fill[pthin] (-0.763,1.547) -- (-0.583,1.299) -- (-0.472,1.452) -- cycle;
\fill[pthick] (-0.180,1.853) -- (0.000,2.101) -- (-0.472,1.948) -- cycle;
\fill[pthick] (-0.472,1.759) -- (-0.180,1.853) -- (-0.472,1.452) -- cycle;
\fill[pthin] (-0.180,1.853) -- (-0.472,1.759) -- (-0.472,1.948) -- cycle;
\fill[pthin] (-0.763,1.853) -- (-0.472,1.759) -- (-0.472,1.948) -- cycle;
\fill[pthick] (-0.472,1.759) -- (-0.763,1.853) -- (-0.472,1.452) -- cycle;
\fill[pthin] (-2.470,1.299) -- (-2.179,1.204) -- (-2.290,1.051) -- cycle;
\fill[pthick] (-2.179,1.204) -- (-2.470,1.299) -- (-1.998,1.452) -- cycle;
\fill[pthick] (-1.998,0.956) -- (-1.998,0.649) -- (-2.290,1.051) -- cycle;
\fill[pthick] (-2.179,1.204) -- (-1.998,0.956) -- (-1.998,1.452) -- cycle;
\fill[pthin] (-1.998,0.956) -- (-2.179,1.204) -- (-2.290,1.051) -- cycle;
\fill[pthick] (-2.762,1.204) -- (-2.470,1.299) -- (-2.762,0.897) -- cycle;
\fill[pthick] (-2.942,0.956) -- (-2.762,1.204) -- (-3.234,1.051) -- cycle;
\fill[pthin] (-2.762,1.204) -- (-2.942,0.956) -- (-2.762,0.897) -- cycle;
\fill[pthick] (-2.290,0.744) -- (-1.998,0.649) -- (-2.290,1.051) -- cycle;
\fill[pthick] (-2.470,0.992) -- (-2.290,0.744) -- (-2.762,0.897) -- cycle;
\fill[pthin] (-2.290,0.744) -- (-2.470,0.992) -- (-2.290,1.051) -- cycle;
\fill[pthin] (-2.470,1.299) -- (-2.470,0.992) -- (-2.290,1.051) -- cycle;
\fill[pthick] (-2.470,0.992) -- (-2.470,1.299) -- (-2.762,0.897) -- cycle;
\fill[pthick] (-2.290,1.853) -- (-2.470,2.101) -- (-1.998,1.948) -- cycle;
\fill[pthick] (-1.998,1.759) -- (-2.290,1.853) -- (-1.998,1.452) -- cycle;
\fill[pthin] (-2.290,1.853) -- (-1.998,1.759) -- (-1.998,1.948) -- cycle;
\fill[pthick] (-1.998,2.444) -- (-1.998,2.751) -- (-2.290,2.349) -- cycle;
\fill[pthick] (-2.179,2.196) -- (-1.998,2.444) -- (-1.998,1.948) -- cycle;
\fill[pthin] (-1.998,2.444) -- (-2.179,2.196) -- (-2.290,2.349) -- cycle;
\fill[pthin] (-2.470,2.101) -- (-2.179,2.196) -- (-2.290,2.349) -- cycle;
\fill[pthick] (-2.179,2.196) -- (-2.470,2.101) -- (-1.998,1.948) -- cycle;
\fill[pthick] (-2.762,1.204) -- (-2.470,1.299) -- (-2.942,1.452) -- cycle;
\fill[pthick] (-3.053,1.299) -- (-2.762,1.204) -- (-3.234,1.051) -- cycle;
\fill[pthin] (-2.762,1.204) -- (-3.053,1.299) -- (-2.942,1.452) -- cycle;
\fill[pthick] (-2.650,1.853) -- (-2.470,2.101) -- (-2.470,1.605) -- cycle;
\fill[pthick] (-2.650,1.547) -- (-2.650,1.853) -- (-2.942,1.452) -- cycle;
\fill[pthin] (-2.650,1.853) -- (-2.650,1.547) -- (-2.470,1.605) -- cycle;
\fill[pthin] (-2.470,1.299) -- (-2.650,1.547) -- (-2.470,1.605) -- cycle;
\fill[pthick] (-2.650,1.547) -- (-2.470,1.299) -- (-2.942,1.452) -- cycle;
\fill[pthin] (-2.470,1.299) -- (-2.290,1.547) -- (-2.470,1.605) -- cycle;
\fill[pthick] (-2.290,1.547) -- (-2.470,1.299) -- (-1.998,1.452) -- cycle;
\fill[pthick] (-2.290,1.853) -- (-2.470,2.101) -- (-2.470,1.605) -- cycle;
\fill[pthick] (-2.290,1.547) -- (-2.290,1.853) -- (-1.998,1.452) -- cycle;
\fill[pthin] (-2.290,1.853) -- (-2.290,1.547) -- (-2.470,1.605) -- cycle;
\fill[pthick] (-1.055,0.649) -- (-1.235,0.401) -- (-1.235,0.897) -- cycle;
\fill[pthick] (-1.055,0.956) -- (-1.055,0.649) -- (-0.763,1.051) -- cycle;
\fill[pthin] (-1.055,0.649) -- (-1.055,0.956) -- (-1.235,0.897) -- cycle;
\fill[pthick] (-1.707,0.744) -- (-1.998,0.649) -- (-1.527,0.496) -- cycle;
\fill[pthick] (-1.415,0.649) -- (-1.707,0.744) -- (-1.235,0.897) -- cycle;
\fill[pthin] (-1.707,0.744) -- (-1.415,0.649) -- (-1.527,0.496) -- cycle;
\fill[pthin] (-1.235,0.401) -- (-1.415,0.649) -- (-1.527,0.496) -- cycle;
\fill[pthick] (-1.415,0.649) -- (-1.235,0.401) -- (-1.235,0.897) -- cycle;
\fill[pthick] (-0.292,0.401) -- (-0.472,0.649) -- (-0.472,0.153) -- cycle;
\fill[pthick] (-0.292,0.095) -- (-0.292,0.401) -- (0.000,0.000) -- cycle;
\fill[pthin] (-0.292,0.401) -- (-0.292,0.095) -- (-0.472,0.153) -- cycle;
\fill[pthick] (-0.944,0.307) -- (-1.235,0.401) -- (-0.763,0.555) -- cycle;
\fill[pthick] (-0.652,0.401) -- (-0.944,0.307) -- (-0.472,0.153) -- cycle;
\fill[pthin] (-0.944,0.307) -- (-0.652,0.401) -- (-0.763,0.555) -- cycle;
\fill[pthin] (-0.472,0.649) -- (-0.652,0.401) -- (-0.763,0.555) -- cycle;
\fill[pthick] (-0.652,0.401) -- (-0.472,0.649) -- (-0.472,0.153) -- cycle;
\fill[pthin] (-0.472,0.649) -- (-0.763,0.744) -- (-0.763,0.555) -- cycle;
\fill[pthick] (-0.763,0.744) -- (-0.472,0.649) -- (-0.763,1.051) -- cycle;
\fill[pthick] (-1.055,0.649) -- (-1.235,0.401) -- (-0.763,0.555) -- cycle;
\fill[pthick] (-0.763,0.744) -- (-1.055,0.649) -- (-0.763,1.051) -- cycle;
\fill[pthin] (-1.055,0.649) -- (-0.763,0.744) -- (-0.763,0.555) -- cycle;
\fill[pthick] (-1.707,1.853) -- (-1.527,2.101) -- (-1.998,1.948) -- cycle;
\fill[pthick] (-1.998,1.759) -- (-1.707,1.853) -- (-1.998,1.452) -- cycle;
\fill[pthin] (-1.707,1.853) -- (-1.998,1.759) -- (-1.998,1.948) -- cycle;
\fill[pthick] (-1.998,2.444) -- (-1.998,2.751) -- (-1.707,2.349) -- cycle;
\fill[pthick] (-1.818,2.196) -- (-1.998,2.444) -- (-1.998,1.948) -- cycle;
\fill[pthin] (-1.998,2.444) -- (-1.818,2.196) -- (-1.707,2.349) -- cycle;
\fill[pthin] (-1.527,2.101) -- (-1.818,2.196) -- (-1.707,2.349) -- cycle;
\fill[pthick] (-1.818,2.196) -- (-1.527,2.101) -- (-1.998,1.948) -- cycle;
\fill[pthick] (-1.235,1.204) -- (-1.527,1.299) -- (-1.055,1.452) -- cycle;
\fill[pthick] (-0.944,1.299) -- (-1.235,1.204) -- (-0.763,1.051) -- cycle;
\fill[pthin] (-1.235,1.204) -- (-0.944,1.299) -- (-1.055,1.452) -- cycle;
\fill[pthick] (-1.346,1.853) -- (-1.527,2.101) -- (-1.527,1.605) -- cycle;
\fill[pthick] (-1.346,1.547) -- (-1.346,1.853) -- (-1.055,1.452) -- cycle;
\fill[pthin] (-1.346,1.853) -- (-1.346,1.547) -- (-1.527,1.605) -- cycle;
\fill[pthin] (-1.527,1.299) -- (-1.346,1.547) -- (-1.527,1.605) -- cycle;
\fill[pthick] (-1.346,1.547) -- (-1.527,1.299) -- (-1.055,1.452) -- cycle;
\fill[pthin] (-1.527,1.299) -- (-1.707,1.547) -- (-1.527,1.605) -- cycle;
\fill[pthick] (-1.707,1.547) -- (-1.527,1.299) -- (-1.998,1.452) -- cycle;
\fill[pthick] (-1.707,1.853) -- (-1.527,2.101) -- (-1.527,1.605) -- cycle;
\fill[pthick] (-1.707,1.547) -- (-1.707,1.853) -- (-1.998,1.452) -- cycle;
\fill[pthin] (-1.707,1.853) -- (-1.707,1.547) -- (-1.527,1.605) -- cycle;
\fill[pthin] (-1.527,1.299) -- (-1.818,1.204) -- (-1.707,1.051) -- cycle;
\fill[pthick] (-1.818,1.204) -- (-1.527,1.299) -- (-1.998,1.452) -- cycle;
\fill[pthick] (-1.998,0.956) -- (-1.998,0.649) -- (-1.707,1.051) -- cycle;
\fill[pthick] (-1.818,1.204) -- (-1.998,0.956) -- (-1.998,1.452) -- cycle;
\fill[pthin] (-1.998,0.956) -- (-1.818,1.204) -- (-1.707,1.051) -- cycle;
\fill[pthick] (-1.235,1.204) -- (-1.527,1.299) -- (-1.235,0.897) -- cycle;
\fill[pthick] (-1.055,0.956) -- (-1.235,1.204) -- (-0.763,1.051) -- cycle;
\fill[pthin] (-1.235,1.204) -- (-1.055,0.956) -- (-1.235,0.897) -- cycle;
\fill[pthick] (-1.707,0.744) -- (-1.998,0.649) -- (-1.707,1.051) -- cycle;
\fill[pthick] (-1.527,0.992) -- (-1.707,0.744) -- (-1.235,0.897) -- cycle;
\fill[pthin] (-1.707,0.744) -- (-1.527,0.992) -- (-1.707,1.051) -- cycle;
\fill[pthin] (-1.527,1.299) -- (-1.527,0.992) -- (-1.707,1.051) -- cycle;
\fill[pthick] (-1.527,0.992) -- (-1.527,1.299) -- (-1.235,0.897) -- cycle;
\fill[pthin] (-2.762,0.401) -- (-2.470,0.307) -- (-2.470,0.496) -- cycle;
\fill[pthick] (-2.470,0.307) -- (-2.762,0.401) -- (-2.470,0.000) -- cycle;
\fill[pthick] (-2.179,0.401) -- (-1.998,0.649) -- (-2.470,0.496) -- cycle;
\fill[pthick] (-2.470,0.307) -- (-2.179,0.401) -- (-2.470,0.000) -- cycle;
\fill[pthin] (-2.179,0.401) -- (-2.470,0.307) -- (-2.470,0.496) -- cycle;
\fill[pthick] (-2.942,0.649) -- (-2.762,0.401) -- (-2.762,0.897) -- cycle;
\fill[pthick] (-2.942,0.956) -- (-2.942,0.649) -- (-3.234,1.051) -- cycle;
\fill[pthin] (-2.942,0.649) -- (-2.942,0.956) -- (-2.762,0.897) -- cycle;
\fill[pthick] (-2.290,0.744) -- (-1.998,0.649) -- (-2.470,0.496) -- cycle;
\fill[pthick] (-2.582,0.649) -- (-2.290,0.744) -- (-2.762,0.897) -- cycle;
\fill[pthin] (-2.290,0.744) -- (-2.582,0.649) -- (-2.470,0.496) -- cycle;
\fill[pthin] (-2.762,0.401) -- (-2.582,0.649) -- (-2.470,0.496) -- cycle;
\fill[pthick] (-2.582,0.649) -- (-2.762,0.401) -- (-2.762,0.897) -- cycle;
\fill[pthick] (-2.942,-0.153) -- (-3.234,-0.248) -- (-2.762,-0.401) -- cycle;
\fill[pthick] (-2.650,-0.248) -- (-2.942,-0.153) -- (-2.470,0.000) -- cycle;
\fill[pthin] (-2.942,-0.153) -- (-2.650,-0.248) -- (-2.762,-0.401) -- cycle;
\fill[pthick] (-3.053,-0.803) -- (-3.234,-1.051) -- (-3.234,-0.555) -- cycle;
\fill[pthick] (-3.053,-0.496) -- (-3.053,-0.803) -- (-2.762,-0.401) -- cycle;
\fill[pthin] (-3.053,-0.803) -- (-3.053,-0.496) -- (-3.234,-0.555) -- cycle;
\fill[pthin] (-3.234,-0.248) -- (-3.053,-0.496) -- (-3.234,-0.555) -- cycle;
\fill[pthick] (-3.053,-0.496) -- (-3.234,-0.248) -- (-2.762,-0.401) -- cycle;
\fill[pthick] (-2.942,0.649) -- (-2.762,0.401) -- (-3.234,0.555) -- cycle;
\fill[pthick] (-3.234,0.744) -- (-2.942,0.649) -- (-3.234,1.051) -- cycle;
\fill[pthin] (-2.942,0.649) -- (-3.234,0.744) -- (-3.234,0.555) -- cycle;
\fill[pthick] (-3.234,0.059) -- (-3.234,-0.248) -- (-2.942,0.153) -- cycle;
\fill[pthick] (-3.053,0.307) -- (-3.234,0.059) -- (-3.234,0.555) -- cycle;
\fill[pthin] (-3.234,0.059) -- (-3.053,0.307) -- (-2.942,0.153) -- cycle;
\fill[pthin] (-2.762,0.401) -- (-3.053,0.307) -- (-2.942,0.153) -- cycle;
\fill[pthick] (-3.053,0.307) -- (-2.762,0.401) -- (-3.234,0.555) -- cycle;
\fill[pthin] (-2.762,0.401) -- (-2.762,0.095) -- (-2.942,0.153) -- cycle;
\fill[pthick] (-2.762,0.095) -- (-2.762,0.401) -- (-2.470,0.000) -- cycle;
\fill[pthick] (-2.942,-0.153) -- (-3.234,-0.248) -- (-2.942,0.153) -- cycle;
\fill[pthick] (-2.762,0.095) -- (-2.942,-0.153) -- (-2.470,0.000) -- cycle;
\fill[pthin] (-2.942,-0.153) -- (-2.762,0.095) -- (-2.942,0.153) -- cycle;
\fill[pthick] (-1.235,0.095) -- (-1.235,0.401) -- (-1.527,0.000) -- cycle;
\fill[pthick] (-1.415,-0.153) -- (-1.235,0.095) -- (-1.235,-0.401) -- cycle;
\fill[pthin] (-1.235,0.095) -- (-1.415,-0.153) -- (-1.527,0.000) -- cycle;
\fill[pthick] (-1.818,0.401) -- (-1.998,0.649) -- (-1.527,0.496) -- cycle;
\fill[pthick] (-1.527,0.307) -- (-1.818,0.401) -- (-1.527,0.000) -- cycle;
\fill[pthin] (-1.818,0.401) -- (-1.527,0.307) -- (-1.527,0.496) -- cycle;
\fill[pthin] (-1.235,0.401) -- (-1.527,0.307) -- (-1.527,0.496) -- cycle;
\fill[pthick] (-1.527,0.307) -- (-1.235,0.401) -- (-1.527,0.000) -- cycle;
\fill[pthick] (-0.472,-0.153) -- (-0.763,-0.248) -- (-0.472,0.153) -- cycle;
\fill[pthick] (-0.292,0.095) -- (-0.472,-0.153) -- (0.000,0.000) -- cycle;
\fill[pthin] (-0.472,-0.153) -- (-0.292,0.095) -- (-0.472,0.153) -- cycle;
\fill[pthick] (-0.944,0.307) -- (-1.235,0.401) -- (-0.944,0.000) -- cycle;
\fill[pthick] (-0.763,0.059) -- (-0.944,0.307) -- (-0.472,0.153) -- cycle;
\fill[pthin] (-0.944,0.307) -- (-0.763,0.059) -- (-0.944,0.000) -- cycle;
\fill[pthin] (-0.763,-0.248) -- (-0.763,0.059) -- (-0.944,0.000) -- cycle;
\fill[pthick] (-0.763,0.059) -- (-0.763,-0.248) -- (-0.472,0.153) -- cycle;
\fill[pthin] (-0.763,-0.248) -- (-1.055,-0.153) -- (-0.944,0.000) -- cycle;
\fill[pthick] (-1.055,-0.153) -- (-0.763,-0.248) -- (-1.235,-0.401) -- cycle;
\fill[pthick] (-1.235,0.095) -- (-1.235,0.401) -- (-0.944,0.000) -- cycle;
\fill[pthick] (-1.055,-0.153) -- (-1.235,0.095) -- (-1.235,-0.401) -- cycle;
\fill[pthin] (-1.235,0.095) -- (-1.055,-0.153) -- (-0.944,0.000) -- cycle;
\fill[pthick] (-2.470,-0.496) -- (-2.470,-0.803) -- (-2.762,-0.401) -- cycle;
\fill[pthick] (-2.650,-0.248) -- (-2.470,-0.496) -- (-2.470,0.000) -- cycle;
\fill[pthin] (-2.470,-0.496) -- (-2.650,-0.248) -- (-2.762,-0.401) -- cycle;
\fill[pthick] (-3.053,-0.803) -- (-3.234,-1.051) -- (-2.762,-0.897) -- cycle;
\fill[pthick] (-2.762,-0.708) -- (-3.053,-0.803) -- (-2.762,-0.401) -- cycle;
\fill[pthin] (-3.053,-0.803) -- (-2.762,-0.708) -- (-2.762,-0.897) -- cycle;
\fill[pthin] (-2.470,-0.803) -- (-2.762,-0.708) -- (-2.762,-0.897) -- cycle;
\fill[pthick] (-2.762,-0.708) -- (-2.470,-0.803) -- (-2.762,-0.401) -- cycle;
\fill[pthick] (-1.707,-0.248) -- (-1.998,-0.153) -- (-1.707,-0.555) -- cycle;
\fill[pthick] (-1.527,-0.496) -- (-1.707,-0.248) -- (-1.235,-0.401) -- cycle;
\fill[pthin] (-1.707,-0.248) -- (-1.527,-0.496) -- (-1.707,-0.555) -- cycle;
\fill[pthick] (-2.179,-0.708) -- (-2.470,-0.803) -- (-2.179,-0.401) -- cycle;
\fill[pthick] (-1.998,-0.460) -- (-2.179,-0.708) -- (-1.707,-0.555) -- cycle;
\fill[pthin] (-2.179,-0.708) -- (-1.998,-0.460) -- (-2.179,-0.401) -- cycle;
\fill[pthin] (-1.998,-0.153) -- (-1.998,-0.460) -- (-2.179,-0.401) -- cycle;
\fill[pthick] (-1.998,-0.460) -- (-1.998,-0.153) -- (-1.707,-0.555) -- cycle;
\fill[pthin] (-1.998,-0.153) -- (-2.290,-0.248) -- (-2.179,-0.401) -- cycle;
\fill[pthick] (-2.290,-0.248) -- (-1.998,-0.153) -- (-2.470,0.000) -- cycle;
\fill[pthick] (-2.470,-0.496) -- (-2.470,-0.803) -- (-2.179,-0.401) -- cycle;
\fill[pthick] (-2.290,-0.248) -- (-2.470,-0.496) -- (-2.470,0.000) -- cycle;
\fill[pthin] (-2.470,-0.496) -- (-2.290,-0.248) -- (-2.179,-0.401) -- cycle;
\fill[pthin] (-1.998,-0.153) -- (-2.179,0.095) -- (-1.998,0.153) -- cycle;
\fill[pthick] (-2.179,0.095) -- (-1.998,-0.153) -- (-2.470,0.000) -- cycle;
\fill[pthick] (-2.179,0.401) -- (-1.998,0.649) -- (-1.998,0.153) -- cycle;
\fill[pthick] (-2.179,0.095) -- (-2.179,0.401) -- (-2.470,0.000) -- cycle;
\fill[pthin] (-2.179,0.401) -- (-2.179,0.095) -- (-1.998,0.153) -- cycle;
\fill[pthick] (-1.707,-0.248) -- (-1.998,-0.153) -- (-1.527,0.000) -- cycle;
\fill[pthick] (-1.415,-0.153) -- (-1.707,-0.248) -- (-1.235,-0.401) -- cycle;
\fill[pthin] (-1.707,-0.248) -- (-1.415,-0.153) -- (-1.527,0.000) -- cycle;
\fill[pthick] (-1.818,0.401) -- (-1.998,0.649) -- (-1.998,0.153) -- cycle;
\fill[pthick] (-1.818,0.095) -- (-1.818,0.401) -- (-1.527,0.000) -- cycle;
\fill[pthin] (-1.818,0.401) -- (-1.818,0.095) -- (-1.998,0.153) -- cycle;
\fill[pthin] (-1.998,-0.153) -- (-1.818,0.095) -- (-1.998,0.153) -- cycle;
\fill[pthick] (-1.818,0.095) -- (-1.998,-0.153) -- (-1.527,0.000) -- cycle;
\fill[pthin] (-1.998,-1.948) -- (-1.818,-1.700) -- (-1.707,-1.853) -- cycle;
\fill[pthick] (-1.818,-1.700) -- (-1.998,-1.948) -- (-1.998,-1.452) -- cycle;
\fill[pthick] (-1.527,-1.605) -- (-1.235,-1.700) -- (-1.707,-1.853) -- cycle;
\fill[pthick] (-1.818,-1.700) -- (-1.527,-1.605) -- (-1.998,-1.452) -- cycle;
\fill[pthin] (-1.527,-1.605) -- (-1.818,-1.700) -- (-1.707,-1.853) -- cycle;
\fill[pthick] (-1.998,-2.255) -- (-1.998,-1.948) -- (-1.707,-2.349) -- cycle;
\fill[pthick] (-1.818,-2.503) -- (-1.998,-2.255) -- (-1.998,-2.751) -- cycle;
\fill[pthin] (-1.998,-2.255) -- (-1.818,-2.503) -- (-1.707,-2.349) -- cycle;
\fill[pthick] (-1.415,-1.948) -- (-1.235,-1.700) -- (-1.707,-1.853) -- cycle;
\fill[pthick] (-1.707,-2.043) -- (-1.415,-1.948) -- (-1.707,-2.349) -- cycle;
\fill[pthin] (-1.415,-1.948) -- (-1.707,-2.043) -- (-1.707,-1.853) -- cycle;
\fill[pthin] (-1.998,-1.948) -- (-1.707,-2.043) -- (-1.707,-1.853) -- cycle;
\fill[pthick] (-1.707,-2.043) -- (-1.998,-1.948) -- (-1.707,-2.349) -- cycle;
\fill[pthick] (-2.470,-1.605) -- (-2.762,-1.700) -- (-2.470,-1.299) -- cycle;
\fill[pthick] (-2.290,-1.357) -- (-2.470,-1.605) -- (-1.998,-1.452) -- cycle;
\fill[pthin] (-2.470,-1.605) -- (-2.290,-1.357) -- (-2.470,-1.299) -- cycle;
\fill[pthick] (-2.942,-1.145) -- (-3.234,-1.051) -- (-2.942,-1.452) -- cycle;
\fill[pthick] (-2.762,-1.393) -- (-2.942,-1.145) -- (-2.470,-1.299) -- cycle;
\fill[pthin] (-2.942,-1.145) -- (-2.762,-1.393) -- (-2.942,-1.452) -- cycle;
\fill[pthin] (-2.762,-1.700) -- (-2.762,-1.393) -- (-2.942,-1.452) -- cycle;
\fill[pthick] (-2.762,-1.393) -- (-2.762,-1.700) -- (-2.470,-1.299) -- cycle;
\fill[pthick] (-1.998,-2.255) -- (-1.998,-1.948) -- (-2.290,-2.349) -- cycle;
\fill[pthick] (-2.179,-2.503) -- (-1.998,-2.255) -- (-1.998,-2.751) -- cycle;
\fill[pthin] (-1.998,-2.255) -- (-2.179,-2.503) -- (-2.290,-2.349) -- cycle;
\fill[pthick] (-2.582,-1.948) -- (-2.762,-1.700) -- (-2.290,-1.853) -- cycle;
\fill[pthick] (-2.290,-2.043) -- (-2.582,-1.948) -- (-2.290,-2.349) -- cycle;
\fill[pthin] (-2.582,-1.948) -- (-2.290,-2.043) -- (-2.290,-1.853) -- cycle;
\fill[pthin] (-1.998,-1.948) -- (-2.290,-2.043) -- (-2.290,-1.853) -- cycle;
\fill[pthick] (-2.290,-2.043) -- (-1.998,-1.948) -- (-2.290,-2.349) -- cycle;
\fill[pthin] (-1.998,-1.948) -- (-2.179,-1.700) -- (-2.290,-1.853) -- cycle;
\fill[pthick] (-2.179,-1.700) -- (-1.998,-1.948) -- (-1.998,-1.452) -- cycle;
\fill[pthick] (-2.470,-1.605) -- (-2.762,-1.700) -- (-2.290,-1.853) -- cycle;
\fill[pthick] (-2.179,-1.700) -- (-2.470,-1.605) -- (-1.998,-1.452) -- cycle;
\fill[pthin] (-2.470,-1.605) -- (-2.179,-1.700) -- (-2.290,-1.853) -- cycle;
\fill[pthick] (-0.944,-0.803) -- (-0.763,-1.051) -- (-1.235,-0.897) -- cycle;
\fill[pthick] (-1.235,-0.708) -- (-0.944,-0.803) -- (-1.235,-0.401) -- cycle;
\fill[pthin] (-0.944,-0.803) -- (-1.235,-0.708) -- (-1.235,-0.897) -- cycle;
\fill[pthick] (-1.235,-1.393) -- (-1.235,-1.700) -- (-0.944,-1.299) -- cycle;
\fill[pthick] (-1.055,-1.145) -- (-1.235,-1.393) -- (-1.235,-0.897) -- cycle;
\fill[pthin] (-1.235,-1.393) -- (-1.055,-1.145) -- (-0.944,-1.299) -- cycle;
\fill[pthin] (-0.763,-1.051) -- (-1.055,-1.145) -- (-0.944,-1.299) -- cycle;
\fill[pthick] (-1.055,-1.145) -- (-0.763,-1.051) -- (-1.235,-0.897) -- cycle;
\fill[pthick] (-0.472,-0.153) -- (-0.763,-0.248) -- (-0.292,-0.401) -- cycle;
\fill[pthick] (-0.180,-0.248) -- (-0.472,-0.153) -- (0.000,0.000) -- cycle;
\fill[pthin] (-0.472,-0.153) -- (-0.180,-0.248) -- (-0.292,-0.401) -- cycle;
\fill[pthick] (-0.583,-0.803) -- (-0.763,-1.051) -- (-0.763,-0.555) -- cycle;
\fill[pthick] (-0.583,-0.496) -- (-0.583,-0.803) -- (-0.292,-0.401) -- cycle;
\fill[pthin] (-0.583,-0.803) -- (-0.583,-0.496) -- (-0.763,-0.555) -- cycle;
\fill[pthin] (-0.763,-0.248) -- (-0.583,-0.496) -- (-0.763,-0.555) -- cycle;
\fill[pthick] (-0.583,-0.496) -- (-0.763,-0.248) -- (-0.292,-0.401) -- cycle;
\fill[pthin] (-0.763,-0.248) -- (-0.944,-0.496) -- (-0.763,-0.555) -- cycle;
\fill[pthick] (-0.944,-0.496) -- (-0.763,-0.248) -- (-1.235,-0.401) -- cycle;
\fill[pthick] (-0.944,-0.803) -- (-0.763,-1.051) -- (-0.763,-0.555) -- cycle;
\fill[pthick] (-0.944,-0.496) -- (-0.944,-0.803) -- (-1.235,-0.401) -- cycle;
\fill[pthin] (-0.944,-0.803) -- (-0.944,-0.496) -- (-0.763,-0.555) -- cycle;
\fill[pthick] (-2.290,-1.051) -- (-2.470,-0.803) -- (-2.470,-1.299) -- cycle;
\fill[pthick] (-2.290,-1.357) -- (-2.290,-1.051) -- (-1.998,-1.452) -- cycle;
\fill[pthin] (-2.290,-1.051) -- (-2.290,-1.357) -- (-2.470,-1.299) -- cycle;
\fill[pthick] (-2.942,-1.145) -- (-3.234,-1.051) -- (-2.762,-0.897) -- cycle;
\fill[pthick] (-2.650,-1.051) -- (-2.942,-1.145) -- (-2.470,-1.299) -- cycle;
\fill[pthin] (-2.942,-1.145) -- (-2.650,-1.051) -- (-2.762,-0.897) -- cycle;
\fill[pthin] (-2.470,-0.803) -- (-2.650,-1.051) -- (-2.762,-0.897) -- cycle;
\fill[pthick] (-2.650,-1.051) -- (-2.470,-0.803) -- (-2.470,-1.299) -- cycle;
\fill[pthick] (-1.527,-0.803) -- (-1.707,-1.051) -- (-1.707,-0.555) -- cycle;
\fill[pthick] (-1.527,-0.496) -- (-1.527,-0.803) -- (-1.235,-0.401) -- cycle;
\fill[pthin] (-1.527,-0.803) -- (-1.527,-0.496) -- (-1.707,-0.555) -- cycle;
\fill[pthick] (-2.179,-0.708) -- (-2.470,-0.803) -- (-1.998,-0.956) -- cycle;
\fill[pthick] (-1.887,-0.803) -- (-2.179,-0.708) -- (-1.707,-0.555) -- cycle;
\fill[pthin] (-2.179,-0.708) -- (-1.887,-0.803) -- (-1.998,-0.956) -- cycle;
\fill[pthin] (-1.707,-1.051) -- (-1.887,-0.803) -- (-1.998,-0.956) -- cycle;
\fill[pthick] (-1.887,-0.803) -- (-1.707,-1.051) -- (-1.707,-0.555) -- cycle;
\fill[pthin] (-1.707,-1.051) -- (-1.998,-1.145) -- (-1.998,-0.956) -- cycle;
\fill[pthick] (-1.998,-1.145) -- (-1.707,-1.051) -- (-1.998,-1.452) -- cycle;
\fill[pthick] (-2.290,-1.051) -- (-2.470,-0.803) -- (-1.998,-0.956) -- cycle;
\fill[pthick] (-1.998,-1.145) -- (-2.290,-1.051) -- (-1.998,-1.452) -- cycle;
\fill[pthin] (-2.290,-1.051) -- (-1.998,-1.145) -- (-1.998,-0.956) -- cycle;
\fill[pthin] (-1.707,-1.051) -- (-1.707,-1.357) -- (-1.527,-1.299) -- cycle;
\fill[pthick] (-1.707,-1.357) -- (-1.707,-1.051) -- (-1.998,-1.452) -- cycle;
\fill[pthick] (-1.527,-1.605) -- (-1.235,-1.700) -- (-1.527,-1.299) -- cycle;
\fill[pthick] (-1.707,-1.357) -- (-1.527,-1.605) -- (-1.998,-1.452) -- cycle;
\fill[pthin] (-1.527,-1.605) -- (-1.707,-1.357) -- (-1.527,-1.299) -- cycle;
\fill[pthick] (-1.527,-0.803) -- (-1.707,-1.051) -- (-1.235,-0.897) -- cycle;
\fill[pthick] (-1.235,-0.708) -- (-1.527,-0.803) -- (-1.235,-0.401) -- cycle;
\fill[pthin] (-1.527,-0.803) -- (-1.235,-0.708) -- (-1.235,-0.897) -- cycle;
\fill[pthick] (-1.235,-1.393) -- (-1.235,-1.700) -- (-1.527,-1.299) -- cycle;
\fill[pthick] (-1.415,-1.145) -- (-1.235,-1.393) -- (-1.235,-0.897) -- cycle;
\fill[pthin] (-1.235,-1.393) -- (-1.415,-1.145) -- (-1.527,-1.299) -- cycle;
\fill[pthin] (-1.707,-1.051) -- (-1.415,-1.145) -- (-1.527,-1.299) -- cycle;
\fill[pthick] (-1.415,-1.145) -- (-1.707,-1.051) -- (-1.235,-0.897) -- cycle;
\fill[pthin] (-1.235,-2.503) -- (-1.055,-2.255) -- (-1.235,-2.196) -- cycle;
\fill[pthick] (-1.055,-2.255) -- (-1.235,-2.503) -- (-0.763,-2.349) -- cycle;
\fill[pthick] (-1.055,-1.948) -- (-1.235,-1.700) -- (-1.235,-2.196) -- cycle;
\fill[pthick] (-1.055,-2.255) -- (-1.055,-1.948) -- (-0.763,-2.349) -- cycle;
\fill[pthin] (-1.055,-1.948) -- (-1.055,-2.255) -- (-1.235,-2.196) -- cycle;
\fill[pthick] (-1.527,-2.597) -- (-1.235,-2.503) -- (-1.707,-2.349) -- cycle;
\fill[pthick] (-1.818,-2.503) -- (-1.527,-2.597) -- (-1.998,-2.751) -- cycle;
\fill[pthin] (-1.527,-2.597) -- (-1.818,-2.503) -- (-1.707,-2.349) -- cycle;
\fill[pthick] (-1.415,-1.948) -- (-1.235,-1.700) -- (-1.235,-2.196) -- cycle;
\fill[pthick] (-1.415,-2.255) -- (-1.415,-1.948) -- (-1.707,-2.349) -- cycle;
\fill[pthin] (-1.415,-1.948) -- (-1.415,-2.255) -- (-1.235,-2.196) -- cycle;
\fill[pthin] (-1.235,-2.503) -- (-1.415,-2.255) -- (-1.235,-2.196) -- cycle;
\fill[pthick] (-1.415,-2.255) -- (-1.235,-2.503) -- (-1.707,-2.349) -- cycle;
\fill[pthick] (-0.763,-2.845) -- (-0.763,-3.152) -- (-0.472,-2.751) -- cycle;
\fill[pthick] (-0.583,-2.597) -- (-0.763,-2.845) -- (-0.763,-2.349) -- cycle;
\fill[pthin] (-0.763,-2.845) -- (-0.583,-2.597) -- (-0.472,-2.751) -- cycle;
\fill[pthick] (-0.180,-3.152) -- (-0.000,-3.400) -- (-0.472,-3.247) -- cycle;
\fill[pthick] (-0.472,-3.057) -- (-0.180,-3.152) -- (-0.472,-2.751) -- cycle;
\fill[pthin] (-0.180,-3.152) -- (-0.472,-3.057) -- (-0.472,-3.247) -- cycle;
\fill[pthin] (-0.763,-3.152) -- (-0.472,-3.057) -- (-0.472,-3.247) -- cycle;
\fill[pthick] (-0.472,-3.057) -- (-0.763,-3.152) -- (-0.472,-2.751) -- cycle;
\fill[pthick] (-1.527,-2.597) -- (-1.235,-2.503) -- (-1.527,-2.904) -- cycle;
\fill[pthick] (-1.707,-2.845) -- (-1.527,-2.597) -- (-1.998,-2.751) -- cycle;
\fill[pthin] (-1.527,-2.597) -- (-1.707,-2.845) -- (-1.527,-2.904) -- cycle;
\fill[pthick] (-1.055,-3.057) -- (-0.763,-3.152) -- (-1.055,-2.751) -- cycle;
\fill[pthick] (-1.235,-2.809) -- (-1.055,-3.057) -- (-1.527,-2.904) -- cycle;
\fill[pthin] (-1.055,-3.057) -- (-1.235,-2.809) -- (-1.055,-2.751) -- cycle;
\fill[pthin] (-1.235,-2.503) -- (-1.235,-2.809) -- (-1.055,-2.751) -- cycle;
\fill[pthick] (-1.235,-2.809) -- (-1.235,-2.503) -- (-1.527,-2.904) -- cycle;
\fill[pthin] (-1.235,-2.503) -- (-0.944,-2.597) -- (-1.055,-2.751) -- cycle;
\fill[pthick] (-0.944,-2.597) -- (-1.235,-2.503) -- (-0.763,-2.349) -- cycle;
\fill[pthick] (-0.763,-2.845) -- (-0.763,-3.152) -- (-1.055,-2.751) -- cycle;
\fill[pthick] (-0.944,-2.597) -- (-0.763,-2.845) -- (-0.763,-2.349) -- cycle;
\fill[pthin] (-0.763,-2.845) -- (-0.944,-2.597) -- (-1.055,-2.751) -- cycle;
\fill[pthick] (-0.472,-1.145) -- (-0.763,-1.051) -- (-0.472,-1.452) -- cycle;
\fill[pthick] (-0.292,-1.393) -- (-0.472,-1.145) -- (-0.000,-1.299) -- cycle;
\fill[pthin] (-0.472,-1.145) -- (-0.292,-1.393) -- (-0.472,-1.452) -- cycle;
\fill[pthick] (-0.944,-1.605) -- (-1.235,-1.700) -- (-0.944,-1.299) -- cycle;
\fill[pthick] (-0.763,-1.357) -- (-0.944,-1.605) -- (-0.472,-1.452) -- cycle;
\fill[pthin] (-0.944,-1.605) -- (-0.763,-1.357) -- (-0.944,-1.299) -- cycle;
\fill[pthin] (-0.763,-1.051) -- (-0.763,-1.357) -- (-0.944,-1.299) -- cycle;
\fill[pthick] (-0.763,-1.357) -- (-0.763,-1.051) -- (-0.472,-1.452) -- cycle;
\fill[pthick] (-0.000,-0.496) -- (-0.000,-0.803) -- (-0.292,-0.401) -- cycle;
\fill[pthick] (-0.180,-0.248) -- (-0.000,-0.496) -- (0.000,0.000) -- cycle;
\fill[pthin] (-0.000,-0.496) -- (-0.180,-0.248) -- (-0.292,-0.401) -- cycle;
\fill[pthick] (-0.583,-0.803) -- (-0.763,-1.051) -- (-0.292,-0.897) -- cycle;
\fill[pthick] (-0.292,-0.708) -- (-0.583,-0.803) -- (-0.292,-0.401) -- cycle;
\fill[pthin] (-0.583,-0.803) -- (-0.292,-0.708) -- (-0.292,-0.897) -- cycle;
\fill[pthin] (-0.000,-0.803) -- (-0.292,-0.708) -- (-0.292,-0.897) -- cycle;
\fill[pthick] (-0.292,-0.708) -- (-0.000,-0.803) -- (-0.292,-0.401) -- cycle;
\fill[pthin] (-0.000,-0.803) -- (-0.180,-1.051) -- (-0.292,-0.897) -- cycle;
\fill[pthick] (-0.180,-1.051) -- (-0.000,-0.803) -- (-0.000,-1.299) -- cycle;
\fill[pthick] (-0.472,-1.145) -- (-0.763,-1.051) -- (-0.292,-0.897) -- cycle;
\fill[pthick] (-0.180,-1.051) -- (-0.472,-1.145) -- (-0.000,-1.299) -- cycle;
\fill[pthin] (-0.472,-1.145) -- (-0.180,-1.051) -- (-0.292,-0.897) -- cycle;
\fill[pthick] (-0.292,-2.503) -- (-0.000,-2.597) -- (-0.472,-2.751) -- cycle;
\fill[pthick] (-0.583,-2.597) -- (-0.292,-2.503) -- (-0.763,-2.349) -- cycle;
\fill[pthin] (-0.292,-2.503) -- (-0.583,-2.597) -- (-0.472,-2.751) -- cycle;
\fill[pthick] (-0.180,-3.152) -- (-0.000,-3.400) -- (-0.000,-2.904) -- cycle;
\fill[pthick] (-0.180,-2.845) -- (-0.180,-3.152) -- (-0.472,-2.751) -- cycle;
\fill[pthin] (-0.180,-3.152) -- (-0.180,-2.845) -- (-0.000,-2.904) -- cycle;
\fill[pthin] (-0.000,-2.597) -- (-0.180,-2.845) -- (-0.000,-2.904) -- cycle;
\fill[pthick] (-0.180,-2.845) -- (-0.000,-2.597) -- (-0.472,-2.751) -- cycle;
\fill[pthick] (-0.292,-1.700) -- (-0.472,-1.948) -- (-0.000,-1.795) -- cycle;
\fill[pthick] (-0.000,-1.605) -- (-0.292,-1.700) -- (-0.000,-1.299) -- cycle;
\fill[pthin] (-0.292,-1.700) -- (-0.000,-1.605) -- (-0.000,-1.795) -- cycle;
\fill[pthick] (-0.000,-2.291) -- (-0.000,-2.597) -- (-0.292,-2.196) -- cycle;
\fill[pthick] (-0.180,-2.043) -- (-0.000,-2.291) -- (-0.000,-1.795) -- cycle;
\fill[pthin] (-0.000,-2.291) -- (-0.180,-2.043) -- (-0.292,-2.196) -- cycle;
\fill[pthin] (-0.472,-1.948) -- (-0.180,-2.043) -- (-0.292,-2.196) -- cycle;
\fill[pthick] (-0.180,-2.043) -- (-0.472,-1.948) -- (-0.000,-1.795) -- cycle;
\fill[pthin] (-0.472,-1.948) -- (-0.472,-2.255) -- (-0.292,-2.196) -- cycle;
\fill[pthick] (-0.472,-2.255) -- (-0.472,-1.948) -- (-0.763,-2.349) -- cycle;
\fill[pthick] (-0.292,-2.503) -- (-0.000,-2.597) -- (-0.292,-2.196) -- cycle;
\fill[pthick] (-0.472,-2.255) -- (-0.292,-2.503) -- (-0.763,-2.349) -- cycle;
\fill[pthin] (-0.292,-2.503) -- (-0.472,-2.255) -- (-0.292,-2.196) -- cycle;
\fill[pthin] (-0.472,-1.948) -- (-0.763,-2.043) -- (-0.763,-1.853) -- cycle;
\fill[pthick] (-0.763,-2.043) -- (-0.472,-1.948) -- (-0.763,-2.349) -- cycle;
\fill[pthick] (-1.055,-1.948) -- (-1.235,-1.700) -- (-0.763,-1.853) -- cycle;
\fill[pthick] (-0.763,-2.043) -- (-1.055,-1.948) -- (-0.763,-2.349) -- cycle;
\fill[pthin] (-1.055,-1.948) -- (-0.763,-2.043) -- (-0.763,-1.853) -- cycle;
\fill[pthick] (-0.292,-1.700) -- (-0.472,-1.948) -- (-0.472,-1.452) -- cycle;
\fill[pthick] (-0.292,-1.393) -- (-0.292,-1.700) -- (-0.000,-1.299) -- cycle;
\fill[pthin] (-0.292,-1.700) -- (-0.292,-1.393) -- (-0.472,-1.452) -- cycle;
\fill[pthick] (-0.944,-1.605) -- (-1.235,-1.700) -- (-0.763,-1.853) -- cycle;
\fill[pthick] (-0.652,-1.700) -- (-0.944,-1.605) -- (-0.472,-1.452) -- cycle;
\fill[pthin] (-0.944,-1.605) -- (-0.652,-1.700) -- (-0.763,-1.853) -- cycle;
\fill[pthin] (-0.472,-1.948) -- (-0.652,-1.700) -- (-0.763,-1.853) -- cycle;
\fill[pthick] (-0.652,-1.700) -- (-0.472,-1.948) -- (-0.472,-1.452) -- cycle;
\fill[pthin] (1.235,-2.503) -- (1.055,-2.255) -- (1.235,-2.196) -- cycle;
\fill[pthick] (1.055,-2.255) -- (1.235,-2.503) -- (0.763,-2.349) -- cycle;
\fill[pthick] (1.055,-1.948) -- (1.235,-1.700) -- (1.235,-2.196) -- cycle;
\fill[pthick] (1.055,-2.255) -- (1.055,-1.948) -- (0.763,-2.349) -- cycle;
\fill[pthin] (1.055,-1.948) -- (1.055,-2.255) -- (1.235,-2.196) -- cycle;
\fill[pthick] (1.527,-2.597) -- (1.235,-2.503) -- (1.707,-2.349) -- cycle;
\fill[pthick] (1.818,-2.503) -- (1.527,-2.597) -- (1.998,-2.751) -- cycle;
\fill[pthin] (1.527,-2.597) -- (1.818,-2.503) -- (1.707,-2.349) -- cycle;
\fill[pthick] (1.415,-1.948) -- (1.235,-1.700) -- (1.235,-2.196) -- cycle;
\fill[pthick] (1.415,-2.255) -- (1.415,-1.948) -- (1.707,-2.349) -- cycle;
\fill[pthin] (1.415,-1.948) -- (1.415,-2.255) -- (1.235,-2.196) -- cycle;
\fill[pthin] (1.235,-2.503) -- (1.415,-2.255) -- (1.235,-2.196) -- cycle;
\fill[pthick] (1.415,-2.255) -- (1.235,-2.503) -- (1.707,-2.349) -- cycle;
\fill[pthick] (0.763,-2.845) -- (0.763,-3.152) -- (0.472,-2.751) -- cycle;
\fill[pthick] (0.583,-2.597) -- (0.763,-2.845) -- (0.763,-2.349) -- cycle;
\fill[pthin] (0.763,-2.845) -- (0.583,-2.597) -- (0.472,-2.751) -- cycle;
\fill[pthick] (0.180,-3.152) -- (-0.000,-3.400) -- (0.472,-3.247) -- cycle;
\fill[pthick] (0.472,-3.057) -- (0.180,-3.152) -- (0.472,-2.751) -- cycle;
\fill[pthin] (0.180,-3.152) -- (0.472,-3.057) -- (0.472,-3.247) -- cycle;
\fill[pthin] (0.763,-3.152) -- (0.472,-3.057) -- (0.472,-3.247) -- cycle;
\fill[pthick] (0.472,-3.057) -- (0.763,-3.152) -- (0.472,-2.751) -- cycle;
\fill[pthick] (1.527,-2.597) -- (1.235,-2.503) -- (1.527,-2.904) -- cycle;
\fill[pthick] (1.707,-2.845) -- (1.527,-2.597) -- (1.998,-2.751) -- cycle;
\fill[pthin] (1.527,-2.597) -- (1.707,-2.845) -- (1.527,-2.904) -- cycle;
\fill[pthick] (1.055,-3.057) -- (0.763,-3.152) -- (1.055,-2.751) -- cycle;
\fill[pthick] (1.235,-2.809) -- (1.055,-3.057) -- (1.527,-2.904) -- cycle;
\fill[pthin] (1.055,-3.057) -- (1.235,-2.809) -- (1.055,-2.751) -- cycle;
\fill[pthin] (1.235,-2.503) -- (1.235,-2.809) -- (1.055,-2.751) -- cycle;
\fill[pthick] (1.235,-2.809) -- (1.235,-2.503) -- (1.527,-2.904) -- cycle;
\fill[pthin] (1.235,-2.503) -- (0.944,-2.597) -- (1.055,-2.751) -- cycle;
\fill[pthick] (0.944,-2.597) -- (1.235,-2.503) -- (0.763,-2.349) -- cycle;
\fill[pthick] (0.763,-2.845) -- (0.763,-3.152) -- (1.055,-2.751) -- cycle;
\fill[pthick] (0.944,-2.597) -- (0.763,-2.845) -- (0.763,-2.349) -- cycle;
\fill[pthin] (0.763,-2.845) -- (0.944,-2.597) -- (1.055,-2.751) -- cycle;
\fill[pthick] (0.472,-1.145) -- (0.763,-1.051) -- (0.472,-1.452) -- cycle;
\fill[pthick] (0.292,-1.393) -- (0.472,-1.145) -- (-0.000,-1.299) -- cycle;
\fill[pthin] (0.472,-1.145) -- (0.292,-1.393) -- (0.472,-1.452) -- cycle;
\fill[pthick] (0.944,-1.605) -- (1.235,-1.700) -- (0.944,-1.299) -- cycle;
\fill[pthick] (0.763,-1.357) -- (0.944,-1.605) -- (0.472,-1.452) -- cycle;
\fill[pthin] (0.944,-1.605) -- (0.763,-1.357) -- (0.944,-1.299) -- cycle;
\fill[pthin] (0.763,-1.051) -- (0.763,-1.357) -- (0.944,-1.299) -- cycle;
\fill[pthick] (0.763,-1.357) -- (0.763,-1.051) -- (0.472,-1.452) -- cycle;
\fill[pthick] (-0.000,-0.496) -- (-0.000,-0.803) -- (0.292,-0.401) -- cycle;
\fill[pthick] (0.180,-0.248) -- (-0.000,-0.496) -- (0.000,0.000) -- cycle;
\fill[pthin] (-0.000,-0.496) -- (0.180,-0.248) -- (0.292,-0.401) -- cycle;
\fill[pthick] (0.583,-0.803) -- (0.763,-1.051) -- (0.292,-0.897) -- cycle;
\fill[pthick] (0.292,-0.708) -- (0.583,-0.803) -- (0.292,-0.401) -- cycle;
\fill[pthin] (0.583,-0.803) -- (0.292,-0.708) -- (0.292,-0.897) -- cycle;
\fill[pthin] (-0.000,-0.803) -- (0.292,-0.708) -- (0.292,-0.897) -- cycle;
\fill[pthick] (0.292,-0.708) -- (-0.000,-0.803) -- (0.292,-0.401) -- cycle;
\fill[pthin] (-0.000,-0.803) -- (0.180,-1.051) -- (0.292,-0.897) -- cycle;
\fill[pthick] (0.180,-1.051) -- (-0.000,-0.803) -- (-0.000,-1.299) -- cycle;
\fill[pthick] (0.472,-1.145) -- (0.763,-1.051) -- (0.292,-0.897) -- cycle;
\fill[pthick] (0.180,-1.051) -- (0.472,-1.145) -- (-0.000,-1.299) -- cycle;
\fill[pthin] (0.472,-1.145) -- (0.180,-1.051) -- (0.292,-0.897) -- cycle;
\fill[pthick] (0.292,-2.503) -- (-0.000,-2.597) -- (0.472,-2.751) -- cycle;
\fill[pthick] (0.583,-2.597) -- (0.292,-2.503) -- (0.763,-2.349) -- cycle;
\fill[pthin] (0.292,-2.503) -- (0.583,-2.597) -- (0.472,-2.751) -- cycle;
\fill[pthick] (0.180,-3.152) -- (-0.000,-3.400) -- (-0.000,-2.904) -- cycle;
\fill[pthick] (0.180,-2.845) -- (0.180,-3.152) -- (0.472,-2.751) -- cycle;
\fill[pthin] (0.180,-3.152) -- (0.180,-2.845) -- (-0.000,-2.904) -- cycle;
\fill[pthin] (-0.000,-2.597) -- (0.180,-2.845) -- (-0.000,-2.904) -- cycle;
\fill[pthick] (0.180,-2.845) -- (-0.000,-2.597) -- (0.472,-2.751) -- cycle;
\fill[pthick] (0.292,-1.700) -- (0.472,-1.948) -- (-0.000,-1.795) -- cycle;
\fill[pthick] (-0.000,-1.605) -- (0.292,-1.700) -- (-0.000,-1.299) -- cycle;
\fill[pthin] (0.292,-1.700) -- (-0.000,-1.605) -- (-0.000,-1.795) -- cycle;
\fill[pthick] (-0.000,-2.291) -- (-0.000,-2.597) -- (0.292,-2.196) -- cycle;
\fill[pthick] (0.180,-2.043) -- (-0.000,-2.291) -- (-0.000,-1.795) -- cycle;
\fill[pthin] (-0.000,-2.291) -- (0.180,-2.043) -- (0.292,-2.196) -- cycle;
\fill[pthin] (0.472,-1.948) -- (0.180,-2.043) -- (0.292,-2.196) -- cycle;
\fill[pthick] (0.180,-2.043) -- (0.472,-1.948) -- (-0.000,-1.795) -- cycle;
\fill[pthin] (0.472,-1.948) -- (0.472,-2.255) -- (0.292,-2.196) -- cycle;
\fill[pthick] (0.472,-2.255) -- (0.472,-1.948) -- (0.763,-2.349) -- cycle;
\fill[pthick] (0.292,-2.503) -- (-0.000,-2.597) -- (0.292,-2.196) -- cycle;
\fill[pthick] (0.472,-2.255) -- (0.292,-2.503) -- (0.763,-2.349) -- cycle;
\fill[pthin] (0.292,-2.503) -- (0.472,-2.255) -- (0.292,-2.196) -- cycle;
\fill[pthin] (0.472,-1.948) -- (0.763,-2.043) -- (0.763,-1.853) -- cycle;
\fill[pthick] (0.763,-2.043) -- (0.472,-1.948) -- (0.763,-2.349) -- cycle;
\fill[pthick] (1.055,-1.948) -- (1.235,-1.700) -- (0.763,-1.853) -- cycle;
\fill[pthick] (0.763,-2.043) -- (1.055,-1.948) -- (0.763,-2.349) -- cycle;
\fill[pthin] (1.055,-1.948) -- (0.763,-2.043) -- (0.763,-1.853) -- cycle;
\fill[pthick] (0.292,-1.700) -- (0.472,-1.948) -- (0.472,-1.452) -- cycle;
\fill[pthick] (0.292,-1.393) -- (0.292,-1.700) -- (-0.000,-1.299) -- cycle;
\fill[pthin] (0.292,-1.700) -- (0.292,-1.393) -- (0.472,-1.452) -- cycle;
\fill[pthick] (0.944,-1.605) -- (1.235,-1.700) -- (0.763,-1.853) -- cycle;
\fill[pthick] (0.652,-1.700) -- (0.944,-1.605) -- (0.472,-1.452) -- cycle;
\fill[pthin] (0.944,-1.605) -- (0.652,-1.700) -- (0.763,-1.853) -- cycle;
\fill[pthin] (0.472,-1.948) -- (0.652,-1.700) -- (0.763,-1.853) -- cycle;
\fill[pthick] (0.652,-1.700) -- (0.472,-1.948) -- (0.472,-1.452) -- cycle;
\fill[pthin] (1.998,-1.948) -- (1.818,-1.700) -- (1.707,-1.853) -- cycle;
\fill[pthick] (1.818,-1.700) -- (1.998,-1.948) -- (1.998,-1.452) -- cycle;
\fill[pthick] (1.527,-1.605) -- (1.235,-1.700) -- (1.707,-1.853) -- cycle;
\fill[pthick] (1.818,-1.700) -- (1.527,-1.605) -- (1.998,-1.452) -- cycle;
\fill[pthin] (1.527,-1.605) -- (1.818,-1.700) -- (1.707,-1.853) -- cycle;
\fill[pthick] (1.998,-2.255) -- (1.998,-1.948) -- (1.707,-2.349) -- cycle;
\fill[pthick] (1.818,-2.503) -- (1.998,-2.255) -- (1.998,-2.751) -- cycle;
\fill[pthin] (1.998,-2.255) -- (1.818,-2.503) -- (1.707,-2.349) -- cycle;
\fill[pthick] (1.415,-1.948) -- (1.235,-1.700) -- (1.707,-1.853) -- cycle;
\fill[pthick] (1.707,-2.043) -- (1.415,-1.948) -- (1.707,-2.349) -- cycle;
\fill[pthin] (1.415,-1.948) -- (1.707,-2.043) -- (1.707,-1.853) -- cycle;
\fill[pthin] (1.998,-1.948) -- (1.707,-2.043) -- (1.707,-1.853) -- cycle;
\fill[pthick] (1.707,-2.043) -- (1.998,-1.948) -- (1.707,-2.349) -- cycle;
\fill[pthick] (2.470,-1.605) -- (2.762,-1.700) -- (2.470,-1.299) -- cycle;
\fill[pthick] (2.290,-1.357) -- (2.470,-1.605) -- (1.998,-1.452) -- cycle;
\fill[pthin] (2.470,-1.605) -- (2.290,-1.357) -- (2.470,-1.299) -- cycle;
\fill[pthick] (2.942,-1.145) -- (3.234,-1.051) -- (2.942,-1.452) -- cycle;
\fill[pthick] (2.762,-1.393) -- (2.942,-1.145) -- (2.470,-1.299) -- cycle;
\fill[pthin] (2.942,-1.145) -- (2.762,-1.393) -- (2.942,-1.452) -- cycle;
\fill[pthin] (2.762,-1.700) -- (2.762,-1.393) -- (2.942,-1.452) -- cycle;
\fill[pthick] (2.762,-1.393) -- (2.762,-1.700) -- (2.470,-1.299) -- cycle;
\fill[pthick] (1.998,-2.255) -- (1.998,-1.948) -- (2.290,-2.349) -- cycle;
\fill[pthick] (2.179,-2.503) -- (1.998,-2.255) -- (1.998,-2.751) -- cycle;
\fill[pthin] (1.998,-2.255) -- (2.179,-2.503) -- (2.290,-2.349) -- cycle;
\fill[pthick] (2.582,-1.948) -- (2.762,-1.700) -- (2.290,-1.853) -- cycle;
\fill[pthick] (2.290,-2.043) -- (2.582,-1.948) -- (2.290,-2.349) -- cycle;
\fill[pthin] (2.582,-1.948) -- (2.290,-2.043) -- (2.290,-1.853) -- cycle;
\fill[pthin] (1.998,-1.948) -- (2.290,-2.043) -- (2.290,-1.853) -- cycle;
\fill[pthick] (2.290,-2.043) -- (1.998,-1.948) -- (2.290,-2.349) -- cycle;
\fill[pthin] (1.998,-1.948) -- (2.179,-1.700) -- (2.290,-1.853) -- cycle;
\fill[pthick] (2.179,-1.700) -- (1.998,-1.948) -- (1.998,-1.452) -- cycle;
\fill[pthick] (2.470,-1.605) -- (2.762,-1.700) -- (2.290,-1.853) -- cycle;
\fill[pthick] (2.179,-1.700) -- (2.470,-1.605) -- (1.998,-1.452) -- cycle;
\fill[pthin] (2.470,-1.605) -- (2.179,-1.700) -- (2.290,-1.853) -- cycle;
\fill[pthick] (0.944,-0.803) -- (0.763,-1.051) -- (1.235,-0.897) -- cycle;
\fill[pthick] (1.235,-0.708) -- (0.944,-0.803) -- (1.235,-0.401) -- cycle;
\fill[pthin] (0.944,-0.803) -- (1.235,-0.708) -- (1.235,-0.897) -- cycle;
\fill[pthick] (1.235,-1.393) -- (1.235,-1.700) -- (0.944,-1.299) -- cycle;
\fill[pthick] (1.055,-1.145) -- (1.235,-1.393) -- (1.235,-0.897) -- cycle;
\fill[pthin] (1.235,-1.393) -- (1.055,-1.145) -- (0.944,-1.299) -- cycle;
\fill[pthin] (0.763,-1.051) -- (1.055,-1.145) -- (0.944,-1.299) -- cycle;
\fill[pthick] (1.055,-1.145) -- (0.763,-1.051) -- (1.235,-0.897) -- cycle;
\fill[pthick] (0.472,-0.153) -- (0.763,-0.248) -- (0.292,-0.401) -- cycle;
\fill[pthick] (0.180,-0.248) -- (0.472,-0.153) -- (0.000,0.000) -- cycle;
\fill[pthin] (0.472,-0.153) -- (0.180,-0.248) -- (0.292,-0.401) -- cycle;
\fill[pthick] (0.583,-0.803) -- (0.763,-1.051) -- (0.763,-0.555) -- cycle;
\fill[pthick] (0.583,-0.496) -- (0.583,-0.803) -- (0.292,-0.401) -- cycle;
\fill[pthin] (0.583,-0.803) -- (0.583,-0.496) -- (0.763,-0.555) -- cycle;
\fill[pthin] (0.763,-0.248) -- (0.583,-0.496) -- (0.763,-0.555) -- cycle;
\fill[pthick] (0.583,-0.496) -- (0.763,-0.248) -- (0.292,-0.401) -- cycle;
\fill[pthin] (0.763,-0.248) -- (0.944,-0.496) -- (0.763,-0.555) -- cycle;
\fill[pthick] (0.944,-0.496) -- (0.763,-0.248) -- (1.235,-0.401) -- cycle;
\fill[pthick] (0.944,-0.803) -- (0.763,-1.051) -- (0.763,-0.555) -- cycle;
\fill[pthick] (0.944,-0.496) -- (0.944,-0.803) -- (1.235,-0.401) -- cycle;
\fill[pthin] (0.944,-0.803) -- (0.944,-0.496) -- (0.763,-0.555) -- cycle;
\fill[pthick] (2.290,-1.051) -- (2.470,-0.803) -- (2.470,-1.299) -- cycle;
\fill[pthick] (2.290,-1.357) -- (2.290,-1.051) -- (1.998,-1.452) -- cycle;
\fill[pthin] (2.290,-1.051) -- (2.290,-1.357) -- (2.470,-1.299) -- cycle;
\fill[pthick] (2.942,-1.145) -- (3.234,-1.051) -- (2.762,-0.897) -- cycle;
\fill[pthick] (2.650,-1.051) -- (2.942,-1.145) -- (2.470,-1.299) -- cycle;
\fill[pthin] (2.942,-1.145) -- (2.650,-1.051) -- (2.762,-0.897) -- cycle;
\fill[pthin] (2.470,-0.803) -- (2.650,-1.051) -- (2.762,-0.897) -- cycle;
\fill[pthick] (2.650,-1.051) -- (2.470,-0.803) -- (2.470,-1.299) -- cycle;
\fill[pthick] (1.527,-0.803) -- (1.707,-1.051) -- (1.707,-0.555) -- cycle;
\fill[pthick] (1.527,-0.496) -- (1.527,-0.803) -- (1.235,-0.401) -- cycle;
\fill[pthin] (1.527,-0.803) -- (1.527,-0.496) -- (1.707,-0.555) -- cycle;
\fill[pthick] (2.179,-0.708) -- (2.470,-0.803) -- (1.998,-0.956) -- cycle;
\fill[pthick] (1.887,-0.803) -- (2.179,-0.708) -- (1.707,-0.555) -- cycle;
\fill[pthin] (2.179,-0.708) -- (1.887,-0.803) -- (1.998,-0.956) -- cycle;
\fill[pthin] (1.707,-1.051) -- (1.887,-0.803) -- (1.998,-0.956) -- cycle;
\fill[pthick] (1.887,-0.803) -- (1.707,-1.051) -- (1.707,-0.555) -- cycle;
\fill[pthin] (1.707,-1.051) -- (1.998,-1.145) -- (1.998,-0.956) -- cycle;
\fill[pthick] (1.998,-1.145) -- (1.707,-1.051) -- (1.998,-1.452) -- cycle;
\fill[pthick] (2.290,-1.051) -- (2.470,-0.803) -- (1.998,-0.956) -- cycle;
\fill[pthick] (1.998,-1.145) -- (2.290,-1.051) -- (1.998,-1.452) -- cycle;
\fill[pthin] (2.290,-1.051) -- (1.998,-1.145) -- (1.998,-0.956) -- cycle;
\fill[pthin] (1.707,-1.051) -- (1.707,-1.357) -- (1.527,-1.299) -- cycle;
\fill[pthick] (1.707,-1.357) -- (1.707,-1.051) -- (1.998,-1.452) -- cycle;
\fill[pthick] (1.527,-1.605) -- (1.235,-1.700) -- (1.527,-1.299) -- cycle;
\fill[pthick] (1.707,-1.357) -- (1.527,-1.605) -- (1.998,-1.452) -- cycle;
\fill[pthin] (1.527,-1.605) -- (1.707,-1.357) -- (1.527,-1.299) -- cycle;
\fill[pthick] (1.527,-0.803) -- (1.707,-1.051) -- (1.235,-0.897) -- cycle;
\fill[pthick] (1.235,-0.708) -- (1.527,-0.803) -- (1.235,-0.401) -- cycle;
\fill[pthin] (1.527,-0.803) -- (1.235,-0.708) -- (1.235,-0.897) -- cycle;
\fill[pthick] (1.235,-1.393) -- (1.235,-1.700) -- (1.527,-1.299) -- cycle;
\fill[pthick] (1.415,-1.145) -- (1.235,-1.393) -- (1.235,-0.897) -- cycle;
\fill[pthin] (1.235,-1.393) -- (1.415,-1.145) -- (1.527,-1.299) -- cycle;
\fill[pthin] (1.707,-1.051) -- (1.415,-1.145) -- (1.527,-1.299) -- cycle;
\fill[pthick] (1.415,-1.145) -- (1.707,-1.051) -- (1.235,-0.897) -- cycle;
\draw[pedge] (2.470,0.496) -- (2.762,0.401) -- (2.470,0.307);
\draw[pedge] (2.470,0.000) -- (2.470,0.307) -- (2.762,0.401);
\draw[pedge] (2.470,0.496) -- (2.179,0.401) -- (1.998,0.649);
\draw[pedge] (2.470,0.000) -- (2.470,0.307) -- (2.179,0.401);
\draw[pedge] (2.470,0.496) -- (2.179,0.401) -- (2.470,0.307);
\draw[pedge] (2.762,0.897) -- (2.942,0.649) -- (2.762,0.401);
\draw[pedge] (3.234,1.051) -- (2.942,0.956) -- (2.942,0.649);
\draw[pedge] (2.762,0.897) -- (2.942,0.649) -- (2.942,0.956);
\draw[pedge] (2.470,0.496) -- (2.290,0.744) -- (1.998,0.649);
\draw[pedge] (2.762,0.897) -- (2.582,0.649) -- (2.290,0.744);
\draw[pedge] (2.470,0.496) -- (2.290,0.744) -- (2.582,0.649);
\draw[pedge] (2.470,0.496) -- (2.762,0.401) -- (2.582,0.649);
\draw[pedge] (2.762,0.897) -- (2.582,0.649) -- (2.762,0.401);
\draw[pedge] (2.762,-0.401) -- (2.942,-0.153) -- (3.234,-0.248);
\draw[pedge] (2.470,0.000) -- (2.650,-0.248) -- (2.942,-0.153);
\draw[pedge] (2.762,-0.401) -- (2.942,-0.153) -- (2.650,-0.248);
\draw[pedge] (3.234,-0.555) -- (3.053,-0.803) -- (3.234,-1.051);
\draw[pedge] (2.762,-0.401) -- (3.053,-0.496) -- (3.053,-0.803);
\draw[pedge] (3.234,-0.555) -- (3.053,-0.803) -- (3.053,-0.496);
\draw[pedge] (3.234,-0.555) -- (3.234,-0.248) -- (3.053,-0.496);
\draw[pedge] (2.762,-0.401) -- (3.053,-0.496) -- (3.234,-0.248);
\draw[pedge] (3.234,0.555) -- (2.942,0.649) -- (2.762,0.401);
\draw[pedge] (3.234,1.051) -- (3.234,0.744) -- (2.942,0.649);
\draw[pedge] (3.234,0.555) -- (2.942,0.649) -- (3.234,0.744);
\draw[pedge] (2.942,0.153) -- (3.234,0.059) -- (3.234,-0.248);
\draw[pedge] (3.234,0.555) -- (3.053,0.307) -- (3.234,0.059);
\draw[pedge] (2.942,0.153) -- (3.234,0.059) -- (3.053,0.307);
\draw[pedge] (2.942,0.153) -- (2.762,0.401) -- (3.053,0.307);
\draw[pedge] (3.234,0.555) -- (3.053,0.307) -- (2.762,0.401);
\draw[pedge] (2.942,0.153) -- (2.762,0.401) -- (2.762,0.095);
\draw[pedge] (2.470,0.000) -- (2.762,0.095) -- (2.762,0.401);
\draw[pedge] (2.942,0.153) -- (2.942,-0.153) -- (3.234,-0.248);
\draw[pedge] (2.470,0.000) -- (2.762,0.095) -- (2.942,-0.153);
\draw[pedge] (2.942,0.153) -- (2.942,-0.153) -- (2.762,0.095);
\draw[pedge] (1.527,0.000) -- (1.235,0.095) -- (1.235,0.401);
\draw[pedge] (1.235,-0.401) -- (1.415,-0.153) -- (1.235,0.095);
\draw[pedge] (1.527,0.000) -- (1.235,0.095) -- (1.415,-0.153);
\draw[pedge] (1.527,0.496) -- (1.818,0.401) -- (1.998,0.649);
\draw[pedge] (1.527,0.000) -- (1.527,0.307) -- (1.818,0.401);
\draw[pedge] (1.527,0.496) -- (1.818,0.401) -- (1.527,0.307);
\draw[pedge] (1.527,0.496) -- (1.235,0.401) -- (1.527,0.307);
\draw[pedge] (1.527,0.000) -- (1.527,0.307) -- (1.235,0.401);
\draw[pedge] (0.472,0.153) -- (0.472,-0.153) -- (0.763,-0.248);
\draw[pedge] (0.000,0.000) -- (0.292,0.095) -- (0.472,-0.153);
\draw[pedge] (0.472,0.153) -- (0.472,-0.153) -- (0.292,0.095);
\draw[pedge] (0.944,0.000) -- (0.944,0.307) -- (1.235,0.401);
\draw[pedge] (0.472,0.153) -- (0.763,0.059) -- (0.944,0.307);
\draw[pedge] (0.944,0.000) -- (0.944,0.307) -- (0.763,0.059);
\draw[pedge] (0.944,0.000) -- (0.763,-0.248) -- (0.763,0.059);
\draw[pedge] (0.472,0.153) -- (0.763,0.059) -- (0.763,-0.248);
\draw[pedge] (0.944,0.000) -- (0.763,-0.248) -- (1.055,-0.153);
\draw[pedge] (1.235,-0.401) -- (1.055,-0.153) -- (0.763,-0.248);
\draw[pedge] (0.944,0.000) -- (1.235,0.095) -- (1.235,0.401);
\draw[pedge] (1.235,-0.401) -- (1.055,-0.153) -- (1.235,0.095);
\draw[pedge] (0.944,0.000) -- (1.235,0.095) -- (1.055,-0.153);
\draw[pedge] (2.762,-0.401) -- (2.470,-0.496) -- (2.470,-0.803);
\draw[pedge] (2.470,0.000) -- (2.650,-0.248) -- (2.470,-0.496);
\draw[pedge] (2.762,-0.401) -- (2.470,-0.496) -- (2.650,-0.248);
\draw[pedge] (2.762,-0.897) -- (3.053,-0.803) -- (3.234,-1.051);
\draw[pedge] (2.762,-0.401) -- (2.762,-0.708) -- (3.053,-0.803);
\draw[pedge] (2.762,-0.897) -- (3.053,-0.803) -- (2.762,-0.708);
\draw[pedge] (2.762,-0.897) -- (2.470,-0.803) -- (2.762,-0.708);
\draw[pedge] (2.762,-0.401) -- (2.762,-0.708) -- (2.470,-0.803);
\draw[pedge] (1.707,-0.555) -- (1.707,-0.248) -- (1.998,-0.153);
\draw[pedge] (1.235,-0.401) -- (1.527,-0.496) -- (1.707,-0.248);
\draw[pedge] (1.707,-0.555) -- (1.707,-0.248) -- (1.527,-0.496);
\draw[pedge] (2.179,-0.401) -- (2.179,-0.708) -- (2.470,-0.803);
\draw[pedge] (1.707,-0.555) -- (1.998,-0.460) -- (2.179,-0.708);
\draw[pedge] (2.179,-0.401) -- (2.179,-0.708) -- (1.998,-0.460);
\draw[pedge] (2.179,-0.401) -- (1.998,-0.153) -- (1.998,-0.460);
\draw[pedge] (1.707,-0.555) -- (1.998,-0.460) -- (1.998,-0.153);
\draw[pedge] (2.179,-0.401) -- (1.998,-0.153) -- (2.290,-0.248);
\draw[pedge] (2.470,0.000) -- (2.290,-0.248) -- (1.998,-0.153);
\draw[pedge] (2.179,-0.401) -- (2.470,-0.496) -- (2.470,-0.803);
\draw[pedge] (2.470,0.000) -- (2.290,-0.248) -- (2.470,-0.496);
\draw[pedge] (2.179,-0.401) -- (2.470,-0.496) -- (2.290,-0.248);
\draw[pedge] (1.998,0.153) -- (1.998,-0.153) -- (2.179,0.095);
\draw[pedge] (2.470,0.000) -- (2.179,0.095) -- (1.998,-0.153);
\draw[pedge] (1.998,0.153) -- (2.179,0.401) -- (1.998,0.649);
\draw[pedge] (2.470,0.000) -- (2.179,0.095) -- (2.179,0.401);
\draw[pedge] (1.998,0.153) -- (2.179,0.401) -- (2.179,0.095);
\draw[pedge] (1.527,0.000) -- (1.707,-0.248) -- (1.998,-0.153);
\draw[pedge] (1.235,-0.401) -- (1.415,-0.153) -- (1.707,-0.248);
\draw[pedge] (1.527,0.000) -- (1.707,-0.248) -- (1.415,-0.153);
\draw[pedge] (1.998,0.153) -- (1.818,0.401) -- (1.998,0.649);
\draw[pedge] (1.527,0.000) -- (1.818,0.095) -- (1.818,0.401);
\draw[pedge] (1.998,0.153) -- (1.818,0.401) -- (1.818,0.095);
\draw[pedge] (1.998,0.153) -- (1.998,-0.153) -- (1.818,0.095);
\draw[pedge] (1.527,0.000) -- (1.818,0.095) -- (1.998,-0.153);
\draw[pedge] (2.290,1.051) -- (2.470,1.299) -- (2.179,1.204);
\draw[pedge] (1.998,1.452) -- (2.179,1.204) -- (2.470,1.299);
\draw[pedge] (2.290,1.051) -- (1.998,0.956) -- (1.998,0.649);
\draw[pedge] (1.998,1.452) -- (2.179,1.204) -- (1.998,0.956);
\draw[pedge] (2.290,1.051) -- (1.998,0.956) -- (2.179,1.204);
\draw[pedge] (2.762,0.897) -- (2.762,1.204) -- (2.470,1.299);
\draw[pedge] (3.234,1.051) -- (2.942,0.956) -- (2.762,1.204);
\draw[pedge] (2.762,0.897) -- (2.762,1.204) -- (2.942,0.956);
\draw[pedge] (2.290,1.051) -- (2.290,0.744) -- (1.998,0.649);
\draw[pedge] (2.762,0.897) -- (2.470,0.992) -- (2.290,0.744);
\draw[pedge] (2.290,1.051) -- (2.290,0.744) -- (2.470,0.992);
\draw[pedge] (2.290,1.051) -- (2.470,1.299) -- (2.470,0.992);
\draw[pedge] (2.762,0.897) -- (2.470,0.992) -- (2.470,1.299);
\draw[pedge] (1.998,1.948) -- (2.290,1.853) -- (2.470,2.101);
\draw[pedge] (1.998,1.452) -- (1.998,1.759) -- (2.290,1.853);
\draw[pedge] (1.998,1.948) -- (2.290,1.853) -- (1.998,1.759);
\draw[pedge] (2.290,2.349) -- (1.998,2.444) -- (1.998,2.751);
\draw[pedge] (1.998,1.948) -- (2.179,2.196) -- (1.998,2.444);
\draw[pedge] (2.290,2.349) -- (1.998,2.444) -- (2.179,2.196);
\draw[pedge] (2.290,2.349) -- (2.470,2.101) -- (2.179,2.196);
\draw[pedge] (1.998,1.948) -- (2.179,2.196) -- (2.470,2.101);
\draw[pedge] (2.942,1.452) -- (2.762,1.204) -- (2.470,1.299);
\draw[pedge] (3.234,1.051) -- (3.053,1.299) -- (2.762,1.204);
\draw[pedge] (2.942,1.452) -- (2.762,1.204) -- (3.053,1.299);
\draw[pedge] (2.470,1.605) -- (2.650,1.853) -- (2.470,2.101);
\draw[pedge] (2.942,1.452) -- (2.650,1.547) -- (2.650,1.853);
\draw[pedge] (2.470,1.605) -- (2.650,1.853) -- (2.650,1.547);
\draw[pedge] (2.470,1.605) -- (2.470,1.299) -- (2.650,1.547);
\draw[pedge] (2.942,1.452) -- (2.650,1.547) -- (2.470,1.299);
\draw[pedge] (2.470,1.605) -- (2.470,1.299) -- (2.290,1.547);
\draw[pedge] (1.998,1.452) -- (2.290,1.547) -- (2.470,1.299);
\draw[pedge] (2.470,1.605) -- (2.290,1.853) -- (2.470,2.101);
\draw[pedge] (1.998,1.452) -- (2.290,1.547) -- (2.290,1.853);
\draw[pedge] (2.470,1.605) -- (2.290,1.853) -- (2.290,1.547);
\draw[pedge] (1.235,0.897) -- (1.055,0.649) -- (1.235,0.401);
\draw[pedge] (0.763,1.051) -- (1.055,0.956) -- (1.055,0.649);
\draw[pedge] (1.235,0.897) -- (1.055,0.649) -- (1.055,0.956);
\draw[pedge] (1.527,0.496) -- (1.707,0.744) -- (1.998,0.649);
\draw[pedge] (1.235,0.897) -- (1.415,0.649) -- (1.707,0.744);
\draw[pedge] (1.527,0.496) -- (1.707,0.744) -- (1.415,0.649);
\draw[pedge] (1.527,0.496) -- (1.235,0.401) -- (1.415,0.649);
\draw[pedge] (1.235,0.897) -- (1.415,0.649) -- (1.235,0.401);
\draw[pedge] (0.472,0.153) -- (0.292,0.401) -- (0.472,0.649);
\draw[pedge] (0.000,0.000) -- (0.292,0.095) -- (0.292,0.401);
\draw[pedge] (0.472,0.153) -- (0.292,0.401) -- (0.292,0.095);
\draw[pedge] (0.763,0.555) -- (0.944,0.307) -- (1.235,0.401);
\draw[pedge] (0.472,0.153) -- (0.652,0.401) -- (0.944,0.307);
\draw[pedge] (0.763,0.555) -- (0.944,0.307) -- (0.652,0.401);
\draw[pedge] (0.763,0.555) -- (0.472,0.649) -- (0.652,0.401);
\draw[pedge] (0.472,0.153) -- (0.652,0.401) -- (0.472,0.649);
\draw[pedge] (0.763,0.555) -- (0.472,0.649) -- (0.763,0.744);
\draw[pedge] (0.763,1.051) -- (0.763,0.744) -- (0.472,0.649);
\draw[pedge] (0.763,0.555) -- (1.055,0.649) -- (1.235,0.401);
\draw[pedge] (0.763,1.051) -- (0.763,0.744) -- (1.055,0.649);
\draw[pedge] (0.763,0.555) -- (1.055,0.649) -- (0.763,0.744);
\draw[pedge] (1.998,1.948) -- (1.707,1.853) -- (1.527,2.101);
\draw[pedge] (1.998,1.452) -- (1.998,1.759) -- (1.707,1.853);
\draw[pedge] (1.998,1.948) -- (1.707,1.853) -- (1.998,1.759);
\draw[pedge] (1.707,2.349) -- (1.998,2.444) -- (1.998,2.751);
\draw[pedge] (1.998,1.948) -- (1.818,2.196) -- (1.998,2.444);
\draw[pedge] (1.707,2.349) -- (1.998,2.444) -- (1.818,2.196);
\draw[pedge] (1.707,2.349) -- (1.527,2.101) -- (1.818,2.196);
\draw[pedge] (1.998,1.948) -- (1.818,2.196) -- (1.527,2.101);
\draw[pedge] (1.055,1.452) -- (1.235,1.204) -- (1.527,1.299);
\draw[pedge] (0.763,1.051) -- (0.944,1.299) -- (1.235,1.204);
\draw[pedge] (1.055,1.452) -- (1.235,1.204) -- (0.944,1.299);
\draw[pedge] (1.527,1.605) -- (1.346,1.853) -- (1.527,2.101);
\draw[pedge] (1.055,1.452) -- (1.346,1.547) -- (1.346,1.853);
\draw[pedge] (1.527,1.605) -- (1.346,1.853) -- (1.346,1.547);
\draw[pedge] (1.527,1.605) -- (1.527,1.299) -- (1.346,1.547);
\draw[pedge] (1.055,1.452) -- (1.346,1.547) -- (1.527,1.299);
\draw[pedge] (1.527,1.605) -- (1.527,1.299) -- (1.707,1.547);
\draw[pedge] (1.998,1.452) -- (1.707,1.547) -- (1.527,1.299);
\draw[pedge] (1.527,1.605) -- (1.707,1.853) -- (1.527,2.101);
\draw[pedge] (1.998,1.452) -- (1.707,1.547) -- (1.707,1.853);
\draw[pedge] (1.527,1.605) -- (1.707,1.853) -- (1.707,1.547);
\draw[pedge] (1.707,1.051) -- (1.527,1.299) -- (1.818,1.204);
\draw[pedge] (1.998,1.452) -- (1.818,1.204) -- (1.527,1.299);
\draw[pedge] (1.707,1.051) -- (1.998,0.956) -- (1.998,0.649);
\draw[pedge] (1.998,1.452) -- (1.818,1.204) -- (1.998,0.956);
\draw[pedge] (1.707,1.051) -- (1.998,0.956) -- (1.818,1.204);
\draw[pedge] (1.235,0.897) -- (1.235,1.204) -- (1.527,1.299);
\draw[pedge] (0.763,1.051) -- (1.055,0.956) -- (1.235,1.204);
\draw[pedge] (1.235,0.897) -- (1.235,1.204) -- (1.055,0.956);
\draw[pedge] (1.707,1.051) -- (1.707,0.744) -- (1.998,0.649);
\draw[pedge] (1.235,0.897) -- (1.527,0.992) -- (1.707,0.744);
\draw[pedge] (1.707,1.051) -- (1.707,0.744) -- (1.527,0.992);
\draw[pedge] (1.707,1.051) -- (1.527,1.299) -- (1.527,0.992);
\draw[pedge] (1.235,0.897) -- (1.527,0.992) -- (1.527,1.299);
\draw[pedge] (0.292,2.503) -- (0.472,2.751) -- (0.472,2.444);
\draw[pedge] (0.763,2.349) -- (0.472,2.444) -- (0.472,2.751);
\draw[pedge] (0.292,2.503) -- (0.292,2.196) -- (0.000,2.101);
\draw[pedge] (0.763,2.349) -- (0.472,2.444) -- (0.292,2.196);
\draw[pedge] (0.292,2.503) -- (0.292,2.196) -- (0.472,2.444);
\draw[pedge] (0.000,2.904) -- (0.292,2.999) -- (0.472,2.751);
\draw[pedge] (0.000,3.400) -- (0.000,3.093) -- (0.292,2.999);
\draw[pedge] (0.000,2.904) -- (0.292,2.999) -- (0.000,3.093);
\draw[pedge] (0.292,2.503) -- (0.000,2.408) -- (0.000,2.101);
\draw[pedge] (0.000,2.904) -- (0.180,2.656) -- (0.000,2.408);
\draw[pedge] (0.292,2.503) -- (0.000,2.408) -- (0.180,2.656);
\draw[pedge] (0.292,2.503) -- (0.472,2.751) -- (0.180,2.656);
\draw[pedge] (0.000,2.904) -- (0.180,2.656) -- (0.472,2.751);
\draw[pedge] (1.235,2.503) -- (1.055,2.751) -- (1.235,2.999);
\draw[pedge] (0.763,2.349) -- (1.055,2.444) -- (1.055,2.751);
\draw[pedge] (1.235,2.503) -- (1.055,2.751) -- (1.055,2.444);
\draw[pedge] (1.527,2.904) -- (1.707,2.656) -- (1.998,2.751);
\draw[pedge] (1.235,2.503) -- (1.415,2.751) -- (1.707,2.656);
\draw[pedge] (1.527,2.904) -- (1.707,2.656) -- (1.415,2.751);
\draw[pedge] (1.527,2.904) -- (1.235,2.999) -- (1.415,2.751);
\draw[pedge] (1.235,2.503) -- (1.415,2.751) -- (1.235,2.999);
\draw[pedge] (0.472,3.247) -- (0.292,2.999) -- (0.472,2.751);
\draw[pedge] (0.000,3.400) -- (0.292,3.305) -- (0.292,2.999);
\draw[pedge] (0.472,3.247) -- (0.292,2.999) -- (0.292,3.305);
\draw[pedge] (0.763,2.845) -- (0.944,3.093) -- (1.235,2.999);
\draw[pedge] (0.472,3.247) -- (0.652,2.999) -- (0.944,3.093);
\draw[pedge] (0.763,2.845) -- (0.944,3.093) -- (0.652,2.999);
\draw[pedge] (0.763,2.845) -- (0.472,2.751) -- (0.652,2.999);
\draw[pedge] (0.472,3.247) -- (0.652,2.999) -- (0.472,2.751);
\draw[pedge] (0.763,2.845) -- (0.472,2.751) -- (0.763,2.656);
\draw[pedge] (0.763,2.349) -- (0.763,2.656) -- (0.472,2.751);
\draw[pedge] (0.763,2.845) -- (1.055,2.751) -- (1.235,2.999);
\draw[pedge] (0.763,2.349) -- (0.763,2.656) -- (1.055,2.751);
\draw[pedge] (0.763,2.845) -- (1.055,2.751) -- (0.763,2.656);
\draw[pedge] (0.472,1.452) -- (0.292,1.204) -- (0.000,1.299);
\draw[pedge] (0.763,1.051) -- (0.583,1.299) -- (0.292,1.204);
\draw[pedge] (0.472,1.452) -- (0.292,1.204) -- (0.583,1.299);
\draw[pedge] (0.000,1.605) -- (0.180,1.853) -- (0.000,2.101);
\draw[pedge] (0.472,1.452) -- (0.180,1.547) -- (0.180,1.853);
\draw[pedge] (0.000,1.605) -- (0.180,1.853) -- (0.180,1.547);
\draw[pedge] (0.000,1.605) -- (0.000,1.299) -- (0.180,1.547);
\draw[pedge] (0.472,1.452) -- (0.180,1.547) -- (0.000,1.299);
\draw[pedge] (0.000,0.496) -- (0.292,0.401) -- (0.472,0.649);
\draw[pedge] (0.000,0.000) -- (0.000,0.307) -- (0.292,0.401);
\draw[pedge] (0.000,0.496) -- (0.292,0.401) -- (0.000,0.307);
\draw[pedge] (0.292,0.897) -- (0.000,0.992) -- (0.000,1.299);
\draw[pedge] (0.000,0.496) -- (0.180,0.744) -- (0.000,0.992);
\draw[pedge] (0.292,0.897) -- (0.000,0.992) -- (0.180,0.744);
\draw[pedge] (0.292,0.897) -- (0.472,0.649) -- (0.180,0.744);
\draw[pedge] (0.000,0.496) -- (0.180,0.744) -- (0.472,0.649);
\draw[pedge] (0.292,0.897) -- (0.472,0.649) -- (0.472,0.956);
\draw[pedge] (0.763,1.051) -- (0.472,0.956) -- (0.472,0.649);
\draw[pedge] (0.292,0.897) -- (0.292,1.204) -- (0.000,1.299);
\draw[pedge] (0.763,1.051) -- (0.472,0.956) -- (0.292,1.204);
\draw[pedge] (0.292,0.897) -- (0.292,1.204) -- (0.472,0.956);
\draw[pedge] (1.235,2.503) -- (1.235,2.196) -- (1.527,2.101);
\draw[pedge] (0.763,2.349) -- (1.055,2.444) -- (1.235,2.196);
\draw[pedge] (1.235,2.503) -- (1.235,2.196) -- (1.055,2.444);
\draw[pedge] (1.707,2.349) -- (1.707,2.656) -- (1.998,2.751);
\draw[pedge] (1.235,2.503) -- (1.527,2.408) -- (1.707,2.656);
\draw[pedge] (1.707,2.349) -- (1.707,2.656) -- (1.527,2.408);
\draw[pedge] (1.707,2.349) -- (1.527,2.101) -- (1.527,2.408);
\draw[pedge] (1.235,2.503) -- (1.527,2.408) -- (1.527,2.101);
\draw[pedge] (1.055,1.452) -- (0.763,1.547) -- (0.763,1.853);
\draw[pedge] (0.763,1.051) -- (0.944,1.299) -- (0.763,1.547);
\draw[pedge] (1.055,1.452) -- (0.763,1.547) -- (0.944,1.299);
\draw[pedge] (1.055,1.948) -- (1.346,1.853) -- (1.527,2.101);
\draw[pedge] (1.055,1.452) -- (1.055,1.759) -- (1.346,1.853);
\draw[pedge] (1.055,1.948) -- (1.346,1.853) -- (1.055,1.759);
\draw[pedge] (1.055,1.948) -- (0.763,1.853) -- (1.055,1.759);
\draw[pedge] (1.055,1.452) -- (1.055,1.759) -- (0.763,1.853);
\draw[pedge] (1.055,1.948) -- (0.763,1.853) -- (0.944,2.101);
\draw[pedge] (0.763,2.349) -- (0.944,2.101) -- (0.763,1.853);
\draw[pedge] (1.055,1.948) -- (1.235,2.196) -- (1.527,2.101);
\draw[pedge] (0.763,2.349) -- (0.944,2.101) -- (1.235,2.196);
\draw[pedge] (1.055,1.948) -- (1.235,2.196) -- (0.944,2.101);
\draw[pedge] (0.472,1.948) -- (0.763,1.853) -- (0.583,2.101);
\draw[pedge] (0.763,2.349) -- (0.583,2.101) -- (0.763,1.853);
\draw[pedge] (0.472,1.948) -- (0.292,2.196) -- (0.000,2.101);
\draw[pedge] (0.763,2.349) -- (0.583,2.101) -- (0.292,2.196);
\draw[pedge] (0.472,1.948) -- (0.292,2.196) -- (0.583,2.101);
\draw[pedge] (0.472,1.452) -- (0.763,1.547) -- (0.763,1.853);
\draw[pedge] (0.763,1.051) -- (0.583,1.299) -- (0.763,1.547);
\draw[pedge] (0.472,1.452) -- (0.763,1.547) -- (0.583,1.299);
\draw[pedge] (0.472,1.948) -- (0.180,1.853) -- (0.000,2.101);
\draw[pedge] (0.472,1.452) -- (0.472,1.759) -- (0.180,1.853);
\draw[pedge] (0.472,1.948) -- (0.180,1.853) -- (0.472,1.759);
\draw[pedge] (0.472,1.948) -- (0.763,1.853) -- (0.472,1.759);
\draw[pedge] (0.472,1.452) -- (0.472,1.759) -- (0.763,1.853);
\draw[pedge] (-0.292,2.503) -- (-0.472,2.751) -- (-0.472,2.444);
\draw[pedge] (-0.763,2.349) -- (-0.472,2.444) -- (-0.472,2.751);
\draw[pedge] (-0.292,2.503) -- (-0.292,2.196) -- (0.000,2.101);
\draw[pedge] (-0.763,2.349) -- (-0.472,2.444) -- (-0.292,2.196);
\draw[pedge] (-0.292,2.503) -- (-0.292,2.196) -- (-0.472,2.444);
\draw[pedge] (0.000,2.904) -- (-0.292,2.999) -- (-0.472,2.751);
\draw[pedge] (0.000,3.400) -- (0.000,3.093) -- (-0.292,2.999);
\draw[pedge] (0.000,2.904) -- (-0.292,2.999) -- (0.000,3.093);
\draw[pedge] (-0.292,2.503) -- (0.000,2.408) -- (0.000,2.101);
\draw[pedge] (0.000,2.904) -- (-0.180,2.656) -- (0.000,2.408);
\draw[pedge] (-0.292,2.503) -- (0.000,2.408) -- (-0.180,2.656);
\draw[pedge] (-0.292,2.503) -- (-0.472,2.751) -- (-0.180,2.656);
\draw[pedge] (0.000,2.904) -- (-0.180,2.656) -- (-0.472,2.751);
\draw[pedge] (-1.235,2.503) -- (-1.055,2.751) -- (-1.235,2.999);
\draw[pedge] (-0.763,2.349) -- (-1.055,2.444) -- (-1.055,2.751);
\draw[pedge] (-1.235,2.503) -- (-1.055,2.751) -- (-1.055,2.444);
\draw[pedge] (-1.527,2.904) -- (-1.707,2.656) -- (-1.998,2.751);
\draw[pedge] (-1.235,2.503) -- (-1.415,2.751) -- (-1.707,2.656);
\draw[pedge] (-1.527,2.904) -- (-1.707,2.656) -- (-1.415,2.751);
\draw[pedge] (-1.527,2.904) -- (-1.235,2.999) -- (-1.415,2.751);
\draw[pedge] (-1.235,2.503) -- (-1.415,2.751) -- (-1.235,2.999);
\draw[pedge] (-0.472,3.247) -- (-0.292,2.999) -- (-0.472,2.751);
\draw[pedge] (0.000,3.400) -- (-0.292,3.305) -- (-0.292,2.999);
\draw[pedge] (-0.472,3.247) -- (-0.292,2.999) -- (-0.292,3.305);
\draw[pedge] (-0.763,2.845) -- (-0.944,3.093) -- (-1.235,2.999);
\draw[pedge] (-0.472,3.247) -- (-0.652,2.999) -- (-0.944,3.093);
\draw[pedge] (-0.763,2.845) -- (-0.944,3.093) -- (-0.652,2.999);
\draw[pedge] (-0.763,2.845) -- (-0.472,2.751) -- (-0.652,2.999);
\draw[pedge] (-0.472,3.247) -- (-0.652,2.999) -- (-0.472,2.751);
\draw[pedge] (-0.763,2.845) -- (-0.472,2.751) -- (-0.763,2.656);
\draw[pedge] (-0.763,2.349) -- (-0.763,2.656) -- (-0.472,2.751);
\draw[pedge] (-0.763,2.845) -- (-1.055,2.751) -- (-1.235,2.999);
\draw[pedge] (-0.763,2.349) -- (-0.763,2.656) -- (-1.055,2.751);
\draw[pedge] (-0.763,2.845) -- (-1.055,2.751) -- (-0.763,2.656);
\draw[pedge] (-0.472,1.452) -- (-0.292,1.204) -- (0.000,1.299);
\draw[pedge] (-0.763,1.051) -- (-0.583,1.299) -- (-0.292,1.204);
\draw[pedge] (-0.472,1.452) -- (-0.292,1.204) -- (-0.583,1.299);
\draw[pedge] (0.000,1.605) -- (-0.180,1.853) -- (0.000,2.101);
\draw[pedge] (-0.472,1.452) -- (-0.180,1.547) -- (-0.180,1.853);
\draw[pedge] (0.000,1.605) -- (-0.180,1.853) -- (-0.180,1.547);
\draw[pedge] (0.000,1.605) -- (0.000,1.299) -- (-0.180,1.547);
\draw[pedge] (-0.472,1.452) -- (-0.180,1.547) -- (0.000,1.299);
\draw[pedge] (0.000,0.496) -- (-0.292,0.401) -- (-0.472,0.649);
\draw[pedge] (0.000,0.000) -- (0.000,0.307) -- (-0.292,0.401);
\draw[pedge] (0.000,0.496) -- (-0.292,0.401) -- (0.000,0.307);
\draw[pedge] (-0.292,0.897) -- (0.000,0.992) -- (0.000,1.299);
\draw[pedge] (0.000,0.496) -- (-0.180,0.744) -- (0.000,0.992);
\draw[pedge] (-0.292,0.897) -- (0.000,0.992) -- (-0.180,0.744);
\draw[pedge] (-0.292,0.897) -- (-0.472,0.649) -- (-0.180,0.744);
\draw[pedge] (0.000,0.496) -- (-0.180,0.744) -- (-0.472,0.649);
\draw[pedge] (-0.292,0.897) -- (-0.472,0.649) -- (-0.472,0.956);
\draw[pedge] (-0.763,1.051) -- (-0.472,0.956) -- (-0.472,0.649);
\draw[pedge] (-0.292,0.897) -- (-0.292,1.204) -- (0.000,1.299);
\draw[pedge] (-0.763,1.051) -- (-0.472,0.956) -- (-0.292,1.204);
\draw[pedge] (-0.292,0.897) -- (-0.292,1.204) -- (-0.472,0.956);
\draw[pedge] (-1.235,2.503) -- (-1.235,2.196) -- (-1.527,2.101);
\draw[pedge] (-0.763,2.349) -- (-1.055,2.444) -- (-1.235,2.196);
\draw[pedge] (-1.235,2.503) -- (-1.235,2.196) -- (-1.055,2.444);
\draw[pedge] (-1.707,2.349) -- (-1.707,2.656) -- (-1.998,2.751);
\draw[pedge] (-1.235,2.503) -- (-1.527,2.408) -- (-1.707,2.656);
\draw[pedge] (-1.707,2.349) -- (-1.707,2.656) -- (-1.527,2.408);
\draw[pedge] (-1.707,2.349) -- (-1.527,2.101) -- (-1.527,2.408);
\draw[pedge] (-1.235,2.503) -- (-1.527,2.408) -- (-1.527,2.101);
\draw[pedge] (-1.055,1.452) -- (-0.763,1.547) -- (-0.763,1.853);
\draw[pedge] (-0.763,1.051) -- (-0.944,1.299) -- (-0.763,1.547);
\draw[pedge] (-1.055,1.452) -- (-0.763,1.547) -- (-0.944,1.299);
\draw[pedge] (-1.055,1.948) -- (-1.346,1.853) -- (-1.527,2.101);
\draw[pedge] (-1.055,1.452) -- (-1.055,1.759) -- (-1.346,1.853);
\draw[pedge] (-1.055,1.948) -- (-1.346,1.853) -- (-1.055,1.759);
\draw[pedge] (-1.055,1.948) -- (-0.763,1.853) -- (-1.055,1.759);
\draw[pedge] (-1.055,1.452) -- (-1.055,1.759) -- (-0.763,1.853);
\draw[pedge] (-1.055,1.948) -- (-0.763,1.853) -- (-0.944,2.101);
\draw[pedge] (-0.763,2.349) -- (-0.944,2.101) -- (-0.763,1.853);
\draw[pedge] (-1.055,1.948) -- (-1.235,2.196) -- (-1.527,2.101);
\draw[pedge] (-0.763,2.349) -- (-0.944,2.101) -- (-1.235,2.196);
\draw[pedge] (-1.055,1.948) -- (-1.235,2.196) -- (-0.944,2.101);
\draw[pedge] (-0.472,1.948) -- (-0.763,1.853) -- (-0.583,2.101);
\draw[pedge] (-0.763,2.349) -- (-0.583,2.101) -- (-0.763,1.853);
\draw[pedge] (-0.472,1.948) -- (-0.292,2.196) -- (0.000,2.101);
\draw[pedge] (-0.763,2.349) -- (-0.583,2.101) -- (-0.292,2.196);
\draw[pedge] (-0.472,1.948) -- (-0.292,2.196) -- (-0.583,2.101);
\draw[pedge] (-0.472,1.452) -- (-0.763,1.547) -- (-0.763,1.853);
\draw[pedge] (-0.763,1.051) -- (-0.583,1.299) -- (-0.763,1.547);
\draw[pedge] (-0.472,1.452) -- (-0.763,1.547) -- (-0.583,1.299);
\draw[pedge] (-0.472,1.948) -- (-0.180,1.853) -- (0.000,2.101);
\draw[pedge] (-0.472,1.452) -- (-0.472,1.759) -- (-0.180,1.853);
\draw[pedge] (-0.472,1.948) -- (-0.180,1.853) -- (-0.472,1.759);
\draw[pedge] (-0.472,1.948) -- (-0.763,1.853) -- (-0.472,1.759);
\draw[pedge] (-0.472,1.452) -- (-0.472,1.759) -- (-0.763,1.853);
\draw[pedge] (-2.290,1.051) -- (-2.470,1.299) -- (-2.179,1.204);
\draw[pedge] (-1.998,1.452) -- (-2.179,1.204) -- (-2.470,1.299);
\draw[pedge] (-2.290,1.051) -- (-1.998,0.956) -- (-1.998,0.649);
\draw[pedge] (-1.998,1.452) -- (-2.179,1.204) -- (-1.998,0.956);
\draw[pedge] (-2.290,1.051) -- (-1.998,0.956) -- (-2.179,1.204);
\draw[pedge] (-2.762,0.897) -- (-2.762,1.204) -- (-2.470,1.299);
\draw[pedge] (-3.234,1.051) -- (-2.942,0.956) -- (-2.762,1.204);
\draw[pedge] (-2.762,0.897) -- (-2.762,1.204) -- (-2.942,0.956);
\draw[pedge] (-2.290,1.051) -- (-2.290,0.744) -- (-1.998,0.649);
\draw[pedge] (-2.762,0.897) -- (-2.470,0.992) -- (-2.290,0.744);
\draw[pedge] (-2.290,1.051) -- (-2.290,0.744) -- (-2.470,0.992);
\draw[pedge] (-2.290,1.051) -- (-2.470,1.299) -- (-2.470,0.992);
\draw[pedge] (-2.762,0.897) -- (-2.470,0.992) -- (-2.470,1.299);
\draw[pedge] (-1.998,1.948) -- (-2.290,1.853) -- (-2.470,2.101);
\draw[pedge] (-1.998,1.452) -- (-1.998,1.759) -- (-2.290,1.853);
\draw[pedge] (-1.998,1.948) -- (-2.290,1.853) -- (-1.998,1.759);
\draw[pedge] (-2.290,2.349) -- (-1.998,2.444) -- (-1.998,2.751);
\draw[pedge] (-1.998,1.948) -- (-2.179,2.196) -- (-1.998,2.444);
\draw[pedge] (-2.290,2.349) -- (-1.998,2.444) -- (-2.179,2.196);
\draw[pedge] (-2.290,2.349) -- (-2.470,2.101) -- (-2.179,2.196);
\draw[pedge] (-1.998,1.948) -- (-2.179,2.196) -- (-2.470,2.101);
\draw[pedge] (-2.942,1.452) -- (-2.762,1.204) -- (-2.470,1.299);
\draw[pedge] (-3.234,1.051) -- (-3.053,1.299) -- (-2.762,1.204);
\draw[pedge] (-2.942,1.452) -- (-2.762,1.204) -- (-3.053,1.299);
\draw[pedge] (-2.470,1.605) -- (-2.650,1.853) -- (-2.470,2.101);
\draw[pedge] (-2.942,1.452) -- (-2.650,1.547) -- (-2.650,1.853);
\draw[pedge] (-2.470,1.605) -- (-2.650,1.853) -- (-2.650,1.547);
\draw[pedge] (-2.470,1.605) -- (-2.470,1.299) -- (-2.650,1.547);
\draw[pedge] (-2.942,1.452) -- (-2.650,1.547) -- (-2.470,1.299);
\draw[pedge] (-2.470,1.605) -- (-2.470,1.299) -- (-2.290,1.547);
\draw[pedge] (-1.998,1.452) -- (-2.290,1.547) -- (-2.470,1.299);
\draw[pedge] (-2.470,1.605) -- (-2.290,1.853) -- (-2.470,2.101);
\draw[pedge] (-1.998,1.452) -- (-2.290,1.547) -- (-2.290,1.853);
\draw[pedge] (-2.470,1.605) -- (-2.290,1.853) -- (-2.290,1.547);
\draw[pedge] (-1.235,0.897) -- (-1.055,0.649) -- (-1.235,0.401);
\draw[pedge] (-0.763,1.051) -- (-1.055,0.956) -- (-1.055,0.649);
\draw[pedge] (-1.235,0.897) -- (-1.055,0.649) -- (-1.055,0.956);
\draw[pedge] (-1.527,0.496) -- (-1.707,0.744) -- (-1.998,0.649);
\draw[pedge] (-1.235,0.897) -- (-1.415,0.649) -- (-1.707,0.744);
\draw[pedge] (-1.527,0.496) -- (-1.707,0.744) -- (-1.415,0.649);
\draw[pedge] (-1.527,0.496) -- (-1.235,0.401) -- (-1.415,0.649);
\draw[pedge] (-1.235,0.897) -- (-1.415,0.649) -- (-1.235,0.401);
\draw[pedge] (-0.472,0.153) -- (-0.292,0.401) -- (-0.472,0.649);
\draw[pedge] (0.000,0.000) -- (-0.292,0.095) -- (-0.292,0.401);
\draw[pedge] (-0.472,0.153) -- (-0.292,0.401) -- (-0.292,0.095);
\draw[pedge] (-0.763,0.555) -- (-0.944,0.307) -- (-1.235,0.401);
\draw[pedge] (-0.472,0.153) -- (-0.652,0.401) -- (-0.944,0.307);
\draw[pedge] (-0.763,0.555) -- (-0.944,0.307) -- (-0.652,0.401);
\draw[pedge] (-0.763,0.555) -- (-0.472,0.649) -- (-0.652,0.401);
\draw[pedge] (-0.472,0.153) -- (-0.652,0.401) -- (-0.472,0.649);
\draw[pedge] (-0.763,0.555) -- (-0.472,0.649) -- (-0.763,0.744);
\draw[pedge] (-0.763,1.051) -- (-0.763,0.744) -- (-0.472,0.649);
\draw[pedge] (-0.763,0.555) -- (-1.055,0.649) -- (-1.235,0.401);
\draw[pedge] (-0.763,1.051) -- (-0.763,0.744) -- (-1.055,0.649);
\draw[pedge] (-0.763,0.555) -- (-1.055,0.649) -- (-0.763,0.744);
\draw[pedge] (-1.998,1.948) -- (-1.707,1.853) -- (-1.527,2.101);
\draw[pedge] (-1.998,1.452) -- (-1.998,1.759) -- (-1.707,1.853);
\draw[pedge] (-1.998,1.948) -- (-1.707,1.853) -- (-1.998,1.759);
\draw[pedge] (-1.707,2.349) -- (-1.998,2.444) -- (-1.998,2.751);
\draw[pedge] (-1.998,1.948) -- (-1.818,2.196) -- (-1.998,2.444);
\draw[pedge] (-1.707,2.349) -- (-1.998,2.444) -- (-1.818,2.196);
\draw[pedge] (-1.707,2.349) -- (-1.527,2.101) -- (-1.818,2.196);
\draw[pedge] (-1.998,1.948) -- (-1.818,2.196) -- (-1.527,2.101);
\draw[pedge] (-1.055,1.452) -- (-1.235,1.204) -- (-1.527,1.299);
\draw[pedge] (-0.763,1.051) -- (-0.944,1.299) -- (-1.235,1.204);
\draw[pedge] (-1.055,1.452) -- (-1.235,1.204) -- (-0.944,1.299);
\draw[pedge] (-1.527,1.605) -- (-1.346,1.853) -- (-1.527,2.101);
\draw[pedge] (-1.055,1.452) -- (-1.346,1.547) -- (-1.346,1.853);
\draw[pedge] (-1.527,1.605) -- (-1.346,1.853) -- (-1.346,1.547);
\draw[pedge] (-1.527,1.605) -- (-1.527,1.299) -- (-1.346,1.547);
\draw[pedge] (-1.055,1.452) -- (-1.346,1.547) -- (-1.527,1.299);
\draw[pedge] (-1.527,1.605) -- (-1.527,1.299) -- (-1.707,1.547);
\draw[pedge] (-1.998,1.452) -- (-1.707,1.547) -- (-1.527,1.299);
\draw[pedge] (-1.527,1.605) -- (-1.707,1.853) -- (-1.527,2.101);
\draw[pedge] (-1.998,1.452) -- (-1.707,1.547) -- (-1.707,1.853);
\draw[pedge] (-1.527,1.605) -- (-1.707,1.853) -- (-1.707,1.547);
\draw[pedge] (-1.707,1.051) -- (-1.527,1.299) -- (-1.818,1.204);
\draw[pedge] (-1.998,1.452) -- (-1.818,1.204) -- (-1.527,1.299);
\draw[pedge] (-1.707,1.051) -- (-1.998,0.956) -- (-1.998,0.649);
\draw[pedge] (-1.998,1.452) -- (-1.818,1.204) -- (-1.998,0.956);
\draw[pedge] (-1.707,1.051) -- (-1.998,0.956) -- (-1.818,1.204);
\draw[pedge] (-1.235,0.897) -- (-1.235,1.204) -- (-1.527,1.299);
\draw[pedge] (-0.763,1.051) -- (-1.055,0.956) -- (-1.235,1.204);
\draw[pedge] (-1.235,0.897) -- (-1.235,1.204) -- (-1.055,0.956);
\draw[pedge] (-1.707,1.051) -- (-1.707,0.744) -- (-1.998,0.649);
\draw[pedge] (-1.235,0.897) -- (-1.527,0.992) -- (-1.707,0.744);
\draw[pedge] (-1.707,1.051) -- (-1.707,0.744) -- (-1.527,0.992);
\draw[pedge] (-1.707,1.051) -- (-1.527,1.299) -- (-1.527,0.992);
\draw[pedge] (-1.235,0.897) -- (-1.527,0.992) -- (-1.527,1.299);
\draw[pedge] (-2.470,0.496) -- (-2.762,0.401) -- (-2.470,0.307);
\draw[pedge] (-2.470,0.000) -- (-2.470,0.307) -- (-2.762,0.401);
\draw[pedge] (-2.470,0.496) -- (-2.179,0.401) -- (-1.998,0.649);
\draw[pedge] (-2.470,0.000) -- (-2.470,0.307) -- (-2.179,0.401);
\draw[pedge] (-2.470,0.496) -- (-2.179,0.401) -- (-2.470,0.307);
\draw[pedge] (-2.762,0.897) -- (-2.942,0.649) -- (-2.762,0.401);
\draw[pedge] (-3.234,1.051) -- (-2.942,0.956) -- (-2.942,0.649);
\draw[pedge] (-2.762,0.897) -- (-2.942,0.649) -- (-2.942,0.956);
\draw[pedge] (-2.470,0.496) -- (-2.290,0.744) -- (-1.998,0.649);
\draw[pedge] (-2.762,0.897) -- (-2.582,0.649) -- (-2.290,0.744);
\draw[pedge] (-2.470,0.496) -- (-2.290,0.744) -- (-2.582,0.649);
\draw[pedge] (-2.470,0.496) -- (-2.762,0.401) -- (-2.582,0.649);
\draw[pedge] (-2.762,0.897) -- (-2.582,0.649) -- (-2.762,0.401);
\draw[pedge] (-2.762,-0.401) -- (-2.942,-0.153) -- (-3.234,-0.248);
\draw[pedge] (-2.470,0.000) -- (-2.650,-0.248) -- (-2.942,-0.153);
\draw[pedge] (-2.762,-0.401) -- (-2.942,-0.153) -- (-2.650,-0.248);
\draw[pedge] (-3.234,-0.555) -- (-3.053,-0.803) -- (-3.234,-1.051);
\draw[pedge] (-2.762,-0.401) -- (-3.053,-0.496) -- (-3.053,-0.803);
\draw[pedge] (-3.234,-0.555) -- (-3.053,-0.803) -- (-3.053,-0.496);
\draw[pedge] (-3.234,-0.555) -- (-3.234,-0.248) -- (-3.053,-0.496);
\draw[pedge] (-2.762,-0.401) -- (-3.053,-0.496) -- (-3.234,-0.248);
\draw[pedge] (-3.234,0.555) -- (-2.942,0.649) -- (-2.762,0.401);
\draw[pedge] (-3.234,1.051) -- (-3.234,0.744) -- (-2.942,0.649);
\draw[pedge] (-3.234,0.555) -- (-2.942,0.649) -- (-3.234,0.744);
\draw[pedge] (-2.942,0.153) -- (-3.234,0.059) -- (-3.234,-0.248);
\draw[pedge] (-3.234,0.555) -- (-3.053,0.307) -- (-3.234,0.059);
\draw[pedge] (-2.942,0.153) -- (-3.234,0.059) -- (-3.053,0.307);
\draw[pedge] (-2.942,0.153) -- (-2.762,0.401) -- (-3.053,0.307);
\draw[pedge] (-3.234,0.555) -- (-3.053,0.307) -- (-2.762,0.401);
\draw[pedge] (-2.942,0.153) -- (-2.762,0.401) -- (-2.762,0.095);
\draw[pedge] (-2.470,0.000) -- (-2.762,0.095) -- (-2.762,0.401);
\draw[pedge] (-2.942,0.153) -- (-2.942,-0.153) -- (-3.234,-0.248);
\draw[pedge] (-2.470,0.000) -- (-2.762,0.095) -- (-2.942,-0.153);
\draw[pedge] (-2.942,0.153) -- (-2.942,-0.153) -- (-2.762,0.095);
\draw[pedge] (-1.527,0.000) -- (-1.235,0.095) -- (-1.235,0.401);
\draw[pedge] (-1.235,-0.401) -- (-1.415,-0.153) -- (-1.235,0.095);
\draw[pedge] (-1.527,0.000) -- (-1.235,0.095) -- (-1.415,-0.153);
\draw[pedge] (-1.527,0.496) -- (-1.818,0.401) -- (-1.998,0.649);
\draw[pedge] (-1.527,0.000) -- (-1.527,0.307) -- (-1.818,0.401);
\draw[pedge] (-1.527,0.496) -- (-1.818,0.401) -- (-1.527,0.307);
\draw[pedge] (-1.527,0.496) -- (-1.235,0.401) -- (-1.527,0.307);
\draw[pedge] (-1.527,0.000) -- (-1.527,0.307) -- (-1.235,0.401);
\draw[pedge] (-0.472,0.153) -- (-0.472,-0.153) -- (-0.763,-0.248);
\draw[pedge] (0.000,0.000) -- (-0.292,0.095) -- (-0.472,-0.153);
\draw[pedge] (-0.472,0.153) -- (-0.472,-0.153) -- (-0.292,0.095);
\draw[pedge] (-0.944,0.000) -- (-0.944,0.307) -- (-1.235,0.401);
\draw[pedge] (-0.472,0.153) -- (-0.763,0.059) -- (-0.944,0.307);
\draw[pedge] (-0.944,0.000) -- (-0.944,0.307) -- (-0.763,0.059);
\draw[pedge] (-0.944,0.000) -- (-0.763,-0.248) -- (-0.763,0.059);
\draw[pedge] (-0.472,0.153) -- (-0.763,0.059) -- (-0.763,-0.248);
\draw[pedge] (-0.944,0.000) -- (-0.763,-0.248) -- (-1.055,-0.153);
\draw[pedge] (-1.235,-0.401) -- (-1.055,-0.153) -- (-0.763,-0.248);
\draw[pedge] (-0.944,0.000) -- (-1.235,0.095) -- (-1.235,0.401);
\draw[pedge] (-1.235,-0.401) -- (-1.055,-0.153) -- (-1.235,0.095);
\draw[pedge] (-0.944,0.000) -- (-1.235,0.095) -- (-1.055,-0.153);
\draw[pedge] (-2.762,-0.401) -- (-2.470,-0.496) -- (-2.470,-0.803);
\draw[pedge] (-2.470,0.000) -- (-2.650,-0.248) -- (-2.470,-0.496);
\draw[pedge] (-2.762,-0.401) -- (-2.470,-0.496) -- (-2.650,-0.248);
\draw[pedge] (-2.762,-0.897) -- (-3.053,-0.803) -- (-3.234,-1.051);
\draw[pedge] (-2.762,-0.401) -- (-2.762,-0.708) -- (-3.053,-0.803);
\draw[pedge] (-2.762,-0.897) -- (-3.053,-0.803) -- (-2.762,-0.708);
\draw[pedge] (-2.762,-0.897) -- (-2.470,-0.803) -- (-2.762,-0.708);
\draw[pedge] (-2.762,-0.401) -- (-2.762,-0.708) -- (-2.470,-0.803);
\draw[pedge] (-1.707,-0.555) -- (-1.707,-0.248) -- (-1.998,-0.153);
\draw[pedge] (-1.235,-0.401) -- (-1.527,-0.496) -- (-1.707,-0.248);
\draw[pedge] (-1.707,-0.555) -- (-1.707,-0.248) -- (-1.527,-0.496);
\draw[pedge] (-2.179,-0.401) -- (-2.179,-0.708) -- (-2.470,-0.803);
\draw[pedge] (-1.707,-0.555) -- (-1.998,-0.460) -- (-2.179,-0.708);
\draw[pedge] (-2.179,-0.401) -- (-2.179,-0.708) -- (-1.998,-0.460);
\draw[pedge] (-2.179,-0.401) -- (-1.998,-0.153) -- (-1.998,-0.460);
\draw[pedge] (-1.707,-0.555) -- (-1.998,-0.460) -- (-1.998,-0.153);
\draw[pedge] (-2.179,-0.401) -- (-1.998,-0.153) -- (-2.290,-0.248);
\draw[pedge] (-2.470,0.000) -- (-2.290,-0.248) -- (-1.998,-0.153);
\draw[pedge] (-2.179,-0.401) -- (-2.470,-0.496) -- (-2.470,-0.803);
\draw[pedge] (-2.470,0.000) -- (-2.290,-0.248) -- (-2.470,-0.496);
\draw[pedge] (-2.179,-0.401) -- (-2.470,-0.496) -- (-2.290,-0.248);
\draw[pedge] (-1.998,0.153) -- (-1.998,-0.153) -- (-2.179,0.095);
\draw[pedge] (-2.470,0.000) -- (-2.179,0.095) -- (-1.998,-0.153);
\draw[pedge] (-1.998,0.153) -- (-2.179,0.401) -- (-1.998,0.649);
\draw[pedge] (-2.470,0.000) -- (-2.179,0.095) -- (-2.179,0.401);
\draw[pedge] (-1.998,0.153) -- (-2.179,0.401) -- (-2.179,0.095);
\draw[pedge] (-1.527,0.000) -- (-1.707,-0.248) -- (-1.998,-0.153);
\draw[pedge] (-1.235,-0.401) -- (-1.415,-0.153) -- (-1.707,-0.248);
\draw[pedge] (-1.527,0.000) -- (-1.707,-0.248) -- (-1.415,-0.153);
\draw[pedge] (-1.998,0.153) -- (-1.818,0.401) -- (-1.998,0.649);
\draw[pedge] (-1.527,0.000) -- (-1.818,0.095) -- (-1.818,0.401);
\draw[pedge] (-1.998,0.153) -- (-1.818,0.401) -- (-1.818,0.095);
\draw[pedge] (-1.998,0.153) -- (-1.998,-0.153) -- (-1.818,0.095);
\draw[pedge] (-1.527,0.000) -- (-1.818,0.095) -- (-1.998,-0.153);
\draw[pedge] (-1.707,-1.853) -- (-1.998,-1.948) -- (-1.818,-1.700);
\draw[pedge] (-1.998,-1.452) -- (-1.818,-1.700) -- (-1.998,-1.948);
\draw[pedge] (-1.707,-1.853) -- (-1.527,-1.605) -- (-1.235,-1.700);
\draw[pedge] (-1.998,-1.452) -- (-1.818,-1.700) -- (-1.527,-1.605);
\draw[pedge] (-1.707,-1.853) -- (-1.527,-1.605) -- (-1.818,-1.700);
\draw[pedge] (-1.707,-2.349) -- (-1.998,-2.255) -- (-1.998,-1.948);
\draw[pedge] (-1.998,-2.751) -- (-1.818,-2.503) -- (-1.998,-2.255);
\draw[pedge] (-1.707,-2.349) -- (-1.998,-2.255) -- (-1.818,-2.503);
\draw[pedge] (-1.707,-1.853) -- (-1.415,-1.948) -- (-1.235,-1.700);
\draw[pedge] (-1.707,-2.349) -- (-1.707,-2.043) -- (-1.415,-1.948);
\draw[pedge] (-1.707,-1.853) -- (-1.415,-1.948) -- (-1.707,-2.043);
\draw[pedge] (-1.707,-1.853) -- (-1.998,-1.948) -- (-1.707,-2.043);
\draw[pedge] (-1.707,-2.349) -- (-1.707,-2.043) -- (-1.998,-1.948);
\draw[pedge] (-2.470,-1.299) -- (-2.470,-1.605) -- (-2.762,-1.700);
\draw[pedge] (-1.998,-1.452) -- (-2.290,-1.357) -- (-2.470,-1.605);
\draw[pedge] (-2.470,-1.299) -- (-2.470,-1.605) -- (-2.290,-1.357);
\draw[pedge] (-2.942,-1.452) -- (-2.942,-1.145) -- (-3.234,-1.051);
\draw[pedge] (-2.470,-1.299) -- (-2.762,-1.393) -- (-2.942,-1.145);
\draw[pedge] (-2.942,-1.452) -- (-2.942,-1.145) -- (-2.762,-1.393);
\draw[pedge] (-2.942,-1.452) -- (-2.762,-1.700) -- (-2.762,-1.393);
\draw[pedge] (-2.470,-1.299) -- (-2.762,-1.393) -- (-2.762,-1.700);
\draw[pedge] (-2.290,-2.349) -- (-1.998,-2.255) -- (-1.998,-1.948);
\draw[pedge] (-1.998,-2.751) -- (-2.179,-2.503) -- (-1.998,-2.255);
\draw[pedge] (-2.290,-2.349) -- (-1.998,-2.255) -- (-2.179,-2.503);
\draw[pedge] (-2.290,-1.853) -- (-2.582,-1.948) -- (-2.762,-1.700);
\draw[pedge] (-2.290,-2.349) -- (-2.290,-2.043) -- (-2.582,-1.948);
\draw[pedge] (-2.290,-1.853) -- (-2.582,-1.948) -- (-2.290,-2.043);
\draw[pedge] (-2.290,-1.853) -- (-1.998,-1.948) -- (-2.290,-2.043);
\draw[pedge] (-2.290,-2.349) -- (-2.290,-2.043) -- (-1.998,-1.948);
\draw[pedge] (-2.290,-1.853) -- (-1.998,-1.948) -- (-2.179,-1.700);
\draw[pedge] (-1.998,-1.452) -- (-2.179,-1.700) -- (-1.998,-1.948);
\draw[pedge] (-2.290,-1.853) -- (-2.470,-1.605) -- (-2.762,-1.700);
\draw[pedge] (-1.998,-1.452) -- (-2.179,-1.700) -- (-2.470,-1.605);
\draw[pedge] (-2.290,-1.853) -- (-2.470,-1.605) -- (-2.179,-1.700);
\draw[pedge] (-1.235,-0.897) -- (-0.944,-0.803) -- (-0.763,-1.051);
\draw[pedge] (-1.235,-0.401) -- (-1.235,-0.708) -- (-0.944,-0.803);
\draw[pedge] (-1.235,-0.897) -- (-0.944,-0.803) -- (-1.235,-0.708);
\draw[pedge] (-0.944,-1.299) -- (-1.235,-1.393) -- (-1.235,-1.700);
\draw[pedge] (-1.235,-0.897) -- (-1.055,-1.145) -- (-1.235,-1.393);
\draw[pedge] (-0.944,-1.299) -- (-1.235,-1.393) -- (-1.055,-1.145);
\draw[pedge] (-0.944,-1.299) -- (-0.763,-1.051) -- (-1.055,-1.145);
\draw[pedge] (-1.235,-0.897) -- (-1.055,-1.145) -- (-0.763,-1.051);
\draw[pedge] (-0.292,-0.401) -- (-0.472,-0.153) -- (-0.763,-0.248);
\draw[pedge] (0.000,0.000) -- (-0.180,-0.248) -- (-0.472,-0.153);
\draw[pedge] (-0.292,-0.401) -- (-0.472,-0.153) -- (-0.180,-0.248);
\draw[pedge] (-0.763,-0.555) -- (-0.583,-0.803) -- (-0.763,-1.051);
\draw[pedge] (-0.292,-0.401) -- (-0.583,-0.496) -- (-0.583,-0.803);
\draw[pedge] (-0.763,-0.555) -- (-0.583,-0.803) -- (-0.583,-0.496);
\draw[pedge] (-0.763,-0.555) -- (-0.763,-0.248) -- (-0.583,-0.496);
\draw[pedge] (-0.292,-0.401) -- (-0.583,-0.496) -- (-0.763,-0.248);
\draw[pedge] (-0.763,-0.555) -- (-0.763,-0.248) -- (-0.944,-0.496);
\draw[pedge] (-1.235,-0.401) -- (-0.944,-0.496) -- (-0.763,-0.248);
\draw[pedge] (-0.763,-0.555) -- (-0.944,-0.803) -- (-0.763,-1.051);
\draw[pedge] (-1.235,-0.401) -- (-0.944,-0.496) -- (-0.944,-0.803);
\draw[pedge] (-0.763,-0.555) -- (-0.944,-0.803) -- (-0.944,-0.496);
\draw[pedge] (-2.470,-1.299) -- (-2.290,-1.051) -- (-2.470,-0.803);
\draw[pedge] (-1.998,-1.452) -- (-2.290,-1.357) -- (-2.290,-1.051);
\draw[pedge] (-2.470,-1.299) -- (-2.290,-1.051) -- (-2.290,-1.357);
\draw[pedge] (-2.762,-0.897) -- (-2.942,-1.145) -- (-3.234,-1.051);
\draw[pedge] (-2.470,-1.299) -- (-2.650,-1.051) -- (-2.942,-1.145);
\draw[pedge] (-2.762,-0.897) -- (-2.942,-1.145) -- (-2.650,-1.051);
\draw[pedge] (-2.762,-0.897) -- (-2.470,-0.803) -- (-2.650,-1.051);
\draw[pedge] (-2.470,-1.299) -- (-2.650,-1.051) -- (-2.470,-0.803);
\draw[pedge] (-1.707,-0.555) -- (-1.527,-0.803) -- (-1.707,-1.051);
\draw[pedge] (-1.235,-0.401) -- (-1.527,-0.496) -- (-1.527,-0.803);
\draw[pedge] (-1.707,-0.555) -- (-1.527,-0.803) -- (-1.527,-0.496);
\draw[pedge] (-1.998,-0.956) -- (-2.179,-0.708) -- (-2.470,-0.803);
\draw[pedge] (-1.707,-0.555) -- (-1.887,-0.803) -- (-2.179,-0.708);
\draw[pedge] (-1.998,-0.956) -- (-2.179,-0.708) -- (-1.887,-0.803);
\draw[pedge] (-1.998,-0.956) -- (-1.707,-1.051) -- (-1.887,-0.803);
\draw[pedge] (-1.707,-0.555) -- (-1.887,-0.803) -- (-1.707,-1.051);
\draw[pedge] (-1.998,-0.956) -- (-1.707,-1.051) -- (-1.998,-1.145);
\draw[pedge] (-1.998,-1.452) -- (-1.998,-1.145) -- (-1.707,-1.051);
\draw[pedge] (-1.998,-0.956) -- (-2.290,-1.051) -- (-2.470,-0.803);
\draw[pedge] (-1.998,-1.452) -- (-1.998,-1.145) -- (-2.290,-1.051);
\draw[pedge] (-1.998,-0.956) -- (-2.290,-1.051) -- (-1.998,-1.145);
\draw[pedge] (-1.527,-1.299) -- (-1.707,-1.051) -- (-1.707,-1.357);
\draw[pedge] (-1.998,-1.452) -- (-1.707,-1.357) -- (-1.707,-1.051);
\draw[pedge] (-1.527,-1.299) -- (-1.527,-1.605) -- (-1.235,-1.700);
\draw[pedge] (-1.998,-1.452) -- (-1.707,-1.357) -- (-1.527,-1.605);
\draw[pedge] (-1.527,-1.299) -- (-1.527,-1.605) -- (-1.707,-1.357);
\draw[pedge] (-1.235,-0.897) -- (-1.527,-0.803) -- (-1.707,-1.051);
\draw[pedge] (-1.235,-0.401) -- (-1.235,-0.708) -- (-1.527,-0.803);
\draw[pedge] (-1.235,-0.897) -- (-1.527,-0.803) -- (-1.235,-0.708);
\draw[pedge] (-1.527,-1.299) -- (-1.235,-1.393) -- (-1.235,-1.700);
\draw[pedge] (-1.235,-0.897) -- (-1.415,-1.145) -- (-1.235,-1.393);
\draw[pedge] (-1.527,-1.299) -- (-1.235,-1.393) -- (-1.415,-1.145);
\draw[pedge] (-1.527,-1.299) -- (-1.707,-1.051) -- (-1.415,-1.145);
\draw[pedge] (-1.235,-0.897) -- (-1.415,-1.145) -- (-1.707,-1.051);
\draw[pedge] (-1.235,-2.196) -- (-1.235,-2.503) -- (-1.055,-2.255);
\draw[pedge] (-0.763,-2.349) -- (-1.055,-2.255) -- (-1.235,-2.503);
\draw[pedge] (-1.235,-2.196) -- (-1.055,-1.948) -- (-1.235,-1.700);
\draw[pedge] (-0.763,-2.349) -- (-1.055,-2.255) -- (-1.055,-1.948);
\draw[pedge] (-1.235,-2.196) -- (-1.055,-1.948) -- (-1.055,-2.255);
\draw[pedge] (-1.707,-2.349) -- (-1.527,-2.597) -- (-1.235,-2.503);
\draw[pedge] (-1.998,-2.751) -- (-1.818,-2.503) -- (-1.527,-2.597);
\draw[pedge] (-1.707,-2.349) -- (-1.527,-2.597) -- (-1.818,-2.503);
\draw[pedge] (-1.235,-2.196) -- (-1.415,-1.948) -- (-1.235,-1.700);
\draw[pedge] (-1.707,-2.349) -- (-1.415,-2.255) -- (-1.415,-1.948);
\draw[pedge] (-1.235,-2.196) -- (-1.415,-1.948) -- (-1.415,-2.255);
\draw[pedge] (-1.235,-2.196) -- (-1.235,-2.503) -- (-1.415,-2.255);
\draw[pedge] (-1.707,-2.349) -- (-1.415,-2.255) -- (-1.235,-2.503);
\draw[pedge] (-0.472,-2.751) -- (-0.763,-2.845) -- (-0.763,-3.152);
\draw[pedge] (-0.763,-2.349) -- (-0.583,-2.597) -- (-0.763,-2.845);
\draw[pedge] (-0.472,-2.751) -- (-0.763,-2.845) -- (-0.583,-2.597);
\draw[pedge] (-0.472,-3.247) -- (-0.180,-3.152) -- (-0.000,-3.400);
\draw[pedge] (-0.472,-2.751) -- (-0.472,-3.057) -- (-0.180,-3.152);
\draw[pedge] (-0.472,-3.247) -- (-0.180,-3.152) -- (-0.472,-3.057);
\draw[pedge] (-0.472,-3.247) -- (-0.763,-3.152) -- (-0.472,-3.057);
\draw[pedge] (-0.472,-2.751) -- (-0.472,-3.057) -- (-0.763,-3.152);
\draw[pedge] (-1.527,-2.904) -- (-1.527,-2.597) -- (-1.235,-2.503);
\draw[pedge] (-1.998,-2.751) -- (-1.707,-2.845) -- (-1.527,-2.597);
\draw[pedge] (-1.527,-2.904) -- (-1.527,-2.597) -- (-1.707,-2.845);
\draw[pedge] (-1.055,-2.751) -- (-1.055,-3.057) -- (-0.763,-3.152);
\draw[pedge] (-1.527,-2.904) -- (-1.235,-2.809) -- (-1.055,-3.057);
\draw[pedge] (-1.055,-2.751) -- (-1.055,-3.057) -- (-1.235,-2.809);
\draw[pedge] (-1.055,-2.751) -- (-1.235,-2.503) -- (-1.235,-2.809);
\draw[pedge] (-1.527,-2.904) -- (-1.235,-2.809) -- (-1.235,-2.503);
\draw[pedge] (-1.055,-2.751) -- (-1.235,-2.503) -- (-0.944,-2.597);
\draw[pedge] (-0.763,-2.349) -- (-0.944,-2.597) -- (-1.235,-2.503);
\draw[pedge] (-1.055,-2.751) -- (-0.763,-2.845) -- (-0.763,-3.152);
\draw[pedge] (-0.763,-2.349) -- (-0.944,-2.597) -- (-0.763,-2.845);
\draw[pedge] (-1.055,-2.751) -- (-0.763,-2.845) -- (-0.944,-2.597);
\draw[pedge] (-0.472,-1.452) -- (-0.472,-1.145) -- (-0.763,-1.051);
\draw[pedge] (-0.000,-1.299) -- (-0.292,-1.393) -- (-0.472,-1.145);
\draw[pedge] (-0.472,-1.452) -- (-0.472,-1.145) -- (-0.292,-1.393);
\draw[pedge] (-0.944,-1.299) -- (-0.944,-1.605) -- (-1.235,-1.700);
\draw[pedge] (-0.472,-1.452) -- (-0.763,-1.357) -- (-0.944,-1.605);
\draw[pedge] (-0.944,-1.299) -- (-0.944,-1.605) -- (-0.763,-1.357);
\draw[pedge] (-0.944,-1.299) -- (-0.763,-1.051) -- (-0.763,-1.357);
\draw[pedge] (-0.472,-1.452) -- (-0.763,-1.357) -- (-0.763,-1.051);
\draw[pedge] (-0.292,-0.401) -- (-0.000,-0.496) -- (-0.000,-0.803);
\draw[pedge] (0.000,0.000) -- (-0.180,-0.248) -- (-0.000,-0.496);
\draw[pedge] (-0.292,-0.401) -- (-0.000,-0.496) -- (-0.180,-0.248);
\draw[pedge] (-0.292,-0.897) -- (-0.583,-0.803) -- (-0.763,-1.051);
\draw[pedge] (-0.292,-0.401) -- (-0.292,-0.708) -- (-0.583,-0.803);
\draw[pedge] (-0.292,-0.897) -- (-0.583,-0.803) -- (-0.292,-0.708);
\draw[pedge] (-0.292,-0.897) -- (-0.000,-0.803) -- (-0.292,-0.708);
\draw[pedge] (-0.292,-0.401) -- (-0.292,-0.708) -- (-0.000,-0.803);
\draw[pedge] (-0.292,-0.897) -- (-0.000,-0.803) -- (-0.180,-1.051);
\draw[pedge] (-0.000,-1.299) -- (-0.180,-1.051) -- (-0.000,-0.803);
\draw[pedge] (-0.292,-0.897) -- (-0.472,-1.145) -- (-0.763,-1.051);
\draw[pedge] (-0.000,-1.299) -- (-0.180,-1.051) -- (-0.472,-1.145);
\draw[pedge] (-0.292,-0.897) -- (-0.472,-1.145) -- (-0.180,-1.051);
\draw[pedge] (-0.472,-2.751) -- (-0.292,-2.503) -- (-0.000,-2.597);
\draw[pedge] (-0.763,-2.349) -- (-0.583,-2.597) -- (-0.292,-2.503);
\draw[pedge] (-0.472,-2.751) -- (-0.292,-2.503) -- (-0.583,-2.597);
\draw[pedge] (-0.000,-2.904) -- (-0.180,-3.152) -- (-0.000,-3.400);
\draw[pedge] (-0.472,-2.751) -- (-0.180,-2.845) -- (-0.180,-3.152);
\draw[pedge] (-0.000,-2.904) -- (-0.180,-3.152) -- (-0.180,-2.845);
\draw[pedge] (-0.000,-2.904) -- (-0.000,-2.597) -- (-0.180,-2.845);
\draw[pedge] (-0.472,-2.751) -- (-0.180,-2.845) -- (-0.000,-2.597);
\draw[pedge] (-0.000,-1.795) -- (-0.292,-1.700) -- (-0.472,-1.948);
\draw[pedge] (-0.000,-1.299) -- (-0.000,-1.605) -- (-0.292,-1.700);
\draw[pedge] (-0.000,-1.795) -- (-0.292,-1.700) -- (-0.000,-1.605);
\draw[pedge] (-0.292,-2.196) -- (-0.000,-2.291) -- (-0.000,-2.597);
\draw[pedge] (-0.000,-1.795) -- (-0.180,-2.043) -- (-0.000,-2.291);
\draw[pedge] (-0.292,-2.196) -- (-0.000,-2.291) -- (-0.180,-2.043);
\draw[pedge] (-0.292,-2.196) -- (-0.472,-1.948) -- (-0.180,-2.043);
\draw[pedge] (-0.000,-1.795) -- (-0.180,-2.043) -- (-0.472,-1.948);
\draw[pedge] (-0.292,-2.196) -- (-0.472,-1.948) -- (-0.472,-2.255);
\draw[pedge] (-0.763,-2.349) -- (-0.472,-2.255) -- (-0.472,-1.948);
\draw[pedge] (-0.292,-2.196) -- (-0.292,-2.503) -- (-0.000,-2.597);
\draw[pedge] (-0.763,-2.349) -- (-0.472,-2.255) -- (-0.292,-2.503);
\draw[pedge] (-0.292,-2.196) -- (-0.292,-2.503) -- (-0.472,-2.255);
\draw[pedge] (-0.763,-1.853) -- (-0.472,-1.948) -- (-0.763,-2.043);
\draw[pedge] (-0.763,-2.349) -- (-0.763,-2.043) -- (-0.472,-1.948);
\draw[pedge] (-0.763,-1.853) -- (-1.055,-1.948) -- (-1.235,-1.700);
\draw[pedge] (-0.763,-2.349) -- (-0.763,-2.043) -- (-1.055,-1.948);
\draw[pedge] (-0.763,-1.853) -- (-1.055,-1.948) -- (-0.763,-2.043);
\draw[pedge] (-0.472,-1.452) -- (-0.292,-1.700) -- (-0.472,-1.948);
\draw[pedge] (-0.000,-1.299) -- (-0.292,-1.393) -- (-0.292,-1.700);
\draw[pedge] (-0.472,-1.452) -- (-0.292,-1.700) -- (-0.292,-1.393);
\draw[pedge] (-0.763,-1.853) -- (-0.944,-1.605) -- (-1.235,-1.700);
\draw[pedge] (-0.472,-1.452) -- (-0.652,-1.700) -- (-0.944,-1.605);
\draw[pedge] (-0.763,-1.853) -- (-0.944,-1.605) -- (-0.652,-1.700);
\draw[pedge] (-0.763,-1.853) -- (-0.472,-1.948) -- (-0.652,-1.700);
\draw[pedge] (-0.472,-1.452) -- (-0.652,-1.700) -- (-0.472,-1.948);
\draw[pedge] (1.235,-2.196) -- (1.235,-2.503) -- (1.055,-2.255);
\draw[pedge] (0.763,-2.349) -- (1.055,-2.255) -- (1.235,-2.503);
\draw[pedge] (1.235,-2.196) -- (1.055,-1.948) -- (1.235,-1.700);
\draw[pedge] (0.763,-2.349) -- (1.055,-2.255) -- (1.055,-1.948);
\draw[pedge] (1.235,-2.196) -- (1.055,-1.948) -- (1.055,-2.255);
\draw[pedge] (1.707,-2.349) -- (1.527,-2.597) -- (1.235,-2.503);
\draw[pedge] (1.998,-2.751) -- (1.818,-2.503) -- (1.527,-2.597);
\draw[pedge] (1.707,-2.349) -- (1.527,-2.597) -- (1.818,-2.503);
\draw[pedge] (1.235,-2.196) -- (1.415,-1.948) -- (1.235,-1.700);
\draw[pedge] (1.707,-2.349) -- (1.415,-2.255) -- (1.415,-1.948);
\draw[pedge] (1.235,-2.196) -- (1.415,-1.948) -- (1.415,-2.255);
\draw[pedge] (1.235,-2.196) -- (1.235,-2.503) -- (1.415,-2.255);
\draw[pedge] (1.707,-2.349) -- (1.415,-2.255) -- (1.235,-2.503);
\draw[pedge] (0.472,-2.751) -- (0.763,-2.845) -- (0.763,-3.152);
\draw[pedge] (0.763,-2.349) -- (0.583,-2.597) -- (0.763,-2.845);
\draw[pedge] (0.472,-2.751) -- (0.763,-2.845) -- (0.583,-2.597);
\draw[pedge] (0.472,-3.247) -- (0.180,-3.152) -- (-0.000,-3.400);
\draw[pedge] (0.472,-2.751) -- (0.472,-3.057) -- (0.180,-3.152);
\draw[pedge] (0.472,-3.247) -- (0.180,-3.152) -- (0.472,-3.057);
\draw[pedge] (0.472,-3.247) -- (0.763,-3.152) -- (0.472,-3.057);
\draw[pedge] (0.472,-2.751) -- (0.472,-3.057) -- (0.763,-3.152);
\draw[pedge] (1.527,-2.904) -- (1.527,-2.597) -- (1.235,-2.503);
\draw[pedge] (1.998,-2.751) -- (1.707,-2.845) -- (1.527,-2.597);
\draw[pedge] (1.527,-2.904) -- (1.527,-2.597) -- (1.707,-2.845);
\draw[pedge] (1.055,-2.751) -- (1.055,-3.057) -- (0.763,-3.152);
\draw[pedge] (1.527,-2.904) -- (1.235,-2.809) -- (1.055,-3.057);
\draw[pedge] (1.055,-2.751) -- (1.055,-3.057) -- (1.235,-2.809);
\draw[pedge] (1.055,-2.751) -- (1.235,-2.503) -- (1.235,-2.809);
\draw[pedge] (1.527,-2.904) -- (1.235,-2.809) -- (1.235,-2.503);
\draw[pedge] (1.055,-2.751) -- (1.235,-2.503) -- (0.944,-2.597);
\draw[pedge] (0.763,-2.349) -- (0.944,-2.597) -- (1.235,-2.503);
\draw[pedge] (1.055,-2.751) -- (0.763,-2.845) -- (0.763,-3.152);
\draw[pedge] (0.763,-2.349) -- (0.944,-2.597) -- (0.763,-2.845);
\draw[pedge] (1.055,-2.751) -- (0.763,-2.845) -- (0.944,-2.597);
\draw[pedge] (0.472,-1.452) -- (0.472,-1.145) -- (0.763,-1.051);
\draw[pedge] (-0.000,-1.299) -- (0.292,-1.393) -- (0.472,-1.145);
\draw[pedge] (0.472,-1.452) -- (0.472,-1.145) -- (0.292,-1.393);
\draw[pedge] (0.944,-1.299) -- (0.944,-1.605) -- (1.235,-1.700);
\draw[pedge] (0.472,-1.452) -- (0.763,-1.357) -- (0.944,-1.605);
\draw[pedge] (0.944,-1.299) -- (0.944,-1.605) -- (0.763,-1.357);
\draw[pedge] (0.944,-1.299) -- (0.763,-1.051) -- (0.763,-1.357);
\draw[pedge] (0.472,-1.452) -- (0.763,-1.357) -- (0.763,-1.051);
\draw[pedge] (0.292,-0.401) -- (-0.000,-0.496) -- (-0.000,-0.803);
\draw[pedge] (0.000,0.000) -- (0.180,-0.248) -- (-0.000,-0.496);
\draw[pedge] (0.292,-0.401) -- (-0.000,-0.496) -- (0.180,-0.248);
\draw[pedge] (0.292,-0.897) -- (0.583,-0.803) -- (0.763,-1.051);
\draw[pedge] (0.292,-0.401) -- (0.292,-0.708) -- (0.583,-0.803);
\draw[pedge] (0.292,-0.897) -- (0.583,-0.803) -- (0.292,-0.708);
\draw[pedge] (0.292,-0.897) -- (-0.000,-0.803) -- (0.292,-0.708);
\draw[pedge] (0.292,-0.401) -- (0.292,-0.708) -- (-0.000,-0.803);
\draw[pedge] (0.292,-0.897) -- (-0.000,-0.803) -- (0.180,-1.051);
\draw[pedge] (-0.000,-1.299) -- (0.180,-1.051) -- (-0.000,-0.803);
\draw[pedge] (0.292,-0.897) -- (0.472,-1.145) -- (0.763,-1.051);
\draw[pedge] (-0.000,-1.299) -- (0.180,-1.051) -- (0.472,-1.145);
\draw[pedge] (0.292,-0.897) -- (0.472,-1.145) -- (0.180,-1.051);
\draw[pedge] (0.472,-2.751) -- (0.292,-2.503) -- (-0.000,-2.597);
\draw[pedge] (0.763,-2.349) -- (0.583,-2.597) -- (0.292,-2.503);
\draw[pedge] (0.472,-2.751) -- (0.292,-2.503) -- (0.583,-2.597);
\draw[pedge] (-0.000,-2.904) -- (0.180,-3.152) -- (-0.000,-3.400);
\draw[pedge] (0.472,-2.751) -- (0.180,-2.845) -- (0.180,-3.152);
\draw[pedge] (-0.000,-2.904) -- (0.180,-3.152) -- (0.180,-2.845);
\draw[pedge] (-0.000,-2.904) -- (-0.000,-2.597) -- (0.180,-2.845);
\draw[pedge] (0.472,-2.751) -- (0.180,-2.845) -- (-0.000,-2.597);
\draw[pedge] (-0.000,-1.795) -- (0.292,-1.700) -- (0.472,-1.948);
\draw[pedge] (-0.000,-1.299) -- (-0.000,-1.605) -- (0.292,-1.700);
\draw[pedge] (-0.000,-1.795) -- (0.292,-1.700) -- (-0.000,-1.605);
\draw[pedge] (0.292,-2.196) -- (-0.000,-2.291) -- (-0.000,-2.597);
\draw[pedge] (-0.000,-1.795) -- (0.180,-2.043) -- (-0.000,-2.291);
\draw[pedge] (0.292,-2.196) -- (-0.000,-2.291) -- (0.180,-2.043);
\draw[pedge] (0.292,-2.196) -- (0.472,-1.948) -- (0.180,-2.043);
\draw[pedge] (-0.000,-1.795) -- (0.180,-2.043) -- (0.472,-1.948);
\draw[pedge] (0.292,-2.196) -- (0.472,-1.948) -- (0.472,-2.255);
\draw[pedge] (0.763,-2.349) -- (0.472,-2.255) -- (0.472,-1.948);
\draw[pedge] (0.292,-2.196) -- (0.292,-2.503) -- (-0.000,-2.597);
\draw[pedge] (0.763,-2.349) -- (0.472,-2.255) -- (0.292,-2.503);
\draw[pedge] (0.292,-2.196) -- (0.292,-2.503) -- (0.472,-2.255);
\draw[pedge] (0.763,-1.853) -- (0.472,-1.948) -- (0.763,-2.043);
\draw[pedge] (0.763,-2.349) -- (0.763,-2.043) -- (0.472,-1.948);
\draw[pedge] (0.763,-1.853) -- (1.055,-1.948) -- (1.235,-1.700);
\draw[pedge] (0.763,-2.349) -- (0.763,-2.043) -- (1.055,-1.948);
\draw[pedge] (0.763,-1.853) -- (1.055,-1.948) -- (0.763,-2.043);
\draw[pedge] (0.472,-1.452) -- (0.292,-1.700) -- (0.472,-1.948);
\draw[pedge] (-0.000,-1.299) -- (0.292,-1.393) -- (0.292,-1.700);
\draw[pedge] (0.472,-1.452) -- (0.292,-1.700) -- (0.292,-1.393);
\draw[pedge] (0.763,-1.853) -- (0.944,-1.605) -- (1.235,-1.700);
\draw[pedge] (0.472,-1.452) -- (0.652,-1.700) -- (0.944,-1.605);
\draw[pedge] (0.763,-1.853) -- (0.944,-1.605) -- (0.652,-1.700);
\draw[pedge] (0.763,-1.853) -- (0.472,-1.948) -- (0.652,-1.700);
\draw[pedge] (0.472,-1.452) -- (0.652,-1.700) -- (0.472,-1.948);
\draw[pedge] (1.707,-1.853) -- (1.998,-1.948) -- (1.818,-1.700);
\draw[pedge] (1.998,-1.452) -- (1.818,-1.700) -- (1.998,-1.948);
\draw[pedge] (1.707,-1.853) -- (1.527,-1.605) -- (1.235,-1.700);
\draw[pedge] (1.998,-1.452) -- (1.818,-1.700) -- (1.527,-1.605);
\draw[pedge] (1.707,-1.853) -- (1.527,-1.605) -- (1.818,-1.700);
\draw[pedge] (1.707,-2.349) -- (1.998,-2.255) -- (1.998,-1.948);
\draw[pedge] (1.998,-2.751) -- (1.818,-2.503) -- (1.998,-2.255);
\draw[pedge] (1.707,-2.349) -- (1.998,-2.255) -- (1.818,-2.503);
\draw[pedge] (1.707,-1.853) -- (1.415,-1.948) -- (1.235,-1.700);
\draw[pedge] (1.707,-2.349) -- (1.707,-2.043) -- (1.415,-1.948);
\draw[pedge] (1.707,-1.853) -- (1.415,-1.948) -- (1.707,-2.043);
\draw[pedge] (1.707,-1.853) -- (1.998,-1.948) -- (1.707,-2.043);
\draw[pedge] (1.707,-2.349) -- (1.707,-2.043) -- (1.998,-1.948);
\draw[pedge] (2.470,-1.299) -- (2.470,-1.605) -- (2.762,-1.700);
\draw[pedge] (1.998,-1.452) -- (2.290,-1.357) -- (2.470,-1.605);
\draw[pedge] (2.470,-1.299) -- (2.470,-1.605) -- (2.290,-1.357);
\draw[pedge] (2.942,-1.452) -- (2.942,-1.145) -- (3.234,-1.051);
\draw[pedge] (2.470,-1.299) -- (2.762,-1.393) -- (2.942,-1.145);
\draw[pedge] (2.942,-1.452) -- (2.942,-1.145) -- (2.762,-1.393);
\draw[pedge] (2.942,-1.452) -- (2.762,-1.700) -- (2.762,-1.393);
\draw[pedge] (2.470,-1.299) -- (2.762,-1.393) -- (2.762,-1.700);
\draw[pedge] (2.290,-2.349) -- (1.998,-2.255) -- (1.998,-1.948);
\draw[pedge] (1.998,-2.751) -- (2.179,-2.503) -- (1.998,-2.255);
\draw[pedge] (2.290,-2.349) -- (1.998,-2.255) -- (2.179,-2.503);
\draw[pedge] (2.290,-1.853) -- (2.582,-1.948) -- (2.762,-1.700);
\draw[pedge] (2.290,-2.349) -- (2.290,-2.043) -- (2.582,-1.948);
\draw[pedge] (2.290,-1.853) -- (2.582,-1.948) -- (2.290,-2.043);
\draw[pedge] (2.290,-1.853) -- (1.998,-1.948) -- (2.290,-2.043);
\draw[pedge] (2.290,-2.349) -- (2.290,-2.043) -- (1.998,-1.948);
\draw[pedge] (2.290,-1.853) -- (1.998,-1.948) -- (2.179,-1.700);
\draw[pedge] (1.998,-1.452) -- (2.179,-1.700) -- (1.998,-1.948);
\draw[pedge] (2.290,-1.853) -- (2.470,-1.605) -- (2.762,-1.700);
\draw[pedge] (1.998,-1.452) -- (2.179,-1.700) -- (2.470,-1.605);
\draw[pedge] (2.290,-1.853) -- (2.470,-1.605) -- (2.179,-1.700);
\draw[pedge] (1.235,-0.897) -- (0.944,-0.803) -- (0.763,-1.051);
\draw[pedge] (1.235,-0.401) -- (1.235,-0.708) -- (0.944,-0.803);
\draw[pedge] (1.235,-0.897) -- (0.944,-0.803) -- (1.235,-0.708);
\draw[pedge] (0.944,-1.299) -- (1.235,-1.393) -- (1.235,-1.700);
\draw[pedge] (1.235,-0.897) -- (1.055,-1.145) -- (1.235,-1.393);
\draw[pedge] (0.944,-1.299) -- (1.235,-1.393) -- (1.055,-1.145);
\draw[pedge] (0.944,-1.299) -- (0.763,-1.051) -- (1.055,-1.145);
\draw[pedge] (1.235,-0.897) -- (1.055,-1.145) -- (0.763,-1.051);
\draw[pedge] (0.292,-0.401) -- (0.472,-0.153) -- (0.763,-0.248);
\draw[pedge] (0.000,0.000) -- (0.180,-0.248) -- (0.472,-0.153);
\draw[pedge] (0.292,-0.401) -- (0.472,-0.153) -- (0.180,-0.248);
\draw[pedge] (0.763,-0.555) -- (0.583,-0.803) -- (0.763,-1.051);
\draw[pedge] (0.292,-0.401) -- (0.583,-0.496) -- (0.583,-0.803);
\draw[pedge] (0.763,-0.555) -- (0.583,-0.803) -- (0.583,-0.496);
\draw[pedge] (0.763,-0.555) -- (0.763,-0.248) -- (0.583,-0.496);
\draw[pedge] (0.292,-0.401) -- (0.583,-0.496) -- (0.763,-0.248);
\draw[pedge] (0.763,-0.555) -- (0.763,-0.248) -- (0.944,-0.496);
\draw[pedge] (1.235,-0.401) -- (0.944,-0.496) -- (0.763,-0.248);
\draw[pedge] (0.763,-0.555) -- (0.944,-0.803) -- (0.763,-1.051);
\draw[pedge] (1.235,-0.401) -- (0.944,-0.496) -- (0.944,-0.803);
\draw[pedge] (0.763,-0.555) -- (0.944,-0.803) -- (0.944,-0.496);
\draw[pedge] (2.470,-1.299) -- (2.290,-1.051) -- (2.470,-0.803);
\draw[pedge] (1.998,-1.452) -- (2.290,-1.357) -- (2.290,-1.051);
\draw[pedge] (2.470,-1.299) -- (2.290,-1.051) -- (2.290,-1.357);
\draw[pedge] (2.762,-0.897) -- (2.942,-1.145) -- (3.234,-1.051);
\draw[pedge] (2.470,-1.299) -- (2.650,-1.051) -- (2.942,-1.145);
\draw[pedge] (2.762,-0.897) -- (2.942,-1.145) -- (2.650,-1.051);
\draw[pedge] (2.762,-0.897) -- (2.470,-0.803) -- (2.650,-1.051);
\draw[pedge] (2.470,-1.299) -- (2.650,-1.051) -- (2.470,-0.803);
\draw[pedge] (1.707,-0.555) -- (1.527,-0.803) -- (1.707,-1.051);
\draw[pedge] (1.235,-0.401) -- (1.527,-0.496) -- (1.527,-0.803);
\draw[pedge] (1.707,-0.555) -- (1.527,-0.803) -- (1.527,-0.496);
\draw[pedge] (1.998,-0.956) -- (2.179,-0.708) -- (2.470,-0.803);
\draw[pedge] (1.707,-0.555) -- (1.887,-0.803) -- (2.179,-0.708);
\draw[pedge] (1.998,-0.956) -- (2.179,-0.708) -- (1.887,-0.803);
\draw[pedge] (1.998,-0.956) -- (1.707,-1.051) -- (1.887,-0.803);
\draw[pedge] (1.707,-0.555) -- (1.887,-0.803) -- (1.707,-1.051);
\draw[pedge] (1.998,-0.956) -- (1.707,-1.051) -- (1.998,-1.145);
\draw[pedge] (1.998,-1.452) -- (1.998,-1.145) -- (1.707,-1.051);
\draw[pedge] (1.998,-0.956) -- (2.290,-1.051) -- (2.470,-0.803);
\draw[pedge] (1.998,-1.452) -- (1.998,-1.145) -- (2.290,-1.051);
\draw[pedge] (1.998,-0.956) -- (2.290,-1.051) -- (1.998,-1.145);
\draw[pedge] (1.527,-1.299) -- (1.707,-1.051) -- (1.707,-1.357);
\draw[pedge] (1.998,-1.452) -- (1.707,-1.357) -- (1.707,-1.051);
\draw[pedge] (1.527,-1.299) -- (1.527,-1.605) -- (1.235,-1.700);
\draw[pedge] (1.998,-1.452) -- (1.707,-1.357) -- (1.527,-1.605);
\draw[pedge] (1.527,-1.299) -- (1.527,-1.605) -- (1.707,-1.357);
\draw[pedge] (1.235,-0.897) -- (1.527,-0.803) -- (1.707,-1.051);
\draw[pedge] (1.235,-0.401) -- (1.235,-0.708) -- (1.527,-0.803);
\draw[pedge] (1.235,-0.897) -- (1.527,-0.803) -- (1.235,-0.708);
\draw[pedge] (1.527,-1.299) -- (1.235,-1.393) -- (1.235,-1.700);
\draw[pedge] (1.235,-0.897) -- (1.415,-1.145) -- (1.235,-1.393);
\draw[pedge] (1.527,-1.299) -- (1.235,-1.393) -- (1.415,-1.145);
\draw[pedge] (1.527,-1.299) -- (1.707,-1.051) -- (1.415,-1.145);
\draw[pedge] (1.235,-0.897) -- (1.415,-1.145) -- (1.707,-1.051);

\end{tikzpicture}
\caption{}
\end{subfigure}
\caption{\textit{а}~— молекула C$_{60}$: усечённый икосаэдр (пятиугольники
закрашены; расчёт); \textit{б}~— мозаика Пенроуза из толстых (жёлтых,
$72^\circ/108^\circ$) и тонких (синих, $36^\circ/144^\circ$) ромбов (расчёт
подразбиением). Мозаика имеет ось 5-го порядка, но не периодична:
отношение числа толстых и тонких ромбов~— иррациональное число
$\tau=(1+\sqrt5)/2$}
\label{fig:c60}\label{fig:penrose}
\end{figure}

\paragraph{Фуллериты и нанотрубки}
Кристаллы из молекул C$_{60}$, C$_{70}$~— \term{фуллериты}. Фуллерит C$_{60}$ при
комнатной температуре имеет ГЦК-решётку с $a=1{,}417$~нм; радиусы
октаэдрических и тетраэдрических пустот 0,207 и 0,112~нм (в них можно
внедрять атомы щелочных металлов; K$_3$C$_{60}$~— сверхпроводник). Фуллерит
C$_{70}$ имеет ГЦК-модификацию с $a=1{,}501$~нм и ГПУ с $a\approx1{,}06$~нм,
$c\approx1{,}72$~нм\ispr{В оригинале для ГПУ-модификации C$_{70}$ указано
$a=1{,}501$~нм; при $c=1{,}718$~нм это дало бы $c/a=1{,}14$, что
невозможно для плотной упаковки молекул. Расстояние между соседними
молекулами в ГЦК-фазе $1{,}501/\sqrt2\approx1{,}06$~нм, и это же значение~—
период $a$ ГПУ-фазы ($c/a\approx1{,}63$).}. Графитовый слой, свёрнутый в цилиндр,~—
\term{углеродная нанотрубка}: диаметр~— единицы нанометров, длина~— до
десятков микрометров; высокая прочность и теплопроводность, большая
допустимая плотность тока.

\paragraph{Квазикристаллы}
В 1984~г. в быстро закалённом (${\sim}10^6$~К/с) сплаве Al--(14--20)~ат.\,\%~Mn
Д.~Шехтман обнаружил дифракционную картину с осью 5-го порядка,
запрещённой для кристаллов. Сейчас известны сотни \term{квазикристаллов},
есть и с осями 8, 10, 12-го порядков. В них нет трансляционной
симметрии, но порядок дальний, «по закону», а не хаотический, как у
аморфных тел,~— поэтому дифракционные максимумы резкие. Модели строят
\emph{проекционным методом}: берут решётку в шестимерном пространстве и
проецируют её узлы, близкие к иррационально ориентированному трёхмерному
подпространству. Двумерная модель~— \term{мозаика Пенроуза}
(рис.~\ref{fig:penrose}): плоскость без пустот и наложений покрывается
ромбами двух сортов~— тонкими ($36^\circ$ и $144^\circ$) и толстыми ($72^\circ$ и
$108^\circ$); отношение их чисел в бесконечной мозаике равно золотому сечению
$\tau=(\sqrt5+1)/2=1{,}618\ldots$, и, поскольку оно иррационально, мозаика не
является двумерным кристаллом.

Свойства квазикристаллов необычны. Они состоят из металлов, но их
электросопротивление при низких температурах аномально велико,
\emph{уменьшается} с ростом температуры и \emph{растёт} с совершенством
структуры (после отжига дефектов). Модули упругости ниже, чем у
кристаллических фаз близкого состава (ближе к аморфным металлам), но
движение дислокаций сильно затруднено, поэтому квазикристаллы хрупки и
интересны как упрочняющие частицы в сплавах.

\subsection{Структура поверхности}\label{ssec:surface}

Поверхность определяет коррозию, трение, износ, адгезию, хрупкость.
Мысленно получим чистую поверхность, рассекая монокристалл по плоскости
$(hkl)$ и удаляя верхнюю часть. Атомы поверхности лишились части связей,
и в приповерхностном слое идут перестройки:
\begin{itemize}
\item \term{релаксация}~— изменение межплоскостного расстояния $d$ в
нескольких верхних слоях по сравнению с объёмным $d_0$ (без изменения
периодичности в плоскости);
\item \term{реконструкция}~— перестройка поверхностных атомов (часто с
участием адсорбированных атомов) с образованием \emph{новой
периодичности} в плоскости поверхности.
\end{itemize}
Структуру поверхности записывают в виде (обозначения Вуда)
\begin{equation}\label{eq:wood}
M(hkl)\text{--}(n\times m)\text{--}C ,
\end{equation}
где $M$~— элемент подложки, $(hkl)$~— грань, $n$ и $m$~— во сколько раз новый
период больше исходного вдоль $\vb a$ и $\vb b$, $C$~— адсорбированный элемент;
приставка «c» означает центрированную ячейку, «R$\alpha$»~— поворот ячейки
на угол $\alpha$ (рис.~\ref{fig:surface}).

\begin{figure}[ht]
\centering
\foreach \typ/\lab in {1/{$(2\times1)$}, 2/{$(2\times2)$}, 3/{$c(2\times2)=(\sqrt2\times\sqrt2)R45^\circ$}}{%
\begin{tikzpicture}[scale=0.55]
  \foreach \i in {0,...,5}{\foreach \j in {0,...,5}{
    \ifnum\typ=1
      \pgfmathsetmacro{\sh}{mod(\i,2)==0 ? 0.18 : -0.18}
      \fill[atomA, draw=atomLine, line width=0.3pt] (\i+\sh,\j) circle (0.3);
    \else
      \fill[atomA, draw=atomLine, line width=0.3pt] (\i,\j) circle (0.3);
    \fi}}
  \ifnum\typ=2
    \foreach \i in {0,2,4}{\foreach \j in {0,2,4}{\fill[fmRed] (\i+0.5,\j+0.5) circle (0.17);}}
    \draw[fmBlue, thick] (0.5,0.5) rectangle (2.5,2.5);
  \fi
  \ifnum\typ=3
    \foreach \i in {0,...,4}{\foreach \j in {0,...,4}{
      \pgfmathtruncatemacro{\pp}{mod(\i+\j,2)}
      \ifnum\pp=0 \fill[fmRed] (\i+0.5,\j+0.5) circle (0.17);\fi}}
    \draw[fmBlue, thick] (0.5,0.5) rectangle (2.5,2.5);
    \draw[fmGreen!60!black, thick, dashed] (1.5,0.5) -- (2.5,1.5) -- (1.5,2.5) -- (0.5,1.5) -- cycle;
  \fi
  \ifnum\typ=1
    \draw[fmBlue, thick] (-0.5,-0.5) rectangle (1.5,0.5);
  \fi
  \node[font=\small] at (2.5,-1.2) {\lab};
\end{tikzpicture}\hspace{5mm}}
\caption{Реконструкция поверхности (схема): \textit{а}~— $(2\times1)$: атомы
сближаются попарно (димеры), период вдоль одной оси удваивается;
\textit{б}~— $(2\times2)$: адсорбированные атомы (красные) в каждой второй
позиции вдоль обеих осей; \textit{в}~— $c(2\times2)$: адатомы в шахматном
порядке; центрированная ячейка (синяя) эквивалентна ячейке
$(\sqrt2\times\sqrt2)$, повёрнутой на $45^\circ$ (зелёная)}
\label{fig:surface}
\end{figure}

Примеры: Ni(111)--$(2\times2)$--O (кислород на никеле); W(110)--c$(9\times5)$--CO;
чистая грань Si(100) имеет структуру $(2\times1)$ (атомы образуют димеры), а
адсорбция водорода восстанавливает исходную периодичность: Si(100)--$(1\times1)$--H.
Чистые грани Au реконструированы: Au(100)--$(5\times20)$, Au(110)--$(1\times2)$
(«пропущенные ряды»), Au(111)--$(22\times\sqrt3)$. У ОЦК-металлов при низких
температурах реконструирована грань (100) W и Mo~— $c(2\times2)$. Релаксация
обычно сводится к \emph{сжатию} верхнего межплоскостного расстояния, тем
большему, чем «рыхлее» грань: у Cu(110) около $-10\,\%$, Al(110) $-5\ldots-10\,\%$,
W(100) $-5\ldots-10\,\%$, Mo(100) около $-11\,\%$, Fe(111) около $-15\,\%$; у плотных граней
релаксация мала (Ni(111) около $-1\,\%$; Al(111) даже слегка расширяется).
Структуру поверхности изучают дифракцией медленных электронов
(энергия 10--200~эВ: такие электроны проникают лишь на несколько
атомных слоёв).

\begin{itogi}
\begin{itemize}
\item Две плотнейшие упаковки~— ГПУ ($ABAB$) и ГЦК ($ABCABC$), компактность
74\,\%, к.\,ч.~=~12; пустот октаэдрических~— по одной, тетраэдрических~— по
две на атом ($0{,}414R$ и $0{,}225R$).
\item ОЦК: к.\,ч.~=~8(+6), компактность 68\,\%; тетрапустоты ($0{,}291R$) крупнее
октапустот ($0{,}155R$).
\item Многие соединения~— плотная упаковка одних атомов с заполнением
пустот другими: NaCl (все октапустоты ГЦК), сфалерит (половина тетрапустот
ГЦК), вюрцит (половина тетрапустот ГПУ), NiAs (все октапустоты ГПУ),
флюорит (все тетрапустоты ГЦК).
\item Сверхструктуры: L1$_2$ (Cu$_3$Au), L1$_0$ (CuAu), B2 (CuZn), D0$_3$
(Fe$_3$Al), L2$_1$ (Cu$_2$MnAl).
\item Квазикристаллы~— дальний порядок без трансляций (оси 5, 8, 10, 12);
поверхность релаксирует и реконструируется.
\end{itemize}
\end{itogi}

\begin{proverka}
\item Почему коэффициент компактности ГЦК и ГПУ одинаков?
\item Сколько октаэдрических и тетраэдрических пустот приходится на атом в
ГЦК- и ОЦК-решётках?
\item Почему углерод в аустените занимает октаэдрические пустоты, а в
феррите~— сплюснутые октаэдрические, хотя тетраэдрические там крупнее?
\item Какую часть тетраэдрических пустот занимают атомы в структуре алмаза
и почему её к.\,ч. равно 4?
\item Чем отличаются структуры сфалерита и вюрцита?
\item Почему CsCl не имеет ОЦК-решётку, хотя атомы расположены по узлам ОЦК?
\item Сколько атомов Fe и Al в ячейке D0$_3$ на рис.~\ref{fig:superstructures}\,\textit{г}?
\item Что означает запись Si(111)--$(7\times7)$?
\end{proverka}
